\documentclass[11pt]{article}
\usepackage[margin=1in]{geometry}
\usepackage{amsmath,amssymb,amsthm,mathtools}
\usepackage{booktabs,longtable,array,graphicx}
\usepackage{xurl}
\usepackage{hyperref}
\usepackage{xcolor}
\usepackage{microtype}
\usepackage{enumitem}
\usepackage[T1]{fontenc}
\usepackage{lmodern}
\setlength{\emergencystretch}{2em}
\setcounter{tocdepth}{1}
\numberwithin{equation}{section}
\raggedbottom
\hypersetup{colorlinks=true,linkcolor=blue,citecolor=blue,urlcolor=blue}
\newtheorem{theorem}{Theorem}[section]
\newtheorem{proposition}[theorem]{Proposition}
\newtheorem{lemma}[theorem]{Lemma}
\newtheorem{corollary}[theorem]{Corollary}
\theoremstyle{remark}
\newtheorem{remark}[theorem]{Remark}
\theoremstyle{definition}
\newtheorem{question}[theorem]{Question}
\newtheorem*{executionhypothesis}{Packet execution hypothesis}
\theoremstyle{plain}
\newtheorem{theoremA}[theorem]{Theorem}
\newtheorem{theoremB}[theorem]{Theorem}
\newcommand{\OL}{\mathcal O_L}
\newcommand{\OM}{\mathcal O_M}
\newcommand{\Ql}{\mathbb Q}
\newcommand{\Z}{\mathbb Z}
\newcommand{\R}{\mathbb R}
\newcommand{\C}{\mathbb C}
\newcommand{\HH}{\mathbb H}
\newcommand{\Gal}{\operatorname{Gal}}
\newcommand{\End}{\operatorname{End}}
\newcommand{\Tr}{\operatorname{Tr}}
\newcommand{\Nm}{\operatorname{N}}
\newcommand{\GL}{\operatorname{GL}}
\newcommand{\SL}{\operatorname{SL}}
\newcommand{\extD}{\ensuremath{D_{\mathrm{CM}}}}
\newcommand{\extAlg}{\operatorname{Alg}}
\newcommand{\extHdg}{\operatorname{Hdg}}
\newcommand{\extHg}{\operatorname{Hg}}
\newcommand{\extsupp}{\operatorname{supp}}
\newcommand{\digest}[2]{\scalebox{0.9}{\texttt{#1#2}}}

\makeatletter
\renewcommand{\l@section}{\@dottedtocline{1}{0em}{2.3em}}
\makeatother
\newcommand{\externalref}[2]{\textnormal{\csname doccode@#1\endcsname.\ifcsname xnumber@#2\endcsname\csname xnumber@#2\endcsname\else ?\fi}}
% Generated from native auxiliary files; do not invent result numbers.
\expandafter\def\csname xnumber@app:arithdetail\endcsname{2}
\expandafter\def\csname xnumber@app:components\endcsname{7}
\expandafter\def\csname xnumber@app:computer\endcsname{13}
\expandafter\def\csname xnumber@app:orders\endcsname{8}
\expandafter\def\csname xnumber@app:prime43\endcsname{3}
\expandafter\def\csname xnumber@cor:12collapse\endcsname{11.5}
\expandafter\def\csname xnumber@cor:Gtransport\endcsname{5.5}
\expandafter\def\csname xnumber@cor:HCspecific\endcsname{5.12}
\expandafter\def\csname xnumber@cor:Weilalg\endcsname{5.9}
\expandafter\def\csname xnumber@cor:d7geometrycomplete\endcsname{9.12}
\expandafter\def\csname xnumber@cor:d7r83all\endcsname{9.9}
\expandafter\def\csname xnumber@cor:d7r83nontrivial\endcsname{9.7}
\expandafter\def\csname xnumber@cor:finitegeometric\endcsname{7.3}
\expandafter\def\csname xnumber@cor:geometry119\endcsname{9.15}
\expandafter\def\csname xnumber@cor:maximalgeometrycomplete\endcsname{11.10}
\expandafter\def\csname xnumber@cor:noAstablepp\endcsname{10.4}
\expandafter\def\csname xnumber@cor:noBstablepp\endcsname{4.4}
\expandafter\def\csname xnumber@cor:noCstablepp\endcsname{10.2}
\expandafter\def\csname xnumber@cor:orientedcollapse\endcsname{11.4}
\expandafter\def\csname xnumber@cor:rosatipi\endcsname{12.6}
\expandafter\def\csname xnumber@eight:cor:mt-splitting\endcsname{4.3}
\expandafter\def\csname xnumber@eight:cor:subfields\endcsname{4.2}
\expandafter\def\csname xnumber@eight:sec:lattice\endcsname{4.1}
\expandafter\def\csname xnumber@eight:sec:main\endcsname{4}
\expandafter\def\csname xnumber@eight:sec:sources\endcsname{4.2}
\expandafter\def\csname xnumber@eight:thm:all\endcsname{4.1}
\expandafter\def\csname xnumber@eq:42test\endcsname{6.4}
\expandafter\def\csname xnumber@eq:CMcriterion\endcsname{5.3}
\expandafter\def\csname xnumber@eq:Deltak\endcsname{8.1}
\expandafter\def\csname xnumber@eq:E1\endcsname{4.1}
\expandafter\def\csname xnumber@eq:HermitianE1\endcsname{5.1}
\expandafter\def\csname xnumber@eq:absnormalpha\endcsname{3.2}
\expandafter\def\csname xnumber@eq:adjdet\endcsname{12.1}
\expandafter\def\csname xnumber@eq:alpha0\endcsname{3.1}
\expandafter\def\csname xnumber@eq:companion\endcsname{6.1}
\expandafter\def\csname xnumber@eq:correct-charpoly\endcsname{12.2}
\expandafter\def\csname xnumber@eq:correct-equiv\endcsname{12.6}
\expandafter\def\csname xnumber@eq:correctschwarzcomposition\endcsname{12.9}
\expandafter\def\csname xnumber@eq:deltagamma\endcsname{12.5}
\expandafter\def\csname xnumber@eq:explicitcovergroup\endcsname{7.1}
\expandafter\def\csname xnumber@eq:galois-exponents\endcsname{12.7}
\expandafter\def\csname xnumber@eq:hgparamformula\endcsname{12.8}
\expandafter\def\csname xnumber@eq:mcmullenF\endcsname{6.3}
\expandafter\def\csname xnumber@eq:minpoly\endcsname{2.1}
\expandafter\def\csname xnumber@eq:normfiber\endcsname{4.1}
\expandafter\def\csname xnumber@eq:normpi\endcsname{12.3}
\expandafter\def\csname xnumber@eq:orbit-test\endcsname{6.5}
\expandafter\def\csname xnumber@eq:pi\endcsname{3.5}
\expandafter\def\csname xnumber@eq:principalpacketpoint\endcsname{3.1}
\expandafter\def\csname xnumber@eq:unitnormquotient\endcsname{11.1}
\expandafter\def\csname xnumber@eq:unitpi\endcsname{12.4}
\expandafter\def\csname xnumber@eq:vertexdisc\endcsname{5.2}
\expandafter\def\csname xnumber@eq:wi\endcsname{6.2}
\expandafter\def\csname xnumber@ext:eq:basis\endcsname{{$\mathcal  R$}}
\expandafter\def\csname xnumber@ext:lem:contraction\endcsname{2.4}
\expandafter\def\csname xnumber@ext:lem:lattice\endcsname{2.5}
\expandafter\def\csname xnumber@ext:sec:evidence\endcsname{2.6}
\expandafter\def\csname xnumber@ext:sec:inputs\endcsname{2.2}
\expandafter\def\csname xnumber@ext:sec:lattice\endcsname{2.4}
\expandafter\def\csname xnumber@ext:sec:main\endcsname{2}
\expandafter\def\csname xnumber@ext:thm:individual\endcsname{2.3}
\expandafter\def\csname xnumber@ext:thm:mixed\endcsname{2.1}
\expandafter\def\csname xnumber@hyp:packetexecution\endcsname{1.2}
\expandafter\def\csname xnumber@lem:adjugate\endcsname{12.5}
\expandafter\def\csname xnumber@lem:degree\endcsname{12.2}
\expandafter\def\csname xnumber@lem:narrowclass\endcsname{11.2}
\expandafter\def\csname xnumber@lem:packetpointfibers\endcsname{3.1}
\expandafter\def\csname xnumber@lem:primitivebinary\endcsname{5.10}
\expandafter\def\csname xnumber@lem:twochambers\endcsname{3.5}
\expandafter\def\csname xnumber@pp:cor:fermatprojector\endcsname{2.3}
\expandafter\def\csname xnumber@pp:cor:jacobianlist\endcsname{12.2}
\expandafter\def\csname xnumber@pp:cor:obstructedmodulifields\endcsname{6.3}
\expandafter\def\csname xnumber@pp:cor:quotientrealforms\endcsname{7.2}
\expandafter\def\csname xnumber@pp:cor:quotientweilhodge\endcsname{10.1}
\expandafter\def\csname xnumber@pp:cor:smallerjacobian\endcsname{13.2}
\expandafter\def\csname xnumber@pp:cor:underlyinggalois\endcsname{4.4}
\expandafter\def\csname xnumber@pp:cor:weilimage\endcsname{8.7}
\expandafter\def\csname xnumber@pp:eq:basis\endcsname{7.1}
\expandafter\def\csname xnumber@pp:eq:curve\endcsname{7.3}
\expandafter\def\csname xnumber@pp:eq:group\endcsname{7.2}
\expandafter\def\csname xnumber@pp:eq:unpolarizedlocaldata\endcsname{7.4}
\expandafter\def\csname xnumber@pp:prop:actualpolarizations\endcsname{7.12}
\expandafter\def\csname xnumber@pp:prop:allunderlyinggalois\endcsname{4.5}
\expandafter\def\csname xnumber@pp:prop:balancedtheta\endcsname{13.1}
\expandafter\def\csname xnumber@pp:prop:bcisogeny\endcsname{7.5}
\expandafter\def\csname xnumber@pp:prop:explicit\endcsname{7.1}
\expandafter\def\csname xnumber@pp:prop:fermatbcminimum\endcsname{3.1}
\expandafter\def\csname xnumber@pp:prop:fermatminpolarization\endcsname{3.2}
\expandafter\def\csname xnumber@pp:prop:fermatrestricted\endcsname{2.2}
\expandafter\def\csname xnumber@pp:prop:homideal\endcsname{7.8}
\expandafter\def\csname xnumber@pp:prop:idealclassmonoid\endcsname{7.10}
\expandafter\def\csname xnumber@pp:prop:markeddistinct\endcsname{7.3}
\expandafter\def\csname xnumber@pp:prop:markedpolarizations\endcsname{7.11}
\expandafter\def\csname xnumber@pp:prop:mindegree\endcsname{7.7}
\expandafter\def\csname xnumber@pp:prop:minimumpolarizedgalois\endcsname{5.2}
\expandafter\def\csname xnumber@pp:prop:picardorder\endcsname{7.9}
\expandafter\def\csname xnumber@pp:prop:polarizationsquare\endcsname{8.5}
\expandafter\def\csname xnumber@pp:prop:sixprincipal\endcsname{8.3}
\expandafter\def\csname xnumber@pp:prop:stabilizer\endcsname{7.2}
\expandafter\def\csname xnumber@pp:prop:trigonaltheta\endcsname{12.3}
\expandafter\def\csname xnumber@pp:prop:twoprimefermat\endcsname{11.1}
\expandafter\def\csname xnumber@pp:prop:unitcurve\endcsname{8.2}
\expandafter\def\csname xnumber@pp:prop:unitfamily\endcsname{8.1}
\expandafter\def\csname xnumber@pp:prop:unmarkeda\endcsname{8.4}
\expandafter\def\csname xnumber@pp:prop:unmarkeddistinct\endcsname{7.4}
\expandafter\def\csname xnumber@pp:prop:unpolarizeddistinct\endcsname{7.6}
\expandafter\def\csname xnumber@pp:rem:obstruction\endcsname{8.8}
\expandafter\def\csname xnumber@pp:rem:unitimage\endcsname{8.6}
\expandafter\def\csname xnumber@pp:sec:fermatintegral\endcsname{2}
\expandafter\def\csname xnumber@pp:sec:thetacertificate\endcsname{12.1}
\expandafter\def\csname xnumber@pp:thm:completequotientmoduli\endcsname{6.4}
\expandafter\def\csname xnumber@pp:thm:completequotientorbits\endcsname{6.2}
\expandafter\def\csname xnumber@pp:thm:definitiondegrees\endcsname{9.1}
\expandafter\def\csname xnumber@pp:thm:fermatallquotients\endcsname{4.1}
\expandafter\def\csname xnumber@pp:thm:fermatunderlying\endcsname{4.3}
\expandafter\def\csname xnumber@pp:thm:introgeometry\endcsname{1.3}
\expandafter\def\csname xnumber@pp:thm:jacobiancomplete\endcsname{12.4}
\expandafter\def\csname xnumber@pp:thm:obstructedisogenies\endcsname{10.2}
\expandafter\def\csname xnumber@pp:thm:obstructedminimum\endcsname{5.1}
\expandafter\def\csname xnumber@pp:thm:quotientdefinitionfield\endcsname{8.1}
\expandafter\def\csname xnumber@pp:thm:quotientgalois\endcsname{4.2}
\expandafter\def\csname xnumber@pp:thm:quotientrayfield\endcsname{6.1}
\expandafter\def\csname xnumber@pp:thm:quotientrealmodels\endcsname{7.1}
\expandafter\def\csname xnumber@pp:thm:realdefinitionfield\endcsname{8.2}
\expandafter\def\csname xnumber@pp:thm:tenprincipal\endcsname{12.1}
\expandafter\def\csname xnumber@prop:2816\endcsname{12.10}
\expandafter\def\csname xnumber@prop:Amaximalpair\endcsname{11.7}
\expandafter\def\csname xnumber@prop:Appisogenous\endcsname{10.3}
\expandafter\def\csname xnumber@prop:BCdistinct\endcsname{5.16}
\expandafter\def\csname xnumber@prop:Cppisogenous\endcsname{10.1}
\expandafter\def\csname xnumber@prop:Cvertex\endcsname{5.15}
\expandafter\def\csname xnumber@prop:DBsame\endcsname{5.17}
\expandafter\def\csname xnumber@prop:E1type\endcsname{4.1}
\expandafter\def\csname xnumber@prop:Kcompatible\endcsname{4.7}
\expandafter\def\csname xnumber@prop:actualcmtype\endcsname{5.7}
\expandafter\def\csname xnumber@prop:alltraceorders\endcsname{8.3}
\expandafter\def\csname xnumber@prop:broadersevenexclusions\endcsname{9.14}
\expandafter\def\csname xnumber@prop:canonical-fixed\endcsname{6.1}
\expandafter\def\csname xnumber@prop:charpoly\endcsname{12.7}
\expandafter\def\csname xnumber@prop:cmrepair\endcsname{12.8}
\expandafter\def\csname xnumber@prop:commonprincipal\endcsname{6.1}
\expandafter\def\csname xnumber@prop:d3ambientpoints\endcsname{3.2}
\expandafter\def\csname xnumber@prop:d7batch10\endcsname{9.5}
\expandafter\def\csname xnumber@prop:d7finalpairs\endcsname{9.6}
\expandafter\def\csname xnumber@prop:d7mask14\endcsname{9.1}
\expandafter\def\csname xnumber@prop:d7mask15_31\endcsname{9.2}
\expandafter\def\csname xnumber@prop:d7mask1_17\endcsname{9.4}
\expandafter\def\csname xnumber@prop:d7mask6_22\endcsname{9.3}
\expandafter\def\csname xnumber@prop:degree-obstruction\endcsname{12.3}
\expandafter\def\csname xnumber@prop:discminusone\endcsname{5.8}
\expandafter\def\csname xnumber@prop:effectiveorbit\endcsname{7.2}
\expandafter\def\csname xnumber@prop:fixedmatrixlattice\endcsname{5.4}
\expandafter\def\csname xnumber@prop:geometricB\endcsname{5.2}
\expandafter\def\csname xnumber@prop:knownCM\endcsname{5.18}
\expandafter\def\csname xnumber@prop:mainarith\endcsname{5.1}
\expandafter\def\csname xnumber@prop:noprincipal\endcsname{4.2}
\expandafter\def\csname xnumber@prop:orientationt\endcsname{7.1}
\expandafter\def\csname xnumber@prop:pi1\endcsname{3.1}
\expandafter\def\csname xnumber@prop:ppisogenous\endcsname{4.3}
\expandafter\def\csname xnumber@prop:principalstabilizer\endcsname{6.3}
\expandafter\def\csname xnumber@prop:residualexclusions\endcsname{11.9}
\expandafter\def\csname xnumber@prop:sixHG\endcsname{12.11}
\expandafter\def\csname xnumber@prop:sixvertexorbits\endcsname{11.6}
\expandafter\def\csname xnumber@prop:survivingorder\endcsname{6.5}
\expandafter\def\csname xnumber@q:jacobians\endcsname{12}
\expandafter\def\csname xnumber@rem:ABprincipallevels\endcsname{6.2}
\expandafter\def\csname xnumber@rem:WKnormalization\endcsname{1.4}
\expandafter\def\csname xnumber@rem:fig5signs\endcsname{5.6}
\expandafter\def\csname xnumber@rem:notgeometry\endcsname{5.2}
\expandafter\def\csname xnumber@rem:ppcover\endcsname{4.5}
\expandafter\def\csname xnumber@rem:pphec\endcsname{4.6}
\expandafter\def\csname xnumber@sec:actualcmtype\endcsname{5.4}
\expandafter\def\csname xnumber@sec:conclusion\endcsname{9}
\expandafter\def\csname xnumber@sec:fullhodge\endcsname{5.6}
\expandafter\def\csname xnumber@sec:hypergeom\endcsname{12.6}
\expandafter\def\csname xnumber@sec:intro\endcsname{1}
\expandafter\def\csname xnumber@sec:latticetransport\endcsname{5.3}
\expandafter\def\csname xnumber@sec:maincertificate\endcsname{5.1}
\expandafter\def\csname xnumber@sec:markman\endcsname{5.5}
\expandafter\def\csname xnumber@sec:polarization\endcsname{4}
\expandafter\def\csname xnumber@sec:setup\endcsname{2}
\expandafter\def\csname xnumber@sec:vertexhit\endcsname{5.2}
\expandafter\def\csname xnumber@six:sec:main\endcsname{3}
\expandafter\def\csname xnumber@six:tab:types\endcsname{1}
\expandafter\def\csname xnumber@six:thm:products\endcsname{3.1}
\expandafter\def\csname xnumber@subsec:arithmeticcensus\endcsname{3.1}
\expandafter\def\csname xnumber@tab:audit\endcsname{1}
\expandafter\def\csname xnumber@thm:42\endcsname{4.1}
\expandafter\def\csname xnumber@thm:A\endcsname{1.2}
\expandafter\def\csname xnumber@thm:Ahodge\endcsname{5.14}
\expandafter\def\csname xnumber@thm:Avertex\endcsname{5.13}
\expandafter\def\csname xnumber@thm:B\endcsname{12.1}
\expandafter\def\csname xnumber@thm:Bvertex\endcsname{5.3}
\expandafter\def\csname xnumber@thm:GLcollapse\endcsname{11.1}
\expandafter\def\csname xnumber@thm:d3classes\endcsname{8.6}
\expandafter\def\csname xnumber@thm:d7ambient\endcsname{8.7}
\expandafter\def\csname xnumber@thm:d7nonmax\endcsname{8.5}
\expandafter\def\csname xnumber@thm:d7r41complete\endcsname{9.11}
\expandafter\def\csname xnumber@thm:d7r83principal\endcsname{9.8}
\expandafter\def\csname xnumber@thm:fullhodge\endcsname{5.11}
\expandafter\def\csname xnumber@thm:fullpacketarithmetic\endcsname{3.3}
\expandafter\def\csname xnumber@thm:fullpacketgeometry\endcsname{3.6}
\expandafter\def\csname xnumber@thm:maximalallorientations\endcsname{9.10}
\expandafter\def\csname xnumber@thm:maxorder\endcsname{8.1}
\expandafter\def\csname xnumber@thm:packetHodge\endcsname{11.8}
\expandafter\def\csname xnumber@thm:packetintro\endcsname{1.1}
\expandafter\def\csname xnumber@thm:principalfamily\endcsname{6.4}
\expandafter\def\csname xnumber@thm:signature\endcsname{5.1}
\expandafter\def\csname xnumber@thm:unitnorm\endcsname{11.3}

\expandafter\def\csname doccode@packet\endcsname{I}
\expandafter\def\csname doccode@quotients\endcsname{II}
\expandafter\def\csname doccode@hodge\endcsname{H}
\expandafter\def\csname doccode@supplement\endcsname{S}

\title{A Fixed CM Packet on McMullen's Curve\\and Its Principal Models}
\author{Rollins Waltchack}
\date{3 November 2026}
\begin{document}
\maketitle
\begin{abstract}
We determine the intersection of a fixed marked CM norm packet with
McMullen's normalized curve. Its arithmetic census has $3520$
trace-labelled integral matrix classes and $368$ quotient points.
Under explicitly stated historical execution assumptions, exactly three
quotient points occur on the curve. We distinguish their nonprincipal
trace polarizations from a common self-dual lattice modification, which
gives principal models on a connected finite-cover RM family. The three
principal models remain distinct after forgetting polarizations and
markings. For the B/C models, the minimum unrestricted isogeny degree is
$7$, while the minimum scalar-polarized degree is $64$. A unit twist
preserves the principal curve image and gives six distinct principal Weil
points. The Hodge consequences concern the specified fibers; identification
with a differently normalized principal moduli locus remains open.
\end{abstract}
\paragraph{Companion references.}
Result numbers prefixed I, II, H and S refer respectively to the packet
paper, quotient paper, supporting Hodge note and computational supplement.
Each document has its own numbering and bibliography.
{\small\tableofcontents}
\clearpage
% !TeX root = ../../packet-and-principal-models.tex
\section{Introduction}\label{sec:intro}
A rational CM Hodge structure does not determine an integral polarized
abelian variety, and neither determines its location on a specified curve.
This paper separates these three questions for one finite packet on
McMullen's curve. We first determine its intersection with the normalized
curve. We then construct principal models of the three surviving fibers
and distinguish their integral maps and polarizations.

The packet is fixed by the extension $M/L$, the ordered embeddings and the
norm generator below. All conjugacy counts use $\SL_2(\mathcal O_L)$.
Replacing this group by $\GL_2(\mathcal O_L)$ or forgetting a field marking
changes the classification problem. No assertion concerns every CM point
of the curve or every Hecke packet of level $43$.

\subsection{The three questions and their answers}
The finite packet has $3520$ trace-labelled matrix classes and $368$
quotient points. Its intersection with the normalized curve is exactly
$A,B,C$, subject to the execution premise stated below. The original
fibers have polarization type $(1,7,7,21,21,21)$, not principal type.
A common self-dual lattice gives isogenous principal models on a connected
finite cover. For the B/C principal models the minimum unrestricted
isogeny degree is $7$, whereas the minimum scalar-polarized degree is $64$.
These two minima measure different requirements on a map.

The table in Theorem~\ref{thm:A} presents all three geometric fibers
together. The main statements distinguish the finite intersection, the
Hodge consequences and the integral principal geometry. The supporting
Hodge note treats powers and mixed products; none of those extensions is
needed to prove the finite intersection.

\subsection{Fixed data and execution premise}
Write $z=\zeta_{42}$, $M=\mathbb Q(z)=\mathbb Q(\zeta_{21})$,
$L=M^+$ and $t=z+z^{-1}$. Order the real embeddings by
$1,5,11,13,17,19$ and put
\[
 \pi=t^4+t^3-4t^2-2t+5,\qquad \alpha_0=1+z+z^4,
 \qquad \alpha_k=z^k\alpha_0,\quad c_k=\alpha_k+\bar\alpha_k.
\]
Then $N_{M/L}(\alpha_0)=\pi$ and $N_{L/\mathbb Q}(\pi)=43$.
The symbols $A,B,C$ denote normalized vertex classes; $A_X$ denotes the
trace-polarized sixfold at $X$, and $B_X$ its specified principal model.
Let $F_t:\mathbb H\to\mathbb H^6$ be McMullen's modular embedding
with the component normalization of Supplement, Section~\externalref{supplement}{app:components}.
The normalized curve image is the image of $F_t$ in
$\SL_2(\mathcal O_L)\backslash\mathbb H^6$. Its quadrilateral
monodromy is the integral, orientation-preserving group used below.
The normalized generators are $X_t=h_tXh_t^{-1}$ with
$h_t=\operatorname{diag}(t,1)$. The fourth quadrilateral vertex is treated
in Proposition~\ref{prop:DBsame}; it is not a fourth quotient class.

\begin{executionhypothesis}\label{hyp:packetexecution}
The packet intersection claims use the retained exact search executions
identified in Supplement, Section~\externalref{supplement}{app:computer}. For the historical commands and
continuation cohort lacking complete end-of-execution guards, we assume
that the archived mathematical sources and runtime were not transiently
replaced during execution. The accepted source correspondence and final
population join are those of the final packet adoption identified there.
Later hashes and small control runs do not establish the missing historical
guards or constitute a fresh execution of the full suite.
This premise concerns the retained execution evidence; it is additional
to the finite-domain and exact-arithmetic arguments in the proofs.
\end{executionhypothesis}

\begin{theorem}[The fixed marked norm-$\pi$ packet]\label{thm:packetintro}
Under the packet execution hypothesis above, the following holds.
Fix the extension $M/L$, its $L$-action and ordered six real embeddings,
and the generator $\pi$ and elements
$\alpha_k=\zeta_{42}^k\alpha_0$ for $0\leq k<42$ specified above.
The integral matrices of determinant $\pi$ and trace
$c_k=\alpha_k+\overline{\alpha}_k$ have exactly $3520$
trace-labelled $\SL_2(\OL)$-conjugacy classes. Their upper fixed tuples
give exactly $368$ distinct points of
$\SL_2(\OL)\backslash\HH^6$.
In the normalized trace-polarized McMullen curve image, precisely three
of these points occur: $A,B,C$. There are $365$ exclusions, consisting
of thirteen maximal-multiplier points and all $352$ proper-multiplier
points. The complete $112$-point $d=7$ subpacket has three hits and
$109$ exclusions.
\end{theorem}
\begin{proof}
The arithmetic assertion is Theorem~\ref{thm:fullpacketarithmetic}.
The geometric assertion is Theorem~\ref{thm:fullpacketgeometry}, using
the retained exact certificates and execution-provenance assumptions
specified in Supplement, Section~\externalref{supplement}{app:computer}.
\end{proof}

\begin{theoremA}\label{thm:A}
Let $z_X$ be the normalized vertex $X\in\{A,B,C\}$ and set
$A_X=\C^6/(\OL z_X+\OL)$, with trace polarization
$E_1((a,b),(a',b'))=\Tr_{L/\Ql}(a'b-ab')$.
Each has polarization type $(1,7,7,21,21,21)$ and a compatible
$K=\Ql(\sqrt{-3})$ action of signature $(3,3)$ and Hermitian
discriminant $-1$. Their actual geometric CM types and Hodge algebras are
\[
\begin{array}{c|c|c}
X&\Phi_X&\operatorname{Hdg}^*(A_X)\\ \hline
A&\{1,5,17,19,29,31\}&\langle\text{divisors}\rangle\\
B&\{1,5,23,25,29,31\}&\langle\text{divisors},W_K(A_B)\rangle\\
C&\{1,5,11,13,17,19\}&\langle\text{divisors},W_K(A_C)\rangle
\end{array}
\]
The A type is induced and the B/C types are primitive.
Consequently the Hodge conjecture holds for all three fibers in every
codimension. Their common self-dual lattice modification supplies
principal models $B_A,B_B,B_C$ on one connected finite-cover RM family,
with degree-$453789$ isogenies preserving the rational CM Hodge
structures, the indicated $K$ signatures and discriminants.
\end{theoremA}

\noindent\emph{Proof locator for Theorem~\ref{thm:A}.}
The trace polarization is computed in Proposition~\ref{prop:E1type}.
The B calculation is in Sections~\ref{sec:actualcmtype}--\ref{sec:fullhodge};
Theorem~\ref{thm:Ahodge} and Proposition~\ref{prop:Cvertex} give A and C.
Proposition~\ref{prop:commonprincipal} and
Theorem~\ref{thm:principalfamily} supply the common principal family.

\begin{theorem}[Principal models and their integral distinctions]
\label{pp:thm:introgeometry}
The specified principal models $B_A,B_B,B_C$ are pairwise nonisomorphic
even as unpolarized complex abelian varieties. For maps $B_B\to B_C$,
the least isogeny degree is $7$. Among these maps pulling back the given
target principal polarization to a positive integer multiple of the
source polarization, the least multiplier is $2$, with exactly six
degree-$64$ maps sharing one kernel. Under the normalization of
Proposition~\ref{pp:prop:homideal}, their integral Hom module is
the target-$R$ ideal $\mathfrak a$ of class order two.

On each of the two underlying varieties $B_B,B_C$, exactly two
isomorphism classes of principal polarizations occur, represented by
the given polarization and its $4-t^2$ twist. These give four distinct
polarized isomorphism classes after forgetting all CM and RM markings.
The same twist extends over the principal family; on the full coarse
principal RM component it is the nontrivial involution
$w_i\mapsto w_i/\sigma_i(4-t^2)$. It preserves the entire principal
curve image. On each of $B_A,B_B,B_C$ with its prescribed RM action
there are exactly two compatible principal-polarization classes;
the resulting six Weil points all lie on this same image and remain
distinct after forgetting all CM and RM markings. For each of the four
B/C polarization pairs, the least scalar multiplier remains $2$,
with six degree-$64$ maps in the direction $B_B\to B_C$; their two possible kernels are complementary
order-$64$ subgroups of $B_B[2]$.
\end{theorem}
\begin{proof}
See Theorem~\ref{thm:principalfamily} and
Propositions~\ref{pp:prop:unpolarizeddistinct},
\ref{pp:prop:bcisogeny}, \ref{pp:prop:mindegree},
\ref{pp:prop:homideal}, \ref{pp:prop:actualpolarizations}
and~\ref{pp:prop:unitfamily}, together with
Propositions~\ref{pp:prop:unitcurve}, \ref{pp:prop:sixprincipal},
\ref{pp:prop:unmarkeda}, \ref{pp:prop:polarizationsquare} and
Corollary~\ref{pp:cor:weilimage}. Arbitrary fine-level lifts are not
part of the involution assertion.
\end{proof}

\begin{remark}[Principal-moduli normalization]\label{rem:WKnormalization}
The original lattice $\mathcal O_L^2$ with its trace form is not a
principally polarized model. The self-dual modification changes the lattice
through an isogeny; passage to a finite cover alone would not do so.
The resulting principal family lies in the component determined by that
self-dual module. Its identification with Mostaed's named principally
polarized locus remains open. The Hodge conclusions at the three special
fibers do not require that further identification, and no global
$K$-action on the family is asserted.
\end{remark}

\subsection{Proof strategy and relation to prior work}
The arithmetic classification includes every module overorder and
orientation, then identifies the two characteristic roots at the level of
quotient points. For geometry, two simultaneous target regions reduce curve
membership to finite exact searches. Both regions are needed: separate
choices of a reflected chamber in different coordinates are not allowed.
The positive vertex realizations and exhaustive empty searches prove the
intersection theorem.

The principal construction instead uses the discriminant module of the
trace lattice. A common nonfree maximal isotropic subgroup is stable under
the three specialized quadratic actions. Its stabilizer gives the finite
cover on which the lattice modification descends. Integral lattice and
unit calculations then distinguish the models, maps and twists.

The Hodge consequences use divisor classes and the split-Weil input of
Schoen and Markman \cite{Schoen1998,Markman2025}. OpenAI's general CM
manuscript \cite{OpenAICM2026}, if valid, implies the rational algebraicity
conclusions; it is not a proof input here. Goodson's factor comparison and
Saller's related discussion \cite{Goodson2024,Saller2026} also precede any
claim of novelty for the rational CM types. The contribution examined here
is the fixed packet and its integral geometry, not new abstract isogeny
classes. Precise historical priority remains a separate literature question.
The reported Aoki route in Saller is not used as a substitute proof.

The trace-dual and unit-norm description of polarizations is classical.
For comparison, Li--Mantovan--Pries prove uniqueness, when a principal
polarization exists, for simple maximal-order CM types under specified
unit-sign hypotheses \cite[Prop.~4.5]{LiMantovanPries2021}. Their
same-field example $y^7=x^3-1$ is a known principally polarized Jacobian
\cite[Prop.~6.11 and Lemma~6.12]{LiMantovanPries2021}.
Our B/C principal models have a proper endomorphism order. The results
here concern their exact integral classes, maps and realization on the
specified curve; neither the general polarization method nor same-field
principal existence is claimed as new.

The constructed principal family is not identified with Mostaed's named
principal normalization \cite{Mostaed2026}. It has a different integral
lattice. Corrections relevant to that distinction are recorded in the
supplement; they are secondary to the positive results of this paper.

% !TeX root = ../../packet-and-principal-models.tex
\section{Arithmetic and Hilbert-modular setup}\label{sec:setup}
Let
\[
 M=\Ql(\zeta_{42}),\qquad L=M^+=\Ql(\zeta_{42}+\zeta_{42}^{-1}),
\]
so $[L:\Ql]=6$ and $[M:L]=2$. Put
\[
 t=\zeta_{42}+\zeta_{42}^{-1}=2\cos(\pi/21).
\]
Then $t$ has minimal polynomial
\begin{equation}\label{eq:minpoly}
 f(T)=T^6+T^5-6T^4-6T^3+8T^2+8T+1.
\end{equation}
We use the six real embeddings indexed by
\[
 k\in\{1,5,11,13,17,19\},\qquad \sigma_k(t)=2\cos(k\pi/21).
\]

For a prime ideal $\mathfrak l\subset\OL$ of absolute norm $\ell$, the Hilbert-modular Hecke correspondence is represented by matrices whose determinant ideal is $\mathfrak l$. When $\mathfrak l$ is narrowly principal and a totally positive generator $\pi$ is fixed, standard representatives have determinant $\pi$. This is also the determinant convention used explicitly in Mostaed's Section 5.3: if $\alpha=\delta\gamma_j$, then $\det_L(\alpha)=\pi$ \cite[Sec.~5.3]{Mostaed2026}.



% !TeX root = ../../packet-and-principal-models.tex
\section{The complete packet: arithmetic and geometry}
The proof separates the finite arithmetic domain from curve membership.
The supplementary order calculations feed into the following full census;
its proof includes noninvertible modules over the originally generated order.
Historical partial exclusion totals are kept in the supplement and play no
role as substitutes for the final two-target cover.

\subsection{The complete marked arithmetic census}\label{subsec:arithmeticcensus}

Throughout this subsection the extension $M/L$, its $L$-action and the
ordered six real embeddings are fixed.  Write $z=\zeta_{42}$ and
\[
 \alpha_k=z^k(1+z+z^4)=A_k+B_kz,\qquad
 c_k=2A_k+tB_k,\qquad 0\leq k<42.
\]
We count integral matrices of trace $c_k$ and determinant $\pi$, up to
$\SL_2(\OL)$ conjugacy, and their upper fixed tuples in
$\SL_2(\OL)\backslash\HH^6$.  Thus the universe is the norm-$\pi$ packet
in this marked CM field, rather than a union over other CM fields.
For an ideal $\mathfrak f\subseteq\OL$, put
$S_{\mathfrak f}=\OL+\mathfrak f\OM$; in particular
$S_{\OL}=\OM$.

\begin{lemma}[Point fibers and the intrinsic multiplier order]
\label{lem:packetpointfibers}
For any fixed trace $c_k$, the map from its integral matrix classes to
quotient upper fixed points has exactly two classes in each fiber,
represented by $T$ and $c_kI-T$.  The relative conductor ideal of the
multiplier order is an invariant of the quotient point.  In particular,
different embedded ideals in Proposition~\externalref{supplement}{prop:alltraceorders}
give disjoint proper-multiplier point populations.
\end{lemma}
\begin{proof}
The two-root argument in Theorem~\externalref{supplement}{thm:d7ambient} applies to every
$c_k$: the characteristic polynomial is irreducible, since its
discriminant $B_k^2(t^2-4)$ is negative at every real embedding.
If two actions fix the same upper tuple, they preserve both conjugate
complex lines and commute.  Injectivity of a real embedding makes this
commutation an identity over $L$.  They therefore belong to the same
quadratic centralizer, where their characteristic polynomial has only
its two roots.  An anti-intertwiner has determinant $-N_{M/L}(a)$,
totally negative, so the two root classes are distinct under $\SL_2$.

The $L$-algebra preserving those lines is intrinsic to the tuple.
Its intersection with $M_2(\OL)$ is the multiplier ring of the
corresponding lattice.  An integral determinant-one change of basis
conjugates both the algebra and this order.  The two identifications of
that algebra with $M$ fixing $L$ differ only by complex conjugation,
which preserves every $S_{\mathfrak f}$.  Hence the exact ideal
$\mathfrak f$, not merely its norm, is invariant.  Extra endomorphisms
of an unmarked abelian variety do not enlarge this two-by-two
$L$-linear centralizer, and automorphisms permuting the marked real
embeddings are not quotient identifications.
\end{proof}

\begin{proposition}[The $144$ selected $d=3$ ambient points]
\label{prop:d3ambientpoints}
The $1280$ matrix classes of Theorem~\externalref{supplement}{thm:d3classes} define exactly
$144$ distinct quotient points.  Their decomposition and the number
of those matrix classes incident to each point are
\[
\begin{array}{c|r|r}
\text{Multiplier conductor}&\text{Points}&\text{Incidences per point}\\\hline
\OL&16&24\\
(7,t-1)&32&16\\
(7,t-1)^2&32&4\\
(43,t-20)&32&4\\
(43,t-31)&32&4.
\end{array}
\]
Their intersection with the $112$-point population of
Theorem~\externalref{supplement}{thm:d7ambient} consists of exactly the sixteen maximal
points; the union contains $240$ distinct ambient points.
\end{proposition}
\begin{proof}
Theorem~\externalref{supplement}{thm:d3classes} supplies the complete module and orientation
lists.  Choose the common basis $(1,gz)$ with
\[
 g=1,\ B_0,\ B_{12},\ B_4,\ B_8
\]
for the five displayed multipliers.  Let $\delta_m$ be the binary
product selected by mask $m$ from
\[
 (\xi_5,\xi_{11},\xi_{13},\xi_{17},-1,u),\qquad u=4-t^2.
\]
Use masks $0,\ldots,31$ for maximal multiplier and $0,\ldots,63$
for a proper multiplier.  Multiplication by any allowed $\alpha_k$
in orientation $\delta=\delta_m$ is
\[
 T_{k,g,\delta}=
 \begin{pmatrix}A_k&-\delta gB_k\\
 B_k/(\delta g)&A_k+tB_k\end{pmatrix}.
\]
Putting $q=g\delta$, the upper tuple is specified exactly by
\begin{equation}\label{eq:principalpacketpoint}
 2x=-qt,\qquad 4h^2=q^2(4-t^2),\qquad |w|^2=q^2,
 \qquad h>0.
\end{equation}
It is independent of the trace label.  The matrix
\[
 U_q=\begin{pmatrix}1&qt\\0&1\end{pmatrix}\in\SL_2(\OL)
\]
sends this tuple to the one for $-q$ and satisfies
$U_q(c_kI-T_{k,g,\delta})U_q^{-1}=T_{k,g,-\delta}$.
Thus root pairing toggles bit $16$.  Canonical masks are $0,\ldots,15$
for $\OM$ and $0,\ldots,15,32,\ldots,47$ for every proper multiplier.
Lemma~\ref{lem:packetpointfibers} proves that no further quotient
identifications occur.  Literal equality of the chosen tuples is a
separate assertion: their real parts force equality of $q$, whose ideal
$(q)=(g)$ separates the five strata.

We also account for the basis orientations in the old $1280$-matrix
list.  A regular basis $(1,\alpha_k)$ changes to $(1,gz)$ by
$Q=\left(\begin{smallmatrix}1&A_k\\0&e\end{smallmatrix}\right)$,
where $e=B_k/g$; for the intermediate basis use
$Q=\operatorname{diag}(1,(t-1)/B_0)$.  The nontrivial unit-square
identities $e=\delta_a v^2$ are
\[
\begin{array}{c|r|l}
 e&a&v\\\hline
 B_{15}/B_0&33&1+5t-5t^3+t^5\\
 B_{18}/B_0&39&3+13t-t^2-11t^3+2t^5\\
 (t-1)/B_0&7&1+18t+3t^2-16t^3-t^4+3t^5.
\end{array}
\]
Each displayed $v$ is a unit.  For regular labels $0,4,8,12$ use
$e=v=1,a=0$; the opposite trace label negates $e$ and toggles bit $16$
of $a$.  The intermediate basis is the same at labels $12$ and $33$.
If $n=m\mathbin{\oplus}a$ and
$h_0=\prod_{i:a_i=1,\,m_i=0}s_i$ for the six listed generators $s_i$,
then $\delta_n\delta_a/\delta_m=h_0^2$.  Consequently
\[
 P=(vh_0)^{-1}D_{\delta_n}QD_{\delta_m}^{-1},
 \qquad D_\delta=\operatorname{diag}(\delta,1),
\]
is integral of determinant one and transports the old action to its
common-basis action.  If needed, compose with $U_{g\delta_n}$ to
select the canonical root partner.  The paired exact implementations
in Supplement, Section~\externalref{supplement}{app:computer} check every unit identity and all
$1280$ matrix and upper-tuple transports.

There are twelve trace labels at maximal multiplier, eight at
multiplier $(7,t-1)$, and two at each of the other three proper
multipliers.  The two-root fibers give the displayed incidences and
$16\cdot24+32\cdot16+96\cdot4=1280$.
The intermediate conductor-seven stratum is counted once.
The proper $d=3$ conductors differ from the two proper $d=7$
conductors, whereas their common maximal tuples coincide.
Lemma~\ref{lem:packetpointfibers} gives the asserted intersection
and $144+112-16=240$.
\end{proof}

\begin{theorem}[The full marked norm-$\pi$ arithmetic census]
\label{thm:fullpacketarithmetic}
For all $42$ traces in Proposition~\externalref{supplement}{prop:alltraceorders}, the
complete integral matrix universe has $3520$ trace-labelled
$\SL_2(\OL)$ classes.  The exact counts, including every module
overorder and determinant orientation, are
\[
\begin{array}{c|r|r|r}
\text{Generated-order conductor}&\text{Traces}&
 \text{Classes per trace}&\text{Subtotal}\\\hline
\OL&16&32&512\\
(7,t-1)&12&96&1152\\
(7,t-1)^2&2&160&320\\
(41,t-12)&2&96&192\\
(41,t-22)&2&96&192\\
(43,t-20)&2&96&192\\
(43,t-31)&2&96&192\\
(83,t-36)&2&160&320\\
(127,t-94)&2&224&448.
\end{array}
\]
Their upper fixed tuples define exactly $368$ distinct ambient
quotient points, with the following intrinsic multiplier decomposition:
\[
\begin{array}{c|r|r|r}
\text{Multiplier conductor}&|\operatorname{Pic}(S_{\mathfrak f})|&
 \text{Points}&\text{Incidences per point}\\\hline
\OL&1&16&84\\
(7,t-1)&1&32&28\\
(7,t-1)^2&1&32&4\\
(41,t-12)&1&32&4\\
(41,t-22)&1&32&4\\
(43,t-20)&1&32&4\\
(43,t-31)&1&32&4\\
(83,t-36)&2&64&4\\
(127,t-94)&3&96&4.
\end{array}
\]
In particular the two nonprincipal classes at conductor $(127,t-94)$
are exchanged by conjugation and together contribute $64$ points.
These are arithmetic ambient points; the theorem makes no new
curve-membership or geometric CM-type assertion.
\end{theorem}
\begin{proof}
The existing local calculations in Theorems~\externalref{supplement}{thm:d7nonmax}
and~\externalref{supplement}{thm:d3classes} leave two embedded orders to consider.
Put
\[
 \mathfrak p'_{41}=(41,t-22),\quad
 \mathfrak p_{127}=(127,t-94),\quad \gamma=B_{10}.
\]
Proposition~\externalref{supplement}{prop:alltraceorders} gives
$(B_5)=\mathfrak p'_{41}$ and $(\gamma)=\mathfrak p_{127}$.
For either prime $\mathfrak p=(p,t-r)$, compute in
\[
 A_{\mathfrak p}=\mathbf F_p[Z]/(Z^2-rZ+1),\qquad
 Q_{\mathfrak p}=A_{\mathfrak p}^{\times}/\mathbf F_p^{\times}.
\]
Exact multiplication of projective pairs $[a:b]$, with
$a^2+rab+b^2\ne0$, gives
\[
\begin{array}{c|l|r|c|r}
\mathfrak p&A_{\mathfrak p}&|Q_{\mathfrak p}|&
 [\varepsilon]&\operatorname{ord}[\varepsilon]\\\hline
\mathfrak p'_{41}&\mathbf F_{41^2}&42&[1:30]&42\\
\mathfrak p_{127}&\mathbf F_{127}\times\mathbf F_{127}&126&[1:80]&42.
\end{array}
\]
Here $\varepsilon=1-z^2$ and $[z]=[\varepsilon]^{22}$ in both
quotients.  At $127$ the roots are $40,54$; the ordered ratio of the
two components sends $\varepsilon$ to $51$ and $1+9z$ to $7$.
For example,
\[
\begin{gathered}
 51^6=32,\quad51^{14}=107,\quad51^{21}=-1,\quad51^{42}=1,\\
 7^{18}=64,\quad7^{42}=107,\quad7^{63}=-1,\quad7^{126}=1
 \pmod{127}
\end{gathered}
\]
certify their respective orders $42$ and $126$.
The complete unit generation $E_M=\langle\varepsilon\rangle\mu_{42}E_L$
proved in Theorem~\externalref{supplement}{thm:d7nonmax} shows that no omitted global
unit enlarges the order-$42$ image.  The conductor exact sequence
and $h(M)=1$ give
\[
 \operatorname{Pic}(S_{\mathfrak p'_{41}})=1,
 \qquad \operatorname{Pic}(S_{\mathfrak p_{127}})\simeq\mathbf Z/3\mathbf Z.
\]
For every proper order here, the unit condition is
$n+22j\equiv0\pmod d$ on $\varepsilon^n z^j v$, $v\in E_L$,
where the image order $d$ is $6$ or $42$.  At $127$ it is $42$,
not $|Q_{\mathfrak p_{127}}|=126$.  The exponent $n$ is therefore
even, so $N(S_{\mathfrak f}^{\times})=E_L^2$; conversely all real
unit squares occur.  For maximal multiplier the image is $E_L^+$.

For completeness of the modules at a prime conductor, normalize a
full local rank-one lattice $I$ so its maximal-order span is $B$.
In the inert case a vector of minimal valuation is a unit of $B$;
in the split case the residue subspace surjects to each coordinate
and contains a vector nonzero in both.  Divide by such a unit.
Then $D+\varpi B\subseteq I\subseteq B$, whose quotient has length
one over the local DVR $D$.  Hence there are exactly the proper
and maximal local types.  Globally this gives
\[
 \operatorname{ICM}(S_{\mathfrak p})=
 \operatorname{Pic}(S_{\mathfrak p})\sqcup\operatorname{Pic}(\OM).
\]
The square-conductor proof in Theorem~\externalref{supplement}{thm:d3classes} adds
exactly the intermediate order at seven.  Thus noninvertible modules
over the originally generated order are included.  Since $h(L)=1$,
each lattice has an integral rank-two basis.  The determinant quotient
$E_L/N(S^{\times})$ supplies $32$ classes for maximal multiplier
and $64$ for each proper Picard class.  These facts prove the first table.

To make the new nonprincipal classes and their root fibers explicit,
use the three distinct conductor-unit cosets
$1,1+9z,1+100z$ at $127$ and the ideals
\[
 J_s=\gamma\OM+(1+sz)\OL,
 \qquad (1+sz,\gamma z)\text{ a basis},\qquad s=0,9,100.
\]
Their residues are units at the conductor; locally they are multiples
of the proper order, and away from it they are maximal.  Thus they
are invertible ideals of the proper multiplier order.  The conductor
sequence identifies their three classes and proves they exhaust it.
The action of $\alpha_k=A_k+B_kz$ is
\[
 T_{k,s}=\begin{pmatrix}
 A_k-sB_k&-\gamma B_k\\
 (B_k/\gamma)(1+st+s^2)&A_k+(t+s)B_k
 \end{pmatrix}.
\]
Its entries are integral and its characteristic polynomial is the
required one; conjugating by $D_{\delta_m}$ supplies all orientations.

Conjugation is inversion on $Q_{\mathfrak p_{127}}$, and
$\overline{1+9z}=85(1+100z)$ modulo the conductor.
Consequently $\overline{J_9}=J_{100}$ and conversely.
The conjugation matrix from the basis of $J_9$ to that of $J_{100}$ is
\[
 C_{100,9}=\begin{pmatrix}
 1+9t&\gamma t\\
 -(109+900t)/\gamma&-1-100t
 \end{pmatrix},\qquad \det C_{100,9}=-1.
\]
It is integral, since $109+900\cdot94\equiv0\pmod{127}$ and
$(\gamma)=(127,t-94)$.  Therefore
$D_{\delta_{m\oplus16}}C_{100,9}D_{\delta_m}^{-1}$ has determinant
one and transports the action on $J_9$ to the conjugate-root action
on $J_{100}$.  The root pairs are
$(J_9,m)\leftrightarrow(J_{100},m\oplus16)$.
All $64$ masks on $J_9$ represent the nonprincipal point population;
the principal class contributes $32$ further points, giving $96$.

For clarity, the nonprincipal conductor-$83$ module may be represented
by $(1+3z,B_2z)$.  In its conductor quotient
$[1+3z]^2=[\varepsilon]^{29}$, and
$\eta=\varepsilon z^{28}$ has norm $u$.
Multiplication by $\eta$ followed by conjugation preserves this ideal
and has determinant $-u$.  Thus its root pairing is
$m\mapsto m\oplus48$.  After twisting, multiply the bridge by the
scalar $u^{-1}$ when bit $32$ of $m$ is unset, and by $1$ otherwise;
the determinant is then one.  Principal modules use $m\mapsto m\oplus16$.
The two portable implementations check all $736$ integral root bridges
against every permitted trace action.

Finally, in any basis $(1+sz,gz)$ put
$q=g\delta$ and $n=1+st+s^2=N_{M/L}(1+sz)>0$.
The exact tuple satisfies
\[
 2x=-\frac{q(t+2s)}n,\qquad
 4h^2=\frac{q^2(4-t^2)}{n^2},\qquad |w|^2=\frac{q^2}n.
\]
It is independent of the allowed trace label.  Together with
Lemma~\ref{lem:packetpointfibers}, this identifies repeated overorders
and proves both the point counts and separation of distinct conductor
strata in the second table.  The fourteen traces acting in the
conductor-seven order give $28$ incidences per point; all $42$ act in
the maximal order.  Each other proper stratum occurs for two traces.
Thus the independent incidence check is
\[
 16\cdot84+32\cdot28+320\cdot4=3520,
 \qquad 16+6\cdot32+64+96=368.
\]
The exact manifests also compare all chosen tuple coefficients, but
quotient distinctness follows from the structural arguments above.
\end{proof}

\begin{remark}[Geometry and representative labels]
The $144$-point and $112$-point populations account for $240$ points of
the full census.  The two further conductor strata contribute
$32+96=128$ points, giving $368$.  Equal rational conductor norms do
not identify the two primes over $41$ or the two over $43$.
The common-basis labels used here need not equal older companion or
ideal-lattice masks, especially at nonprincipal conductor $83$.
An individual geometric decision requires an exact representative
transport; a proved whole-stratum conclusion transfers through the
equality of complete intrinsic populations.  This arithmetic census
adds no geometric CM type, polarization or Hodge conclusion.
\end{remark}


\subsection{From two regions to the full curve}
\begin{lemma}[Two simultaneous target regions]\label{lem:twochambers}
Let $Q$ be the labelled quadrilateral chamber in McMullen's source
half-plane, let $r$ be reflection in its $AB$ side, and let $Q_i$ be
its six normalized coordinate quadrilaterals. Write $r_i$ for reflection
in the labelled $A_iB_i$ side and put
\[
 P_0=\prod_{i=1}^6 Q_i,\qquad P_1=\prod_{i=1}^6 r_i(Q_i).
\]
If a tuple $p\in\HH^6$ represents a point of the normalized McMullen
curve, some $H\in\SL_2(\OL)$ sends $p$ into $P_0\cup P_1$.
Consequently, complete empty searches for both target regions exclude
$p$ from the curve. Each search may use its own explicitly verified
$\SL_2(\OL)$ translate of $p$.
\end{lemma}
\begin{proof}
Let $W$ be the quadrilateral reflection group and let $\Delta_0$ be
its orientation-preserving subgroup. McMullen's construction identifies
the holomorphic quadrilateral group with this index-two subgroup
\cite[\S4]{McMullen2026}. Thus $Q\cup rQ$ is a fundamental region
for $\Delta_0$, with boundary identifications allowed. The normalized
modular embedding $F$ is obtained from the labelled triangle maps by
Schwarz continuation. The same continuation across the $AB$ side gives
$F_i(rz)=r_i(F_i(z))$ in every coordinate
\cite[\S6, Theorem~6.3 and the proof of Theorem~1.9]{McMullen2026}.
It follows that $F(Q)\subset P_0$ and $F(rQ)\subset P_1$.
The region choice is simultaneous in all six coordinates; no independent
coordinate reflections are used.

Write $p=KF(z)$ with $K\in\SL_2(\OL)$. Choose
$\eta\in\Delta_0$ such that $\eta z\in Q\cup rQ$.
Equivariance gives
$\rho(\eta)K^{-1}p=F(\eta z)\in P_0\cup P_1$,
where the normalized monodromy lift $\rho(\eta)$ is integral of
determinant one. If search $j$ uses $p_j=G_jp$ with
$G_j\in\SL_2(\OL)$, then $HG_j^{-1}$ carries $p_j$ to the same
target. Both empty exhaustive searches contradict this necessity.
This proof makes no assertion that any $r_i$ normalizes the arithmetic
group or permutes the packet.
\end{proof}

\begin{theorem}[Complete fixed marked geometric packet]\label{thm:fullpacketgeometry}
Assume the packet execution hypothesis stated in the introduction.
Let $\mathcal S$ be the $368$ ambient quotient points of
Theorem~\ref{thm:fullpacketarithmetic}, with its fixed $M/L$, norm
generator, $L$-action and ordered embeddings. Let $\mathcal X$ be the
normalized McMullen curve image in $\SL_2(\OL)\backslash\HH^6$.
Then
\[
 \mathcal S\cap\mathcal X=\{A,B,C\}.
\]
Thus exactly $365$ points are excluded, including every point of proper
multiplier order. The three labels in the arithmetic certificate are
$\texttt{OM:pic0:07}$, $\texttt{OM:pic0:15}$ and
$\texttt{OM:pic0:00}$ for $A,B,C$, respectively.
\end{theorem}
\begin{proof}
The arithmetic census fixes the complete domain, including all module
classes, determinant sectors, intermediate overorders and root fibers.
The positive vertex realizations in Theorem~\ref{thm:A} and
Theorem~\externalref{supplement}{thm:d7ambient} give three distinct points in the intersection.
It remains to exclude each of the other $365$ points.

For the original target $P_0$, the retained legacy calculations cover
$116$ distinct points. The original-target portable certificate covers
the disjoint complement of $249$ nonvertex labels. Their union is exactly
$\mathcal S\setminus\{A,B,C\}$; the thirteen maximal exclusions are already
inside the $116$ and are not counted twice. The original-target coverage
uses the complete numerical searches, not the full-curve conclusions
of the historical subpacket propositions. At nonprincipal conductor
$83$, its bridge is equality of the entire $32$-point population, not
an unproved permutation of the old and new mask labels.

The reflected certificate reconstructs all $365$ nonvertex labels and
tests the simultaneous target $P_1$. In each target the source transport
is separately verified integral of determinant one. Fresh exact target
vertices and outward rational height/modulus bounds give finite
positive trace-form ellipsoids for all possible entries. Inclusive
lattice recursion, the affine determinant equation, both simultaneous
matrix signs and the separate $c=0$ branch exhaust the candidate set.
Closed side predicates retain boundary equality. These arguments apply
to nonprincipal modules and all orientations, without a maximal-order
or mask-zero hypothesis.

The primary quotient-field path and the separate Sturm path agree on
the complete preparations, entries, rows, candidates and classifications.
Every survivor set is empty. The retained reports bind each plan and
result to its source, target, runtime and full row journal. The
whole-executable correspondence and actual portable controls justify
the source relocation; controls are not a sampling proof for untested
labels. The certificates and historical execution assumptions are
specified in Supplement, Section~\externalref{supplement}{app:computer}.

If a remaining point lay on the curve, Lemma~\ref{lem:twochambers} would
give an integral determinant-one matrix sending it into $P_0$ or $P_1$.
For a source replaced by $G_jp$ in target $j$, the corresponding matrix
is $HG_j^{-1}$. The complete search for that target would retain it,
contradicting its empty result. The positive realizations therefore
give the asserted equality. No arithmetic-normalizer property of the
reflection and no independent per-coordinate chamber choice is used.
\end{proof}


% !TeX root = ../../packet-and-principal-models.tex
\section{The explicit polarization on the normalized Hilbert modular model}\label{sec:polarization}
The preceding arguments use $X_L=\mathbb H^6/\mathrm{SL}_2(\OL)$ as an analytic Hilbert modular quotient. To apply a theorem about polarized abelian varieties we specify the polarization directly rather than infer principal polarization from the notation $X_L$.

For $z=(z_1,\ldots,z_6)\in\mathbb H^6$, let $\sigma_1,\ldots,\sigma_6$ be the real embeddings of $L$ and set
\[
 \Lambda_z=\{(\sigma_i(a)z_i+\sigma_i(b))_{i=1}^6:a,b\in\OL\}\subset\mathbb C^6,
 \qquad A_z=\mathbb C^6/\Lambda_z.
\]
On the abstract lattice $\OL^2$ define
\begin{equation}\label{eq:E1}
 E_1((a,b),(a',b'))=\operatorname{Tr}_{L/\Ql}(a'b-ab').
\end{equation}
This form is integral and alternating. With the displayed sign convention, at one real place, if $v=(a,b)$ and $J_z$ is multiplication by $i$ in the coordinate $az+b$, then
\[
 E_1(v,J_zv)=a^2\,\operatorname{Im}(z)+\frac{(a\operatorname{Re}(z)+b)^2}{\operatorname{Im}(z)}>0
\]
for $v\neq0$. Summing over the six embeddings verifies the Riemann bilinear positivity criterion, so $E_1$ is a polarization. The $\OL$-action is self-adjoint for $E_1$, and every matrix in $\mathrm{SL}_2(\OL)$ preserves \eqref{eq:E1}. Thus this construction descends to the quotient stack $[\mathbb H^6/\mathrm{SL}_2(\OL)]$ (equivalently, to the Hilbert modular orbifold $X_L$) as a polarized RM family.  After passage to any neat finite-index subgroup it becomes an honest abelian scheme over the corresponding finite cover.  In the arguments below only the explicit fiber at $z_B$ is needed.

\begin{proposition}[Polarization type]\label{prop:E1type}
For the integral basis $1,t,\ldots,t^5$ of $\OL$, the polarization $E_1$ has type
\[
 \boxed{(1,7,7,21,21,21)}.
\]
In particular, the polarization used below is not principal. Its degree is
\[
 (1\cdot7\cdot7\cdot21\cdot21\cdot21)^2=\operatorname{disc}(L)^2.
\]
\end{proposition}
\begin{proof}
Let
\[
 T=(\operatorname{Tr}_{L/\Ql}(t^{i+j}))_{0\le i,j\le5}.
\]
In the induced basis of $\OL^2$, the alternating matrix of $E_1$ is
$\left(\begin{smallmatrix}0&-T\\T&0\end{smallmatrix}\right)$. Exact Smith reduction gives
\[
 \operatorname{SNF}(T)=\operatorname{diag}(1,7,7,21,21,21),
 \qquad \det T=3^3 7^5=453789=\operatorname{disc}(L).
\]
The elementary divisors of the alternating form are therefore the six displayed entries, and the polarization degree is the square of their product.
\end{proof}

\begin{proposition}[No compatible principal polarization on the explicit $\OL^2$ lattice]\label{prop:noprincipal}
For the fixed real-multiplication structure on $H_1(A_z,\Z)\cong\OL^2$, there is no principal polarization whose Rosati involution fixes $L$ pointwise.
\end{proposition}
\begin{proof}
An $L$-self-adjoint alternating form on the rank-two $L$-space $L^2$ is necessarily of the form
\[
 E_\lambda((a,b),(a',b'))=\Tr_{L/\Ql}\!\left(\lambda(a'b-ab')\right)
\]
for a unique $\lambda\in L$.  On the lattice $\OL^2$ it is integral exactly when $\lambda\in\mathfrak D_L^{-1}$, and it is unimodular exactly when $\lambda\mathfrak D_L=\OL$.  With $z\in\HH^6$, the Riemann positivity condition is $\sigma_i(\lambda)>0$ for all six real embeddings.  Thus a compatible principal polarization exists precisely when the codifferent has a totally positive generator.

Because $\OL=\Z[t]$ with minimal polynomial $f$, the different is $(f'(t))$.  In the embedding order $k=1,5,11,13,17,19$, the signs of $f'(\sigma_k(t))$ are
\[
 (+,-,+,-,+,-).
\]
Supplement, Section~\externalref{supplement}{app:orders} computes the image of the unit-sign map $E_L\to\{\pm1\}^6$: it is the rank-five subgroup generated by the sign vectors of $\xi_5,\xi_{11},\xi_{13},\xi_{17},\xi_{19}$ and $-1$.  Exact row reduction over $\mathbf F_2$ shows that $(+,-,+,-,+,-)$ is not in this subgroup.  Hence no unit multiple of $1/f'(t)$ is totally positive, so $\mathfrak D_L^{-1}$ has no totally positive generator.
\end{proof}

\paragraph{Homology convention.}
The column $(a,b)$ represents $az+b$.  If a Mobius matrix
$M=\left(\begin{smallmatrix}r&s\\u&v\end{smallmatrix}\right)$ fixes $z$,
then multiplication by $uz+v$ on these columns is $M^{\mathsf T}$, as
proved in Proposition~\ref{prop:fixedmatrixlattice}.  Accordingly, a
Mobius-polynomial representative $S_X$ of $\sqrt{-3}$ acts on homology as
$\Sigma_X=S_X^{\mathsf T}$.  All lattice stability and Hermitian forms
below use $\Sigma_X$.

\begin{proposition}[A principally polarized model in the isogeny class]\label{prop:ppisogenous}
Let $\Lambda=\mathcal O_L^2$ be the homology lattice of $A_B$, let
\[
 \Lambda^\#=\{v\in L^2:E_1(v,\Lambda)\subset\mathbb Z\}
              =\mathfrak D_L^{-1}\Lambda,
\]
and let $\Sigma_B=(I+2B_t^{14})^{\mathsf T}$.  There is an $\mathcal O_L$-submodule
\[
 H\subset \Lambda^\#/\Lambda
\]
which is maximal isotropic for the discriminant pairing and stable under the action of $\mathcal O_K$.  Here $\sqrt{-3}$ acts as $\Sigma_B$, while the integral generator
\[
 U:=(B_t^{14})^{\mathsf T}=\frac{\Sigma_B-I}{2}
\]
satisfies $U^2+U+I=0$ and hence identifies $\mathcal O_K$ with $\mathbb Z[U]$.  If $\Lambda^{\mathrm{pp}}$ denotes the inverse image of $H$ in $\Lambda^\#$, then
\[
 [\Lambda^{\mathrm{pp}}:\Lambda]=\operatorname{disc}(L)=453789,
 \qquad (\Lambda^{\mathrm{pp}})^\#=\Lambda^{\mathrm{pp}}.
\]
Consequently
\[
 A_B^{\mathrm{pp}}:=\mathbb C^6/\Lambda_B^{\mathrm{pp}}
\]
is a principally polarized abelian sixfold, the natural map
\[
 A_B\longrightarrow A_B^{\mathrm{pp}}
\]
is an $\mathcal O_L$-linear isogeny of degree $453789$, and the $K$-action descends to $A_B^{\mathrm{pp}}$.  The rational Hodge structure, $K$-signature $(3,3)$, Weil discriminant $-1$, and validity of the Hodge conjecture are unchanged under this isogeny.
\end{proposition}
\begin{proof}
Since $\mathcal O_L=\mathbb Z[t]$, Dedekind factorization applied to~\eqref{eq:minpoly} gives
\[
 f(T)\equiv (T^3-T^2+T+1)^2\pmod 3,
 \qquad
 f(T)\equiv (T-1)^6\pmod 7.
\]
Both ramification indices are prime to the residue characteristic, hence the ramification is tame.  Thus
\[
 (3)=\mathfrak p_3^2,\quad N\mathfrak p_3=3^3,
 \qquad
 (7)=\mathfrak p_7^6,\quad N\mathfrak p_7=7,
\]
and the different is
\[
 \mathfrak D_L=\mathfrak p_3\mathfrak p_7^5.
\]
Multiplication by a generator of the different identifies the discriminant module with
\[
 \Lambda^\#/\Lambda\cong
 (\mathcal O_L/\mathfrak D_L)^2
 \cong
 (\mathcal O_L/\mathfrak p_3)^2\oplus
 (\mathcal O_L/\mathfrak p_7^5)^2.
\]

At $\mathfrak p_3$, exact reduction of the displayed matrix $\Sigma_B$ gives
\[
 \Sigma_B\equiv0\pmod{\mathfrak p_3}.
\]
Hence any one-dimensional $\mathcal O_L/\mathfrak p_3\cong\mathbb F_{27}$ subspace is $\Sigma_B$-stable; take the first-coordinate line $H_3$.  Since $U=(\Sigma_B-I)/2\equiv1\pmod{\mathfrak p_3}$, this line is already $\mathcal O_K=\mathbb Z[U]$-stable.

At $\mathfrak p_7^5$, the quotient has characteristic $7$ and $\Sigma_B^2=-3I=4I$.  Moreover $U^2+U+I=0$ gives $\operatorname{tr}(U)=-1$, hence
\[
 \operatorname{tr}(\Sigma_B)=\operatorname{tr}(I+2U)=0.
\]
Since $4$ is a unit, the idempotents
\[
 e_+=\frac{\Sigma_B+2I}{4},\qquad e_-=\frac{2I-\Sigma_B}{4}
\]
split $(\mathcal O_L/\mathfrak p_7^5)^2$ into the $+2$ and $-2$ eigensummands.  Because the local ring $\mathcal O_L/\mathfrak p_7^5$ is local, the images of these idempotents are finite projective and hence free.  Moreover
\[
 \operatorname{tr}(e_+)=\frac{\operatorname{tr}(\Sigma_B)+4}{4}=1,
 \qquad
 \operatorname{tr}(e_-)=1,
\]
so both images have rank one.  Let $H_7=\operatorname{im}(e_+)$.  Because $\Sigma_B^\dagger=-\Sigma_B$, for $v,w\in H_7$ one has
\[
 0=E_1(\Sigma_Bv,w)+E_1(v,\Sigma_Bw)=4E_1(v,w)
\]
in the local discriminant pairing, so $H_7$ is isotropic.  The line $H_3$ is isotropic for the same alternating pairing.  By the Chinese remainder theorem
\[
 H=H_3\oplus H_7
\]
is isotropic and $\mathcal O_L$-stable, and
\[
 |H|=3^3\,7^5=453789=\operatorname{disc}(L)
       =\sqrt{|\Lambda^\#/\Lambda|}.
\]
Thus $H$ is maximal isotropic and its inverse image $\Lambda^{\mathrm{pp}}$ is self-dual for $E_1$.

The action of $\Sigma_B$ preserves $H$.  The full ring $\mathcal O_K=\mathbb Z[U]$ preserves it as well.  Indeed $U=(B_t^{14})^{\mathsf T}$ preserves both $\Lambda$ and $\Lambda^\#$; on the discriminant module one also has $U=(\Sigma_B-I)/2$, and multiplication by $2$ is invertible on both the $3$- and $7$-primary parts, so $\Sigma_B$-stability of $H$ implies $U$-stability.  Hence $\Lambda^{\mathrm{pp}}$ is $\mathcal O_K$-stable.  The complex structure, the rational $K$-action on $V:=\Lambda\otimes\Ql=L^2$, and the rational alternating form $E_1$ have not changed.  Hence the $K$-Hermitian form $H(v,w)=E_1(\Sigma_Bv,w)+\sqrt{-3}\,E_1(v,w)$ on the rational $K$-space $V$ is literally the same form before and after the lattice enlargement.  In particular its determinant class in $\Ql^\times/N_{K/\Ql}(K^\times)$ remains $-1$, not merely up to an appeal to isogeny invariance.  The CM type and all rational Hodge classes are likewise unchanged.  The inclusion $\Lambda\subset\Lambda^{\mathrm{pp}}$ has index $|H|$, which is the degree of the induced isogeny.  Finally the Hodge conjecture is invariant under isogeny (pullback and pushforward along a finite isogeny preserve algebraicity and induce inverse isomorphisms on rational cohomology up to the degree).
\end{proof}

\begin{corollary}[The principal model is not the same orbifold vertex]\label{cor:noBstablepp}
No principally polarized self-dual overlattice
\[
 \Lambda\subset\Lambda'\subset\Lambda^\#
\]
can be stable under the full vertex automorphism acting as $B_t^{\mathsf T}$ on homology.  In particular the principally polarized model of Proposition~\ref{prop:ppisogenous} cannot be obtained by merely changing the polarization while retaining the same integral order-$42$ $B_t$-vertex structure.
\end{corollary}
\begin{proof}
Every self-dual overlattice corresponds to a maximal isotropic subgroup,
whose $\mathfrak p_3$-part has order $27$. Stability under the actual
homology automorphism $B_t^{\mathsf T}$ implies stability under its inverse
and therefore under their sum $tI$. Hence this subgroup is an
$\mathbb F_{27}$-line in $(\mathcal O_L/\mathfrak p_3)^2$.  The characteristic polynomial of $B_t$ is $X^2-tX+1$.  Its discriminant $t^2-4$ is a nonsquare in $\mathbb F_{27}$: exact reduction modulo $g_3(T)=T^3-T^2+T+1$ gives
\[
 (t^2-4)^{13}=-1.
\]
Hence $X^2-tX+1$ is irreducible over $\mathbb F_{27}$ and $B_t$ has no invariant $\mathbb F_{27}$-line.  Thus no maximal isotropic subgroup, and hence no principally polarized overlattice, can be $B_t^{\mathsf T}$-stable.
\end{proof}

\begin{remark}[From special fibers to a common principal family]\label{rem:ppcover}
Proposition~\ref{prop:commonprincipal} below supplies a common nonfree
maximal isotropic subgroup for the three specialized $K$-actions.
Proposition~\ref{prop:principalstabilizer} determines its ambient index
$224$ and its index $28$ in the explicitly parametrizing quadrilateral
group, before auxiliary level. Theorem~\ref{thm:principalfamily} descends
its inverse-image lattice sheaf and principal polarization over a connected
neat cover. Proposition~\ref{prop:survivingorder} determines the surviving
order inside each selected CM field. The value $537824$ in the earlier
free-line construction concerns a different subgroup and is not the
index of this common nonfree level.
\end{remark}

\begin{remark}[What the principally polarized model does not prove]\label{rem:pphec}
The special-fiber subgroup $H$ gives a principally polarized model
$A_B^{\mathrm{pp}}$; Remark~\ref{rem:ppcover} describes the later common
choice used for the family. Corollary~\ref{cor:noBstablepp}
shows that any such modification necessarily loses the integral order-$42$
vertex automorphism at the prime above $3$, so it would not be the original
orbifold $B$-vertex in another polarization disguise.  We do not identify
$A_B^{\mathrm{pp}}$ with a point on Mostaed's original curve $V$ inside the
principally polarized Hilbert--Blumenthal moduli space.
\end{remark}

Thus the Hodge-theoretic conclusion below does not require a map from $X_L$ to the principally polarized moduli space $\mathcal A_6$: it uses the explicit polarized abelian variety $(A_z,E_1)$ at the certified vertex.  In fact Proposition~\ref{prop:noprincipal} shows that, for this fixed $\OL$-module and real-multiplication action, the missing principal polarization is not merely unconstructed.  Markman's theorem is stated for polarized abelian sixfolds, not only principally polarized ones.

At the normalized $B$-vertex put
\[
 U=(B_t^{14})^{\mathsf T},\qquad \Sigma_B=I+2U.
\]
Since $B_t^{\mathsf T}\in\mathrm{SL}_2(\OL)$, its $E_1$-adjoint is $(B_t^{-1})^{\mathsf T}$. The elliptic relation gives $U^2+U+I=0$, and hence
\[
 \Sigma_B^2=-3I,\qquad \Sigma_B^\dagger=I+2U^{-1}=I+2U^2=-\Sigma_B,\qquad \Sigma_B^\dagger \Sigma_B=3I.
\]

\begin{proposition}[Polarization-compatible $K$-action]\label{prop:Kcompatible}
The assignment $\sqrt{-3}\mapsto \Sigma_B$ defines an embedding
\[
 K=\Ql(\sqrt{-3})\hookrightarrow\operatorname{End}^0(A_z)
\]
at the certified $B_t$-vertex and is compatible with $E_1$: for every $k\in K$,
\[
 \eta(k)^\dagger\eta(k)=N_{K/\Ql}(k)I.
\]
\end{proposition}
\begin{proof}
The identity $\Sigma_B^2=-3I$ gives the field embedding. Since the vertex is fixed by $B_t$, the actual homology action is $B_t^{\mathsf T}$; every polynomial in it, in particular $\Sigma_B$, preserves the complex structure. Integrality gives an actual endomorphism. The relation $\Sigma_B^\dagger=-\Sigma_B$ gives
\[
 (a-b\Sigma_B)(a+b\Sigma_B)=a^2+3b^2=N_{K/\Ql}(a+b\sqrt{-3}),
\]
and the claim follows for all $a,b\in\Ql$.
\end{proof}



% !TeX root = ../../packet-and-principal-models.tex
\section{The three geometric fibers}
We now supply the positive realizations and Hodge data used above. The B
companion bridge fixes the geometric marking explicitly. The A and C
calculations below use their actual vertex actions; their CM types are not
inferred from an arithmetic companion signature.

\subsection{The arithmetic certificate used in Theorem~\ref{thm:A}}\label{sec:maincertificate}
Only a small part of the $\ell=43$ arithmetic is needed for the main theorem.  The full norm-fiber classification, signature enumeration, conjugacy reduction, component analysis, and unit calculation are deferred to Supplement, Section~\externalref{supplement}{app:arithdetail}.

Put
\[
 \pi=t^4+t^3-4t^2-2t+5,
 \qquad
 \alpha_0=1+\zeta_{42}+\zeta_{42}^4,
 \qquad
 \alpha_{11}=\zeta_{42}^{11}\alpha_0,
\]
and $c_{11}=\alpha_{11}+\overline{\alpha}_{11}$.  Define
\[
 A_{c_{11}}=\begin{pmatrix}0&-\pi\\1&c_{11}\end{pmatrix},
 \qquad
 B=\begin{pmatrix}0&-1\\1&t\end{pmatrix}.
\]

\begin{proposition}[Arithmetic and component certificate]\label{prop:mainarith}
The following assertions hold.
\begin{enumerate}[label=(\roman*)]
\item $N_{M/L}(\alpha_0)=\pi$, $N_{L/\Ql}(\pi)=43$, and $\pi$ is totally positive.  Thus $A_{c_{11}}$ is an integral determinant-$\pi$ Hecke matrix with characteristic polynomial $X^2-c_{11}X+\pi$.
\item The six Galois triangles in McMullen's construction have orientations, relative to the identity component and ordered by $k=1,5,11,13,17,19$,
\[
 (+,+,-,-,-,-).
\]
The unit $t$ has the same sign vector.  Hence
\[
 h_t=\operatorname{diag}(t,1),\qquad B_t=h_tBh_t^{-1}
\]
changes all six components to $\mathbb H$ and normalizes $\mathrm{SL}_2(\OL)$.
\item The point defined by $A_{c_{11}}$ belongs to the maximal-order branch $\OL[\alpha_{11}]=\OM$; no additional integral-conjugacy or unit-norm ambiguity is required for the exact certificate of Section~\ref{sec:vertexhit}.
\end{enumerate}
\end{proposition}
\begin{proof}
Part (i) is the exact cyclotomic calculation of Supplement, Section~\externalref{supplement}{app:prime43}.  For (ii), McMullen constructs the component maps by comparing the orientations of the six Galois triangles and sends an oppositely oriented component to $-\mathbb H$ while preserving the equivariance relation; see \cite[Prop.~6.1, Cor.~6.2, Thm.~6.3 and Eq.~(6.1)]{McMullen2026}.  The six signs are certified in Supplement, Section~\externalref{supplement}{app:components}; $N_{L/\Ql}(t)=1$, so $t\in\OL^\times$ and conjugation by $h_t$ preserves $\mathrm{SL}_2(\OL)$.  Part (iii) follows from the exact order-discriminant and unit-norm calculation in Supplement, Section~\externalref{supplement}{app:orders}.
\end{proof}

This proposition is deliberately a certificate, not a reconstruction of the entire search.  The proof of Theorem~\ref{thm:A} from this point onward uses only the explicit matrices and the polarized Hodge structure at their common vertex.

\begin{proposition}[Geometric meaning of the normalized $B$-vertex]\label{prop:geometricB}
Let $F:\HH\to(\pm\HH)^6$ be McMullen's modular embedding constructed from the six Galois triangles, with the identity embedding as the first component.  Let $z_B\in\HH$ be the vertex fixed by the order-$42$ generator $B$ in the identity triangle and put
\[
 F_t=h_t\circ F,\qquad h_t=\operatorname{diag}(t,1).
\]
Then $F_t$ takes values in $\HH^6$, is equivariant for the conjugated integral representation $g\mapsto h_tgh_t^{-1}$, and $F_t(z_B)$ is the componentwise upper-half-plane fixed point of
\[
 B_t=h_tBh_t^{-1}.
\]
Consequently the image of $F_t(z_B)$ in $X_L$ is a point of the normalized McMullen geodesic, independently of any Hecke calculation.
\end{proposition}
\begin{proof}
McMullen's Theorem~6.3 constructs each component $f_i$ by the unique vertex-respecting conformal map from the identity triangle to the $i$-th Galois triangle.  When the two orientations agree, $f_i$ takes values in $\HH$; when they disagree, it takes values in $-\HH$, and in both cases
\[
 f_i(g_1z)=g_i f_i(z)
\]
for every triangle-group element $g$ \cite[Thm.~6.3, Eq.~(6.1)]{McMullen2026}.  In particular $f_i(z_B)$ is the vertex fixed by $B_i$.  Proposition~\ref{prop:mainarith} gives the component signs $(+,+,-,-,-,-)$ and the identical sign vector for the real determinants $\sigma_i(t)$ of $h_t$.  Since
\[
 \operatorname{Im}(h_{t,i}w)=\frac{\sigma_i(t)\operatorname{Im}(w)}{|c_iw+d_i|^2}
\]
for a real fractional-linear map of determinant $\sigma_i(t)$, $h_t$ carries the raw component of $F$ to $\HH^6$.  As $t\in\OL^\times$, conjugation by $h_t$ preserves $\mathrm{SL}_2(\OL)$.  Finally, conjugating the equation $B_i f_i(z_B)=f_i(z_B)$ by $h_{t,i}$ shows that the $i$-th coordinate of $F_t(z_B)$ is fixed by $(B_t)_i$.
\end{proof}

\subsection{An exact $k=11$ vertex intersection}\label{sec:vertexhit}
The finite orbit search described above contains a particularly rigid case.  It is not merely a numerical near-hit: after the integral component normalization $h_t=\operatorname{diag}(t,1)$, one of the $k=11$ Hecke fixed points is exactly the elliptic $B$-vertex of McMullen's normalized fundamental quadrilateral.

Let
\[
 B=\begin{pmatrix}0&-1\\1&t\end{pmatrix},\qquad B_t=h_tBh_t^{-1},
\]
and let $A_{c_{11}}=\left(\begin{smallmatrix}0&-\pi\\1&c_{11}\end{smallmatrix}\right)$ be the companion representative associated with $\alpha_{11}=\zeta_{42}^{11}\alpha_0$.  Define
\[
 G=\begin{pmatrix}g_{11}&g_{12}\\g_{21}&g_{22}\end{pmatrix}
\]
with
\begin{align*}
 g_{11}&=2t^5-t^4-9t^3+3t^2+7t+1,\\
 g_{12}&=(t^2-t-1)(t^3-4t-1),\\
 g_{21}&=-2(t^5-5t^3-t^2+5t+3),\\
 g_{22}&=-(t-2)(3t^4+7t^3-4t^2-16t-7).
\end{align*}

\begin{theorem}[Exact $B_t$-vertex certificate]\label{thm:Bvertex}
In $M_2(L)$ one has
\[
 \det G=1,
 \qquad
 [GA_{c_{11}}G^{-1},B_t]=0.
\]
More precisely,
\[
 GA_{c_{11}}G^{-1}=xI+yB_t
\]
with
\[
 x=(t^2-2)(t^3-t^2-6t-2),\qquad
 y=t(t^2-2)(t^2+t-1).
\]
Consequently the Hecke fixed point represented by $A_{c_{11}}$ is
$\mathrm{SL}_2(\OL)$-equivalent to the geometrically defined point $F_t(z_B)$ of Proposition~\ref{prop:geometricB}.  In particular the corresponding quotient point lies on the normalized McMullen curve in $X_L$.
\end{theorem}
\begin{proof}
All identities are polynomial identities in $L=\Ql[t]/(f)$.  Direct reduction modulo $f$ gives $\det G=1$ and the displayed linear expression in $B_t$.  Since $B_t$ is elliptic, every non-scalar element of the $L$-algebra $L[B_t]$ has the same pair of fixed points at each real embedding.  The chosen upper-half-plane fixed point is therefore carried by $G$ from the companion Hecke fixed point to the $B_t$ vertex.  The accompanying exact certificate performs no floating-point comparison.
\end{proof}


\subsection{From fixed matrices to endomorphisms of the explicit lattice}\label{sec:latticetransport}
The preceding vertex certificate is an identity in the Hilbert modular matrix model.  We now record explicitly why it is also an identity of endomorphisms of the polarized abelian variety constructed above in Section~\ref{sec:polarization}.  This removes a possible ambiguity between a Möbius fixed-point computation and the actual rational Hodge structure.

Let
\[
 A=\begin{pmatrix}a&b\\ c&d\end{pmatrix}\in M_2(\OL)
\]
and suppose that $z=(z_i)\in\mathbb H^6$ is fixed componentwise by $A$, so that
\[
 z_i=\frac{\sigma_i(a)z_i+\sigma_i(b)}{\sigma_i(c)z_i+\sigma_i(d)}.
\]
Put $\lambda_i=\sigma_i(c)z_i+\sigma_i(d)$ and let $m_\lambda$ denote the diagonal complex-linear map with multipliers $\lambda_i$.

\begin{proposition}[Fixed matrices act on the explicit lattice]\label{prop:fixedmatrixlattice}
The map $m_\lambda$ preserves $\Lambda_z$ and hence induces an endomorphism of $A_z$.  Relative to the $\OL$-basis $(z,1)$ of the rank-two $\OL$-module defining $\Lambda_z$, its coordinate matrix is $A^{\mathsf T}$.  If $\det_L(A)=\delta\ne0$, then
\[
 [\Lambda_z:m_\lambda(\Lambda_z)]=|N_{L/\Ql}(\delta)|,
\]
so the induced isogeny has degree $|N_{L/\Ql}(\delta)|$.  In particular, $A_{c_{11}}$ induces a degree-$43$ endomorphism, while $B_t$ induces an automorphism.
\end{proposition}
\begin{proof}
The fixed-point equation gives
\[
 \lambda_i z_i=\sigma_i(a)z_i+\sigma_i(b),\qquad
 \lambda_i=\sigma_i(c)z_i+\sigma_i(d).
\]
Thus multiplication by $\lambda$ sends the two $\OL$-generators $z$ and $1$ to the $\OL$-linear combinations encoded by the columns of $A^{\mathsf T}$, proving lattice preservation and the matrix assertion.  The index of the image of the rank-two $\OL$-lattice under an $L$-linear matrix is the absolute norm of its $L$-determinant.  For $A_{c_{11}}$ this is $N_{L/\Ql}(\pi)=43$, while $\det B_t=1$.
\end{proof}

Now let $\gamma=\left(\begin{smallmatrix}r&s\\u&v\end{smallmatrix}\right)\in\mathrm{SL}_2(\OL)$ and put $z'=\gamma z$.  The standard analytic isomorphism
\[
 q_\gamma:A_z\longrightarrow A_{z'},\qquad
 (w_i)_i\longmapsto\left(\frac{w_i}{\sigma_i(u)z_i+\sigma_i(v)}\right)_i
\]
is an isomorphism of lattices.  On the abstract $\OL^2$ coordinates it is represented by $\gamma^{-\mathsf T}$.  Since $\det\gamma=1$, it preserves the determinant pairing and hence the trace polarization $E_1$.

\begin{corollary}[The exact $G$-certificate transports the endomorphism]\label{cor:Gtransport}
Under the polarized isomorphism $q_G$ from the $k=11$ companion fixed point to the normalized $B_t$-vertex, the endomorphism represented by $A_{c_{11}}^{\mathsf T}$ is carried to
\[
 (GA_{c_{11}}G^{-1})^{\mathsf T}=(xI+yB_t)^{\mathsf T}.
\]
Consequently the field generated by the corrected Hecke endomorphism is identified inside the actual endomorphism algebra of the certified vertex with the field generated by $B_t^{\mathsf T}$. The $K$-operator $\Sigma_B=(I+2B_t^{14})^{\mathsf T}$ used below is therefore an endomorphism of the same polarized rational Hodge structure appearing in the exact Hecke certificate.
\end{corollary}
\begin{proof}
In lattice coordinates, conjugation by $q_G$ sends $A^{\mathsf T}$ to
\[
 G^{-\mathsf T}A^{\mathsf T}G^{\mathsf T}=(GAG^{-1})^{\mathsf T}.
\]
Apply Theorem~\ref{thm:Bvertex}.  Preservation of $E_1$ follows from $G\in\mathrm{SL}_2(\OL)$.
\end{proof}


\subsection{The actual CM type at the vertex}\label{sec:actualcmtype}

There is a small orientation point here which is essential.  McMullen's construction
associates to the $i$-th real embedding a triangle $T_i\subset\mathbb H$; when the
orientation of $T_i$ agrees with $T_1$, the component map lands in $\mathbb H$, and
when the orientations disagree it lands in $-\mathbb H$ \cite[Thm.~6.3 and Eq.~(6.1)]{McMullen2026}.
Using McMullen's printed matrices $A,B,C$ and their six real embeddings, the oriented
vertex triples have relative orientation signs
\[
 (+,+,-,-,-,-)
\]
for $k=1,5,11,13,17,19$.  These are exactly the signs of
$\sigma_k(t)$, where $t=2\cos(\pi/21)$.  Thus the fixed $B$-vertex contributes the
eigenvalue $\zeta^{k}$ in the first two components and $\zeta^{-k}$ in the last four.
This gives a geometric derivation of the CM-type signs below, independent of the
normalizing matrix $h_t$ alone.  The signs are certified by rational isolating intervals
for all six real roots of the minimal polynomial of $t$ and exact rational interval
evaluation of the three fixed points; no floating-point sign decision enters the proof.

\begin{remark}[The caption of McMullen's Figure~5]\label{rem:fig5signs}
McMullen's general construction is explicitly formulated in
$(\pm\mathbb H)^6$ and sends an oppositely oriented Galois triangle to
$-\mathbb H$ \cite[Thm.~6.3, Eq.~(6.1)]{McMullen2026}.  The caption of Figure~5
writes $\widetilde f_0:\mathbb H\to\mathbb H^6$.  For the printed matrix model,
however, the direct certified orientation calculation above gives the mixed sign vector
$(+,+,-,-,-,-)$.  We therefore use the matrices and Theorem~6.3, rather than the
caption alone, to determine the raw component of $Y_L$.  Proposition~\externalref{supplement}{prop:orientationt}
then transfers that raw component to $X_L$ by the integral matrix $h_t$.
\end{remark}

\begin{proposition}[CM type and Weil signature]\label{prop:actualcmtype}
The CM type of the vertex in Theorem~\ref{thm:Bvertex} is
\[
 \Phi=\{1,5,31,29,25,23\}\subset (\Z/42\Z)^\times.
\]
For the quadratic subfield $K=\Ql(\sqrt{-3})=\Ql(\zeta_3)$, the restrictions of these six embeddings have residues
\[
 1,2,1,2,1,2\pmod 3.
\]
Hence the $K$-action on $H^{1,0}$ has multiplicities $(3,3)$: the certified vertex is an abelian sixfold of $K$-Weil type.
\end{proposition}
\begin{proof}
The sign statement above selects the exponents $1,5,-11,-13,-17,-19$ modulo $42$, which are exactly the displayed set.  Restriction to $K=\Ql(\zeta_3)$ is determined by the exponent modulo $3$, giving three copies of each of the two embeddings of $K$.
\end{proof}


\subsection{The Weil discriminant and Markman's theorem}\label{sec:markman}
Let $s=\sqrt{-3}\in K$, acting as $\Sigma_B$. With respect to $E_1$, define
\begin{equation}\label{eq:HermitianE1}
 H(v,w)=E_1(\Sigma_B v,w)+sE_1(v,w).
\end{equation}
The identities $\Sigma_B^\dagger=-\Sigma_B$ and $\Sigma_B^2=-3I$ imply $H(\Sigma_B v,w)=sH(v,w)$ and $H(v,\Sigma_B w)=\bar sH(v,w)$, so \eqref{eq:HermitianE1} is a $K$-Hermitian form for the polarized Weil structure (up to the conventional nonzero rational normalization, which does not change the determinant class in rank six).

Let $\eta\in L$ be the lower-left entry of the Mobius representative $S_B=I+2B_t^{14}$. Exact reduction gives
\[
 \eta=2(t^2-2)(2t^3+t^2-7t-4),\qquad N_{L/\Ql}(\eta)=-1728.
\]
The six vectors $f_j=(0,t^j)$, $0\le j\le5$, form a $K$-basis of $L^2$. Since $E_1(f_i,f_j)=0$, the Hermitian matrix in this basis is
\[
 \bigl(-\operatorname{Tr}_{L/\Ql}(\eta t^{i+j})\bigr)_{0\le i,j\le5}.
\]
The basis assertion follows because $\eta\ne0$: the vectors $f_j,\Sigma_Bf_j$ span $L^2$ over $\Ql$. Since the rank is six, the minus sign does not change the determinant. Therefore
\begin{equation}\label{eq:vertexdisc}
 \det H=\operatorname{disc}(L)N_{L/\Ql}(\eta)
 =(3^3 7^5)(-1728)=-2^6 3^6 7^5.
\end{equation}
Every identity here is exact in $\Ql[t]/(f)$.

\begin{proposition}[Weil discriminant $-1$]\label{prop:discminusone}
For the explicit polarization $E_1$, the $K$-Hermitian discriminant of the sixfold of Theorem~\ref{thm:Bvertex} is the class of $-1$ in
\[
 \Ql^\times/N_{K/\Ql}(K^\times).
\]
\end{proposition}
\begin{proof}
Equation~\eqref{eq:vertexdisc} represents the determinant class. Now
\[
 2^6=N_{K/\Ql}(8),\qquad 3^6=N_{K/\Ql}(27),\qquad 7=N_{K/\Ql}(2+\sqrt{-3}).
\]
Therefore $2^6 3^6 7^5$ is a $K/\Ql$ norm, leaving precisely the class $-1$.
\end{proof}

Markman's Section~1.1 defines a polarized abelian variety of Weil type as a triple $(A,\eta,h)$ in which $h\in H^{1,1}(A,\Ql)$ is an ample class satisfying $\eta(k)^*h=N_{K/\Ql}(k)h$; no principality hypothesis is imposed \cite[Sec.~1.1]{Markman2025}.  In the present example the hypotheses match term by term: $\dim A_B=6$; Proposition~\ref{prop:actualcmtype} gives multiplicities $(3,3)$; $E_1$ is an ample integral, hence rational, $(1,1)$-class; Proposition~\ref{prop:Kcompatible} gives the required norm-similitude; and Proposition~\ref{prop:discminusone} gives discriminant class $-1$ for the associated $K$-Hermitian form.  Markman's Theorem~1.5.1 therefore applies exactly and states that the Hodge--Weil classes are algebraic \cite[Thm.~1.5.1]{Markman2025}.  The theorem does not exclude CM points.

\begin{corollary}[Algebraicity of the exceptional Weil space]\label{cor:Weilalg}
For the $k=11$ vertex of Theorem~\ref{thm:Bvertex}, the two-dimensional rational Hodge--Weil space
\[
 W_K(A)\subset H^6(A,\Ql)\cap H^{3,3}(A)
\]
is algebraic.
\end{corollary}
\begin{proof}
Proposition~\ref{prop:actualcmtype} gives Weil type for $K=\Ql(\sqrt{-3})$, Proposition~\ref{prop:discminusone} gives discriminant $-1$, and Markman's theorem applies.
\end{proof}


\subsection{All Hodge classes on the $k=11$ CM sixfold}\label{sec:fullhodge}
We finally determine the Hodge classes of this particular CM sixfold.  This is stronger than merely identifying the exceptional Weil subspace.

Let $G=(\Z/42\Z)^\times=\Gal(M/\Ql)$ and let $e_a$ denote the one-dimensional complex eigenspace of $H^1(A,\C)$ indexed by $a\in G$.  For a subset $S\subset G$ put $e_S=\bigwedge_{a\in S}e_a$.  If $|S|=2p$, the standard CM-torus character criterion says that $e_S$ is a Hodge line in $H^{2p}(A,\C)$ precisely when
\begin{equation}\label{eq:CMcriterion}
 |gS\cap\Phi|=p\qquad\text{for every }g\in G.
\end{equation}
To make the character calculation explicit, let $T=\operatorname{Res}_{M/\Ql}\mathbf G_m$ and let $\mu$ be the Hodge cocharacter of the CM type $\Phi$.  The Mumford--Tate torus is the smallest $\Ql$-subtorus of $T$ through which all Galois conjugates $g\mu$ factor.  If $\chi_S$ is the $T_\C$-character of $e_S$, then
\[
 \langle\chi_S,g\mu\rangle=|gS\cap\Phi|.
\]
For $|S|=2p$, the line $e_S$ has Hodge type $(p,p)$ exactly when this pairing equals the weight $p$ for every Galois conjugate of the Hodge cocharacter.  Hence
\[
 |gS\cap\Phi|=p\quad\text{for every }g\in G,
\]
which is \eqref{eq:CMcriterion}.  Equivalently, after subtracting the central weight cocharacter, the character $\chi_S$ is trivial on the Hodge group.  This is the elementary character form of the CM-torus description used in the general theory of Hodge classes on CM abelian varieties; compare \cite{MilneHodge2020}.

Complex conjugation on $G$ is $a\mapsto-a$.  The six conjugate pairs are
\[
 \{1,41\},\{5,37\},\{11,31\},\{13,29\},\{17,25\},\{19,23\}.
\]
Each corresponding degree-two Hodge line is a divisor line, and products of divisors correspond to unions of these pairs.

\begin{lemma}[Fourier classification for the primitive CM type]\label{lem:primitivebinary}
Let
\[
 \Phi=\{1,5,23,25,29,31\}\subset G
\]
and let $S\subset G$ have cardinality $2p$.  Then
\[
 |gS\cap\Phi|=p\quad\text{for every }g\in G
\]
if and only if either $S$ is a union of $p$ antipodal pairs
$\{a,-a\}$, or $p=3$ and
\[
 S=\{a\in G:a\equiv1\pmod3\}
 \quad\text{or}\quad
 S=\{a\in G:a\equiv2\pmod3\}.
\]
Consequently the numbers of solutions for $p=1,\ldots,6$ are
\[
 \binom6p+2\delta_{p,3}=6,15,22,15,6,1.
\]
The same classification holds for every Galois translate of $\Phi$.
\end{lemma}
\begin{proof}
Use the Chinese-remainder coordinates
\[
 G\simeq C_2\times C_6,\qquad
 a\bmod3=(-1)^e,\qquad a\bmod7=3^n,
\]
where $e\in\{0,1\}$ and $n\in\mathbb Z/6\mathbb Z$.  In these coordinates
complex conjugation $a\mapsto-a$ is translation by $(1,3)$.  Put
$\varepsilon=2\mathbf 1_\Phi-1$.  With columns indexed by $n=0,\ldots,5$,
direct CRT reduction gives
\[
 \bigl(\varepsilon(0,n)\bigr)_n=(1,1,-1,-1,1,-1),
 \qquad
 \bigl(\varepsilon(1,n)\bigr)_n=(1,-1,1,-1,-1,1).
\]
In particular $\varepsilon(x+(1,3))=-\varepsilon(x)$.

Let $\zeta_6$ be a primitive sixth root of unity and use the Fourier
convention
\[
 \widehat f(r,s)=\sum_{e=0}^1\sum_{n=0}^5
 f(e,n)(-1)^{re}\zeta_6^{-sn},
 \qquad r\in\mathbb Z/2\mathbb Z,\quad s\in\mathbb Z/6\mathbb Z.
\]
The character $\chi_{r,s}$ is odd under complex conjugation exactly when
$r+s$ is odd.  Since $\varepsilon$ is odd, its even Fourier coefficients
vanish.  For the six odd characters, summing the two displayed rows and using
$\zeta_6^2-\zeta_6+1=0$ gives the exact table
\[
\begin{array}{c|cccccc}
 (r,s)&(0,1)&(0,3)&(0,5)&(1,0)&(1,2)&(1,4)\\ \hline
 \widehat\varepsilon(r,s)
 &4&4&4&0&-4(2\zeta_6-1)&4(2\zeta_6-1).
\end{array}
\]
Thus the only odd character missing from the Fourier support of
$\varepsilon$ is
\[
 \eta=\chi_{1,0},\qquad \eta(e,n)=(-1)^e,
\]
the quadratic character modulo $3$.

Let $F=\mathbf1_S$ and define the correlation
\[
 C_F(g)=\sum_{x\in G}F(x)\varepsilon(g^{-1}x)
       =2|S\cap g\Phi|-|S|.
\]
The Hodge condition is equivalent to $C_F=0$, since $g$ ranges through all
of $G$.  Character orthogonality, with the convention above, gives
\[
 \widehat C_F(\chi)=\widehat F(\chi)
                    \widehat\varepsilon(\chi^{-1}).
\]
The odd support in the table is invariant under inversion.  Therefore every
odd Fourier coefficient of $F$ vanishes except possibly the $\eta$-coefficient.
Equivalently the odd part of $F$ has the form
\[
 F^-(x):=\frac{F(x)-F(-x)}2=c\eta(x)
\]
for some real number $c$.  Because $F$ is binary-valued,
$F(x)-F(-x)\in\{-1,0,1\}$.  If $c=0$, $F$ is constant on each antipodal
pair, so $S$ is a union of such pairs.  If $c\ne0$, the difference is nonzero
and has the same absolute value on all six pairs; hence $c=\pm\tfrac12$ and
each pair contributes exactly one element.  Solving for $F$ gives
$F=(1+\eta)/2$ or $F=(1-\eta)/2$, the two residue fibers modulo $3$.
They both have six elements, so they occur only when $p=3$.

There are $\binom6p$ unions of $p$ antipodal pairs.  Finally, if
$\Phi'=u\Phi$, then
$\widehat\varepsilon_{\Phi'}(\chi)=\chi(u)^{-1}
\widehat\varepsilon_\Phi(\chi)$, so its Fourier support is unchanged.  This
proves the translated statement as well.
\end{proof}

The same finite Galois calculation shows that the stabilizer of $\Phi$ in $G$ is trivial.  Hence this CM type is primitive: it is not pulled back from a CM type on any proper CM subfield of $M$.  In particular, the example is not produced by the obvious induction of a lower-dimensional CM type.

\begin{theorem}[Hodge ring of the certified CM sixfold]\label{thm:fullhodge}
For the CM type
\[
 \Phi=\{1,5,23,25,29,31\},
\]
the numbers of complex Hodge eigenlines in codimensions $p=1,\ldots,6$ are
\[
 6,\ 15,\ 22,\ 15,\ 6,\ 1.
\]
In every codimension except $p=3$, all of these lines are products of divisor lines.  In codimension $3$, exactly $20$ of the $22$ lines are products of three divisor lines; the two remaining lines are indexed by
\[
 S_+=\{1,13,19,25,31,37\},\qquad
 S_-=\{5,11,17,23,29,41\}.
\]
These are precisely the two complex eigenspaces of $\bigwedge_K^6 H^1(A,\Ql)$; their rational span is $W_K(A)$.  Consequently the full rational Hodge algebra of $A$ is generated by divisor classes together with $W_K(A)$.
\end{theorem}
\begin{proof}
Lemma~\ref{lem:primitivebinary} classifies every binary solution of
criterion~\eqref{eq:CMcriterion}, giving the displayed Hodge counts.  The six
degree-two solutions are precisely the antipodal pairs.  Their exterior
products are nonzero exactly when the chosen pairs are distinct, so they give
$\binom6p$, namely $6,15,20,15,6,1$, divisor-product eigenlines.  The only two
additional lines occur in the middle degree.  The two sets displayed above
are exactly the fibers $a\equiv1\pmod3$ and $a\equiv2\pmod3$, hence the two
eigenspaces obtained after complexifying $\bigwedge_K^6H^1$.  This identifies
their rational Galois-stable span with the two-dimensional Weil space.  The
registered exact subset enumerations independently verify the classification
and counts, but are not the sole explanation.
Although the eigenline description is made after complexification, the divisor eigenlines are
permuted by $G$ and their full span is the complexification of the rational divisor space;
likewise $S_+$ and $S_-$ form one Galois-stable pair.  Hence the resulting spanning
statement descends from complex eigenspaces to the rational Hodge algebra.
\end{proof}

\begin{corollary}[Hodge conjecture for the certified sixfold]\label{cor:HCspecific}
The Hodge conjecture holds in every codimension for the specific polarized CM abelian sixfold represented by the $k=11$ normalized $B$-vertex.
\end{corollary}
\begin{proof}
Divisor classes are algebraic by the Lefschetz $(1,1)$ theorem \cite[\S2, Remarks~(iii), p.~49]{DeligneHodge2006}, hence so are all their products.  Corollary~\ref{cor:Weilalg} gives algebraicity of the only two Hodge eigenlines not contained in the divisor algebra.  Theorem~\ref{thm:fullhodge} exhausts all Hodge classes in all degrees.
\end{proof}



\subsection{The A and C realizations and the fourth vertex}
\begin{theorem}[A second exact Hecke vertex]\label{thm:Avertex}
Let $A$ be McMullen's order-$14$ lift of the projective order-$7$ quadrilateral vertex generator in the index-two $(7,21,21,21)$ model, and put
\[
 A_t=h_tAh_t^{-1},\qquad h_t=\operatorname{diag}(t,1).
\]
For every one of the four $d=7$ arithmetic signature candidates
\[
 k\in\{2,14,23,35\}
\]
there are explicit $x_k,y_k\in L$ such that
\[
 T_k=x_kI+y_kA_t\in M_2(\OL),
 \qquad
 \operatorname{tr}(T_k)=c_k,
 \qquad
 \det(T_k)=\pi.
\]
Thus $T_k$ is an integral degree-$43$ Hecke representative and has exactly the same componentwise fixed points as $A_t$.  Consequently all four corrected $d=7$ trace classes admit a Hecke representative fixing the normalized $A$-vertex of McMullen's curve.

The CM type of the resulting $A$-vertex rational Hodge structure is
\[
 \Phi_A=\{1,5,31,29,17,19\}.
\]
Its restriction to $\Ql(\sqrt{-7})$ has multiplicities $(2,4)$, so the $A$-vertex is \emph{not} of $\Ql(\sqrt{-7})$-Weil type despite the arithmetic $d=7$ coefficient-sign filter.  Its restriction to $K=\Ql(\sqrt{-3})$, however, has multiplicities $(3,3)$.
\end{theorem}
\begin{proof}
McMullen's matrix $A$ has trace
\[
 \operatorname{tr}(A)=t^3-3t=\zeta_{42}^3+\zeta_{42}^{-3}.
\]
Hence $L[A_t]=M$.  Writing $\lambda_A=\zeta_{42}^3$, each $\alpha_k$ has a unique expansion
\[
 \alpha_k=x_k+y_k\lambda_A.
\]
For the two conjugate pairs these coefficients are
\[
\begin{aligned}
 x_2&=\frac{3t^5-t^4-19t^3+5t^2+29t-10}{7},\\
 y_2&=-\frac{4t^5+t^4-30t^3-5t^2+48t-4}{7},\\
 x_{14}&=\frac{3t^5-t^4-26t^3-2t^2+43t+4}{7},\\
 y_{14}&=\frac{(t-1)(3t^4+9t^3-3t^2-26t-18)}{7},
\end{aligned}
\]
and
\[
 (x_{23},y_{23})=-(x_2,y_2),\qquad
 (x_{35},y_{35})=-(x_{14},y_{14}).
\]
Although the displayed coefficients have denominator $7$ in the power basis, exact reduction in $L=\Ql[t]/(f)$ shows that all four matrices $T_k=x_kI+y_kA_t$ have entries in $\OL$ and satisfy the displayed trace and determinant identities.  Since $y_k\ne0$, the fixed-point equation of $T_k$ is a nonzero scalar multiple of that of $A_t$, so their fixed points agree at every real embedding.  Because $N_{L/\Ql}(\pi)=43$ is prime, an integral matrix of determinant $\pi$ has the required local Hecke elementary-divisor type.

For the CM type, the normalized $A_t$ fixed point lies in $\HH$ at every real place.  If $c_{A,t}$ denotes the lower-left entry of $A_t$, the eigenvalue on the $H^{1,0}$ fixed line has imaginary part $c_{A,t}\operatorname{Im}(z)$.  Exact interval evaluation at the six real embeddings, compared with the signs of $\sin(k\pi/7)$ for the eigenvalues $\zeta_{42}^{\pm3k}$, selects the extensions
\[
 1,5,-11,-13,17,19\pmod{42},
\]
which gives the displayed $\Phi_A$.  Reduction modulo $7$ gives two quadratic-residue and four nonresidue restrictions, while reduction modulo $3$ gives three of each embedding.  This proves the two signature assertions.
\end{proof}

\begin{theorem}[Hodge conjecture at the normalized $A$-vertex]\label{thm:Ahodge}
For the normalized $A$-vertex, every rational Hodge class is generated by divisor classes.  In particular the Hodge conjecture holds for the associated polarized abelian sixfold in every codimension.

Moreover, for the $K=\Ql(\sqrt{-3})$ Weil structure its Hermitian discriminant is the class of $-1$.
\end{theorem}
\begin{proof}
The element $\sqrt{-3}=1+2\zeta_{42}^{14}$ has a unique expression $x+y\zeta_{42}^3$ over $L$.  Replacing $\zeta_{42}^3$ by $A_t$ gives an integral Mobius representative $S_A$ with $S_A^2=-3I$. Its actual homology action is $\Sigma_A=S_A^{\mathsf T}$. For the trace polarization $E_1$, the corresponding Hermitian form
\[
 H_A(v,w)=E_1(\Sigma_Av,w)+\sqrt{-3}\,E_1(v,w)
\]
has Gram matrix $(-\operatorname{Tr}_{L/\Ql}((S_A)_{21}t^{i+j}))$ in the $K$-basis $(0,t^j)$, and exact determinant
\[
 \det H_A=-2^6 3^6 7^5\,127.
\]
Here $2^6=N(8)$, $3^6=N(27)$, $7=N(2+\sqrt{-3})$, and
\[
 127=N(10+3\sqrt{-3}),
\]
so the determinant class is $-1$.

It remains to determine all Hodge classes.  Applying the CM-character criterion to
\[
 \Phi_A=\{1,5,31,29,17,19\}
\]
gives respectively
\[
 12,\ 51,\ 88,\ 51,\ 12,\ 1
\]
complex Hodge eigenlines in codimensions $1,\ldots,6$.  There are exactly twelve degree-two Hodge eigenlines, indexed by
\[
\begin{gathered}
 \{1,13\},\ \{1,41\},\ \{5,23\},\ \{5,37\},\
 \{11,17\},\ \{11,31\},\\
 \{13,29\},\ \{17,25\},\ \{19,23\},\ \{19,37\},\
 \{25,31\},\ \{29,41\}.
\end{gathered}
\]
By the Lefschetz $(1,1)$ theorem their rational Galois-stable span consists of divisor classes.  In particular the two $K=\Ql(\sqrt{-3})$ Weil eigenlines are already divisor products; for example
\[
 \{a:a\equiv1\!\!\pmod3\}
 =\{1,13\}\sqcup\{19,37\}\sqcup\{25,31\},
\]
with the conjugate line decomposing as
\[
 \{5,23\}\sqcup\{11,17\}\sqcup\{29,41\}.
\]
Exhaustive exact enumeration of disjoint products of the twelve degree-two eigenlines gives
\[
 12,\ 51,\ 88,\ 51,\ 12,\ 1
\]
distinct eigenlines in the corresponding codimensions, exactly exhausting the preceding Hodge counts.  Thus the complexified Hodge algebra is generated in degree two, and Galois stability descends this statement to the rational Hodge algebra.  All such generators are algebraic divisors, proving the Hodge conjecture in every codimension.
\end{proof}

\begin{proposition}[The normalized $C$-vertex and saturation of the $43$-packet]\label{prop:Cvertex}
Let $C_t=h_tCh_t^{-1}$ be the normalized second order-$42$ generator.  Then
\[
 \operatorname{tr}(C_t)=t,\qquad \det(C_t)=1,\qquad \OL[C_t]=\OM.
\]
For every $k=0,\ldots,41$, substitution $\zeta_{42}\mapsto C_t$ gives
\[
 T_{C,k}=\alpha_k(C_t)\in M_2(\OL),\qquad
 \operatorname{tr}(T_{C,k})=c_k,\qquad
 \det(T_{C,k})=\pi.
\]
Thus every corrected $\ell=43$ trace class has a representative fixing the
normalized $C$-vertex.  The identical statement holds at $B_t$.

The CM type at $C$ is
\[
 \Phi_C=\{1,5,11,13,17,19\}.
\]
It has $K=\Ql(\sqrt{-3})$-signature $(3,3)$ and $K$-Hermitian discriminant
$-1$.  The Hodge eigenline counts are
\[
 6,\ 15,\ 22,\ 15,\ 6,\ 1,
\]
and the only two non-divisor lines are the two complex lines of $W_K(A_C)$.
Consequently the Hodge conjecture holds for the polarized $C$-vertex in every
codimension.
\end{proposition}
\begin{proof}
The relative discriminant of $X^2-tX+1$ is $t^2-4$, and exact norm
computation gives $N_{L/\Ql}(t^2-4)=1$.  Hence the order generated by either
$B_t$ or $C_t$ is the maximal order $\OM$.  Evaluating the integral element
$\alpha_k=\zeta_{42}^k(1+\zeta_{42}+\zeta_{42}^4)$ at either matrix therefore
gives an integral endomorphism.  Its $L$-trace and determinant are
$c_k$ and $N_{M/L}(\alpha_k)=\pi$.  Since $N_{L/\Ql}(\pi)=43$ is prime, the
matrix has the required local Hecke elementary-divisor type.

Exact interval evaluation shows that the lower-left entry of $C_t$ is positive
at all six real embeddings, so the upper fixed lines select the exponents
$1,5,11,13,17,19$.  Put $S_C=I+2C_t^{14}$.  Then
the actual homology action $\Sigma_C=S_C^{\mathsf T}$ satisfies $\Sigma_C^2=-3I$, $\Sigma_C^\dagger=-\Sigma_C$. In the $K$-basis $(0,t^j)$ the Hermitian Gram matrix is $(-\operatorname{Tr}_{L/\Ql}((S_C)_{21}t^{i+j}))$, and its determinant is
\[
 -2^6 3^6 7^5\,127.
\]
The factors outside the sign are norms because
$7=N(2+\sqrt{-3})$ and $127=N(10+3\sqrt{-3})$, leaving discriminant class
$-1$.  The exact CM-character enumeration then gives the displayed Hodge
counts and exactly two exceptional middle-degree lines, namely the two
$K$-Weil eigenspaces.  Markman's theorem applies.
\end{proof}

\begin{proposition}[The $B$- and $C$-vertices are distinct holomorphic classes]\label{prop:BCdistinct}
The normalized $B$- and $C$-vertices are not
$\mathrm{SL}_2(\OL)$-equivalent, although $B_t$ and $C_t$ are in the same
unoriented $\mathrm{GL}_2(\OL)$ similarity class.
\end{proposition}
\begin{proof}
The maximal order $\OM$ has class number one and $h(L)=1$, so the maximal-order
module classification gives a single $\mathrm{GL}_2(\OL)$ similarity class.
Let $P$ conjugate $B_t$ to $C_t$.  At the real embeddings $k=1,5$ the chosen
upper fixed lines carry the same eigenvalue $\zeta_{42}^k$, while at
$k=11,13,17,19$ the $B_t$ upper line carries $\zeta_{42}^{-k}$ and the
$C_t$ upper line carries $\zeta_{42}^k$.  Thus
\[
 \operatorname{sgn}(\det P)=(+,+,-,-,-,-).
\]
Every other integral conjugator differs by a centralizer unit in $\OM^\times$,
whose determinant is a relative unit norm and is totally positive.  Hence this
sign vector cannot be changed.  No integral conjugator is totally positive,
and in particular none lies in $\mathrm{SL}_2(\OL)$.
\end{proof}

\begin{proposition}[The fourth quadrilateral vertex is the $B$-class]\label{prop:DBsame}
Let $D$ be McMullen's fourth printed generator for the symmetric quadrilateral $Q(7,21,21,21)$ and put $D_t=h_tDh_t^{-1}$.  Then
\[
 \operatorname{tr}(D_t)=t,\qquad \det(D_t)=1,\qquad \OL[D_t]=\OM,
\]
and $D_t$ is $\mathrm{SL}_2(\OL)$-conjugate to $B_t$.  Consequently the four elliptic vertices $A,B,C,D$ of McMullen's quadrilateral determine exactly three holomorphic classes in the normalized Hilbert modular quotient: $A$, $B=D$, and $C$.
\end{proposition}
\begin{proof}
McMullen's explicit matrix for $D$ has trace $t$ and determinant one.  Since the relative discriminant of $X^2-tX+1$ is $t^2-4$ and $N_{L/\Ql}(t^2-4)=1$, both $\OL[B_t]$ and $\OL[D_t]$ equal the maximal order $\OM$.  The maximal-order integral similarity theorem therefore puts $B_t$ and $D_t$ in one $\mathrm{GL}_2(\OL)$ class.

Let $B=\left(\begin{smallmatrix}0&-1\\1&t\end{smallmatrix}\right)$ be the companion matrix for $X^2-tX+1$.  The rational matrix $P_0=(e_1,D_te_1)$ satisfies $P_0B=D_tP_0$, and exact rational-interval evaluation gives
\[
 \operatorname{sgn}(\det P_0)=(+,+,-,-,-,-).
\]
The same sign vector occurs for $h_t=\operatorname{diag}(t,1)$, which satisfies $h_tB=B_th_t$.  Therefore $P_0h_t^{-1}$ conjugates $B_t$ to $D_t$ and has totally positive determinant.  Any two conjugators between $B_t$ and $D_t$ differ through the centralizer $M^\times$, whose determinant is a relative norm and hence positive at every real place.  Thus an integral conjugator supplied by the maximal-order classification lies in the same positive orientation class; Corollary~\externalref{supplement}{cor:orientedcollapse} adjusts it through an integral centralizer unit to determinant one.  Hence $B_t$ and $D_t$ are $\mathrm{SL}_2(\OL)$-conjugate.
\end{proof}

\subsection{Comparison of the underlying isogeny classes}
\begin{proposition}[Comparison with known CM factors]\label{prop:knownCM}
Let $J_n$ be the Jacobian of the smooth projective model of
$y^2=x^n-1$, and let $r_n$ be the endomorphism induced by
$(x,y)\mapsto(\zeta_nx,y)$. Put
$X_{21}=\operatorname{im}(\Phi_3(r_{21})\Phi_7(r_{21}):J_{21}\to J_{21})$.
Forgetting the chosen polarizations and fixed RM markings, the three fibers
of Theorem~\ref{thm:A} satisfy
\[
 A_A\sim_{\mathbf C}J_7^2,\qquad
 A_B\sim_{\mathbf C}A_C\sim_{\mathbf C}X_{21}.
\]
Here $X_{21}$ is a simple abelian sixfold. Thus the three marked quotient
classes belong to exactly two underlying complex isogeny classes.
\end{proposition}
\begin{proof}
For odd $n$, the differentials $x^{j-1}dx/y$, $1\leq j\leq(n-1)/2$,
have $r_n$-characters $\zeta_n^j$. We use pushforward on divisor classes;
equivariance of the Abel--Jacobi map identifies these with the Jacobian
tangent eigenvalues. The polynomial $\Phi_3\Phi_7$ kills precisely the
nonprimitive characters of $J_{21}$, so $X_{21}$ has dimension six and
CM type $S=\{1,2,4,5,8,10\}$ for $\xi=\zeta_{21}$.
This is the known primitive factor in
\cite[Eq.~(4.2), Lem.~4.2, Prop.~4.1 and Sec.~5.3]{Goodson2024}.
Its type has trivial stabilizer in $(\mathbf Z/21\mathbf Z)^\times$,
hence is primitive and gives a simple variety
\cite[Prop.~3.13]{MilneCM2020}.

Since $z=\zeta_{42}=-\xi^{11}$, the action
$z\mapsto-r_{21}^{11}|_{X_{21}}$ has type
$\{21+22j\bmod42:j\in S\}=\Phi_B$.
Precomposing this action by $z\mapsto z^{19}$ gives
$19\Phi_B=\Phi_C$. This last automorphism acts nontrivially on $L$;
it does not identify the two prescribed RM markings.

For A, on $J_7^2$ define commuting endomorphisms
\[
 R=\operatorname{diag}(r_7,r_7),\quad
 U=\begin{pmatrix}0&-1\\1&-1\end{pmatrix},\quad
 T=-R^2U^2=\begin{pmatrix}r_7^2&-r_7^2\\r_7^2&0\end{pmatrix}.
\]
The identities $R^7=I$ and $U^2+U+I=0$ give
$T^6=R^5$, $T^{21}=-I$ and $T^{42}=I$.
Its tangent eigenvalues have $z$-exponents
$21+12j+28b\pmod{42}$ for $j=1,2,3$ and $b=1,2$.
Their set is $\{1,5,17,19,29,31\}=\Phi_A$.
All are primitive, so $\Phi_{42}(T)=0$ and $z\mapsto T$ gives
a faithful $M$-action. No compatibility with the product polarization is
asserted.

Each compared sixfold has rational homology of rank one over the
degree-twelve field $M$. Equal specified CM types identify their rational
complex structures; the corresponding full lattices in $M$ are
commensurable. Multiplication by a suitable positive integer therefore
gives an $M$-equivariant complex isogeny, which is algebraic. This is the
CM classification up to isogeny \cite[Prop.~3.12]{MilneCM2020}, and it
does not require the varieties to be simple. Finally, $J_7^2$ is not
simple, so its isogeny class differs from that of $X_{21}$.
\end{proof}


% !TeX root = ../../packet-and-principal-models.tex
\section{A common principal lattice and family}
The construction is a change of integral lattice, not a declaration that
the original trace polarization is principal. The subgroup at $7$ is
nonfree; requiring a common free line would incorrectly rule it out.
The special-fiber K-actions need not extend to a single global K-action.

\begin{proposition}[A common marked principal lattice at the three vertices]\label{prop:commonprincipal}
Use the common marking $\Lambda=\OL^2$, with $(a,b)\mapsto az_X+b$ at
$X=A,B,C$. There is one self-dual overlattice
$\Lambda\subset\Lambda^{\mathrm{common}}\subset\Lambda^\#$ of index
$453789$ stable under $\OL$ and all three specialized $\mathcal O_K$-actions
$\Sigma_X=S_X^{\mathsf T}$. With the respective complex structures it gives
three principally polarized isogenous special fibers with unchanged rational
CM Hodge structures.
\end{proposition}
\begin{proof}
Put $R=\OL/\mathfrak p_7^5=\mathbb F_7[u]/(u^5)$, where $u=t-1$.
The actual homology residues are
\[
 \Sigma_A\equiv\begin{pmatrix}1&3\\1&6\end{pmatrix},\qquad
 \Sigma_B\equiv\begin{pmatrix}6&2\\5&1\end{pmatrix},\qquad
 \Sigma_C\equiv\begin{pmatrix}5&5\\0&2\end{pmatrix}\pmod u.
\]
They all preserve the line of $v=(3,1)$. Define
\[
 H_7=u^2Rv+u^3R^2.
\]
Taking $w=(1,0)$, the vectors $v,w$ are an $R$-basis and
$H_7=u^2Rv\oplus u^3Rw$ has $\mathbb F_7$-basis
\[
 u^2v,\ u^3v,\ u^4v,\ u^3w,\ u^4w.
\]
Thus $|H_7|=7^5$.
For each $X$, write $\Sigma_Xv=\lambda_Xv+uq_X$; this immediately proves
$\Sigma_XH_7\subset H_7$. The determinant of any two elements of $H_7$
vanishes modulo $u^5$.

At $\mathfrak p_3$, all three $\Sigma_X$ vanish. Choose the same
first-coordinate $\mathbb F_{27}$-line $H_3$, and put $H=H_3\oplus H_7$.
To verify isotropy for the actual discriminant pairing, let $d=f'(t)$.
Multiplication by $d$ identifies $\Lambda^\#/\Lambda$ with
$(\OL/(d))^2$. Representatives $x,y\in\Lambda$ of elements of $H$ have
$\det(x,y)\in(d)$, by the two local determinant calculations. Hence
\[
 E_1(d^{-1}x,d^{-1}y)
 =-\operatorname{Tr}_{L/\Ql}\!\left(\frac{\det(x,y)}{d^2}\right)
 \in\mathbb Z,
\]
since the trace argument lies in the codifferent $d^{-1}\OL$.
As $|H|=3^3 7^5=\sqrt{|\Lambda^\#/\Lambda|}$, H is maximal isotropic.
Exact power-basis reduction verifies that all three
$U_X=(\Sigma_X-I)/2$ are integral. They preserve $\Lambda^\#$ as
$L$-linear integral matrices; since 2 is invertible on the discriminant
module, they preserve $H$. Thus its inverse image is simultaneously
$\OL$- and $\mathcal O_K$-stable and self-dual. The unchanged positive
Riemann form supplies the asserted principal polarizations.
\end{proof}

\begin{remark}[Free lines and the scope of the common construction]\label{rem:ABprincipallevels}
The common $H_7$ above is nonfree. Exhaustion of $\mathbb P^1(\mathbb F_7)$
shows that each pair of the three homology operators has only the common
residue line $(3,1)$. Its free lifts modulo $u^2$ have the form
$(3+au,1)$, with respective unique coefficients
\[
 a_A=3,\qquad a_B=6,\qquad a_C=2.
\]
Consequently no pair admits a common invariant \emph{free rank-one summand}
on $R^2$. This restricted obstruction does not exclude the common nonfree
maximal isotropic subgroup. Proposition~\ref{prop:commonprincipal} concerns
three specialized $K$-actions on one marked lattice. The global construction
below gives a family with RM and these three special fibers, without
asserting a single global $K$-action or a realization of Mostaed's literal
principally polarized locus.
\end{remark}

\begin{proposition}[Ambient and curve stabilizers of the common level]\label{prop:principalstabilizer}
Let $\Gamma=\mathrm{SL}_2(\OL)$, let $j:\Delta_0\to\Gamma$ be the
normalized integral representation of McMullen's quadrilateral group, and
let $H$ be the subgroup in Proposition~\ref{prop:commonprincipal}, with
$H_3=\mathbb F_{27}(1,0)$. For the transition representation
$\rho(g)=g^{-\mathsf T}$ put
\[
 \Gamma_H=\{g\in\Gamma:\rho(g)H=H\},
 \qquad \Delta_H=j^{-1}(\Gamma_H).
\]
Then
\[
 [\Gamma:\Gamma_H]=224,\qquad
 [\Delta_0:\Delta_H]=[\bar\Delta_0:\bar\Delta_H]=28.
\]
The exact ideal congruence level of $\Gamma_H$ is
$\mathfrak p_3\mathfrak p_7$: it contains $\Gamma(\mathfrak p_3\mathfrak p_7)$,
and $\Gamma(J)\subset\Gamma_H$ implies $J\subset\mathfrak p_3\mathfrak p_7$.
The number $28$ is the covering degree over the explicit parametrizing
orbifold $\mathbb H/\bar\Delta_0$, before an auxiliary neat refinement.
\end{proposition}
\begin{proof}
With the marking $(a,b)\mapsto az+b$, the isomorphism from the fiber at $z$
to the fiber at $gz$ is $w\mapsto w/(cz+d)$. Its homology matrix is
$g^{-\mathsf T}$. This transition action is a representation; it differs
from the fixed-point endomorphism matrix $g^{\mathsf T}$.

For any residue line $\ell\subset\mathbb F_7^2$, write
$H(\ell)=u^2Rv+u^3R^2$, where $v$ lifts $\ell$. An invertible matrix $T$
over $R$ preserves $H(\ell)$ precisely when its reduction preserves $\ell$.
Necessity follows by projecting to $u^2R^2/u^3R^2$. Conversely,
$Tv=\lambda v+uq$ gives $TH(\ell)\subset H(\ell)$, and finiteness gives
equality. Thus $\Gamma_H$ is the inverse image of two line stabilizers in
$\mathrm{SL}_2(\mathbb F_{27})\times\mathrm{SL}_2(\mathbb F_7)$.
Reduction is surjective: the finite groups are generated by elementary
matrices, whose entries lift by the Chinese remainder theorem. Hence the
ambient index is $(27+1)(7+1)=224$.

This description proves containment of $\Gamma(\mathfrak p_3\mathfrak p_7)$.
If $J$ is not contained in one of these primes, an elementary unipotent
with entry in $J$ can move the corresponding residue line. At $7$ use the
integral determinant-one basis with columns $(3,1),(-1,0)$. Applying
$\rho^{-1}$ gives an element of $\Gamma(J)$ outside $\Gamma_H$, proving
the asserted exactness of the level.

McMullen's generators $A,B,C,D$ generate $\Delta_0$ and satisfy $ABCD=I$
\cite[Sec.~4]{McMullen2026}, so $A,B,C$ suffice. Exact reduction of
$\rho(j(A)),\rho(j(B)),\rho(j(C))$ on the $224$-element set
$\mathbb P^1(\mathbb F_{27})\times\mathbb P^1(\mathbb F_7)$ fixes
$(3,1)$ in the second factor. Exhaustive closure of
$(\mathbb F_{27}(1,0),\mathbb F_7(3,1))$ is exactly
\[
 \mathbb P^1(\mathbb F_{27})\times\{\mathbb F_7(3,1)\}.
\]
The two separate implementations in Supplement, Section~\externalref{supplement}{app:computer} verify
the permutations and this entire finite orbit. Forward closure suffices
because inverses are positive powers of finite permutations. Orbit--stabilizer
therefore gives $28$, without a surjectivity assumption for the curve
group. Both $\Delta_0$ and $\Delta_H$ contain $-I$, so projectivization
preserves this index.
\end{proof}

\begin{theorem}[A connected principally polarized RM family]\label{thm:principalfamily}
The normalized modular embedding admits a connected finite cover $C'$
of its parametrizing orbifold and algebraic RM abelian schemes
$\mathcal A,\mathcal B$ of relative dimension six over $C'$, with an
$\OL$-linear isogeny
\[
 q:\mathcal A\longrightarrow\mathcal B,\qquad
 \deg(q)=453789,\qquad q^*\lambda_{\mathcal B}=\lambda_{\mathcal A}.
\]
Here $\lambda_{\mathcal A}$ is the original trace polarization and
$\lambda_{\mathcal B}$ is principal. The same connected cover contains
lifts of all three specified normalized vertices whose $\mathcal B$-fibers
are the common-lattice models of Proposition~\ref{prop:commonprincipal}.
Their specialized $\mathcal O_K$-actions and rational CM Hodge structures
are preserved, so the Hodge conjecture holds for these three fibers.

The moduli image lies in the principally polarized Hilbert--Blumenthal
component determined by the self-dual module $\Lambda_{\mathrm{common}}$.
No identification with Mostaed's named curve or global $K$-action is
included. The degree of a chosen auxiliary fine-level cover is not
asserted to be $28$.
\end{theorem}
\begin{proof}
Choose a neat finite-index subgroup $\Gamma_{\rm aux}\subset\Gamma_H$;
it exists by \cite[Thm.~3.6]{MilneModuli2011}, since $\Gamma_H$ is
arithmetic in $\operatorname{Res}_{L/\mathbb Q}\mathrm{SL}_2$.
Set $\Delta'=j^{-1}(\Gamma_{\rm aux})$ and
$C'=\mathbb H/\bar\Delta'$. The group $\bar\Delta'$ is torsion-free and
finite index in the cocompact quadrilateral group, so $C'$ is a connected
compact Riemann surface, hence a smooth projective algebraic curve.
Its map to the parametrizing orbifold is finite; the coarse map may
ramify at elliptic points. Since the neat group excludes $-I$, its
geometric degree is measured using projective group indices.

Over $\mathbb H$ use the fixed normalized embedding $F_t$ and the constant
lattices $\Lambda=\OL^2$ and $\Lambda'=\Lambda_{\mathrm{common}}$.
The transition action $\rho j$ descends both to integral local systems.
Their homology Hodge filtration $F^0$ is the kernel of the holomorphic map
\[
 (a,b)\longmapsto
 \bigl(\sigma_i(a)F_{t,i}(z)+\sigma_i(b)\bigr)_{i=1}^6.
\]
It has rank six and Hodge types $\{(-1,0),(0,-1)\}$; Griffiths
transversality is automatic for these types. The form $E_1$ is flat,
integral and positive by the Riemann-form identity already proved, and
is unimodular on $\Lambda'$. Thus both are polarizable integral
variations of Hodge structure. The discriminant local system
$\Lambda^\#/\Lambda$ has invariant subgroup $H$, and its inverse-image
lattice sheaf is exactly $\Lambda'$. This identifies the monodromy descent
of the common level.

Milne's equivalence for polarizable integral variations of these homology
types over smooth complex algebraic varieties
\cite[Thm.~7.13]{MilneModuli2011} algebraizes both families; the proof
explicitly includes nonprincipal polarizations. Full faithfulness
algebraizes the lattice inclusion, the $\OL$-actions and the
polarizations. The inclusion has index $453789$, and both forms are
$E_1$, giving the degree and pullback identity. Real multiplication is
self-adjoint and $\lambda_{\mathcal B}$ is principal.

Finally take the original three universal-cover parameters $z_A,z_B,z_C$.
Their images lie on the same connected curve $\mathbb H/\bar\Delta'$.
The constant marking $\Lambda'$ makes these fibers precisely the
three models in Proposition~\ref{prop:commonprincipal}. No choice of
components of a fiber product is involved. The specialized $K$-actions
need not be one global Hodge endomorphism. The RM family defines the
stated moduli map, and additional level gives a fine moduli interpretation.
\end{proof}
\begin{proposition}[The surviving order in the selected CM field]\label{prop:survivingorder}
For each of the three common principal fibers $B_X$, let
$\iota_X:M\hookrightarrow\operatorname{End}^0(B_X)$ be the selected CM
field embedding. Then
\[
 \operatorname{End}(B_X)\cap\iota_X(M)
 =\iota_X(R_K),\qquad
 R_K=\OL[\mathcal O_K]=\OL+\mathfrak p_3\OM.
\]
This order has index $27$ in $\OM$ and conductor $\mathfrak p_3\OM$.
The retained powers of the actual vertex automorphism are $0,7$ modulo
$14$ at $A$, and $0,7,14,21,28,35$ modulo $42$ at $B,C$.
The intersection assertion does not identify the full endomorphism
ring at the imprimitive $A$-fiber, or the generic endomorphism ring
of the family.
\end{proposition}
\begin{proof}
Put $z=\zeta_{42}$, so $\OM=\OL+\OL z$. Exact arithmetic gives
\[
 z^{14}=a+bz,\quad
 \begin{cases}
 a=-t^5+4t^3-3t-1,\\
 b=-t^5+t^4+6t^3-2t^2-8t-2.
 \end{cases}
\]
Both implementations verify $N_{L/\mathbb Q}(b)=-27$ and
$b\equiv0\pmod{\mathfrak p_3}$. Thus $(b)=\mathfrak p_3$, proving the
order formula and index. If $x=a_0+b_0z$ and both $x,xz$ lie in $R_K$,
their $z$ coefficients force $a_0,b_0\in\mathfrak p_3$; the reverse
inclusion is immediate. This proves the conductor assertion.

The intersection order $R'$ contains $R_K$ by the common-lattice theorem
and lies in $\OM$, since lattice endomorphisms are integral over $\mathbb Z$.
The quotient $\OM/R_K\simeq\OL/\mathfrak p_3$ is a simple $\OL$-module.
Consequently $R'$ is either $R_K$ or $\OM$. The latter would preserve the
full vertex automorphism, contradicted by the $\mathfrak p_3$ obstruction
proved above for all three vertices. Hence $R'=R_K$. Exhausting the
$14$ or $42$ powers on the two finite local subgroups gives the stated
surviving powers, with separate exact checks in both implementations.
\end{proof}

\begin{remark}[Imprimitive CM type at the $A$-vertex]
The stabilizer of $\Phi_A$ in $G=(\mathbb Z/42\mathbb Z)^\times$ is $\{1,29\}$.  Its fixed field is $\Ql(\zeta_7)$, and $\Phi_A$ is induced from the three-element CM type on that sextic CM field.  Thus the larger divisor algebra at the $A$-vertex is consistent with the fact that this CM type is imprimitive; in particular the associated sixfold is isogenous to the square of a CM threefold.
\end{remark}


% !TeX root = ../../packet-and-principal-models.tex
\section{Integral models, polarizations and maps}
We first give coordinates for the common principal component. We then
separate marked, polarized and unpolarized equivalence. The ideal and unit
calculations are also the explicitly cited arithmetic inputs to Paper II.


\begin{proposition}[An explicit principal RM component]
\label{pp:prop:explicit}
Write $O=\OL$, $u=t-1$, $d=f'(t)$, and
\[
 \beta=-t^5+t^4+6t^3-2t^2-8t-2,\qquad
 a=(\beta u^2)^{-1},\quad b=u^{-3},\quad
 \delta=ab=(\beta u^5)^{-1}.
\]
The common self-dual lattice $P=\Lambda_{\mathrm{common}}$ has $O$-basis
\begin{equation}\label{pp:eq:basis}
 T=\begin{pmatrix}a&3b\\0&b\end{pmatrix},\qquad P=TO^2.
\end{equation}
Its determinant $\delta$ generates the codifferent and has signs
$(+,-,+,-,+,-)$ in the ordered embeddings $1,5,11,13,17,19$.
A principal RM presentation of the family in
Theorem~\ref{thm:principalfamily} has homology module
\[
 P_{\rm std}=Oe_1+\delta Oe_2
\]
with the trace form $E_1$, period lattice $Ow+\delta O$, and arithmetic group
\begin{equation}\label{pp:eq:group}
 \Gamma_{\rm std}
 =\left\{\begin{pmatrix}r&\delta s\\v/\delta&q\end{pmatrix}:
 r,s,v,q\in O,\ rq-sv=1\right\}
 =\operatorname{diag}(\delta,1)\SL_2(O)
   \operatorname{diag}(\delta,1)^{-1}.
\end{equation}
Its period map on the universal covering curve is
\begin{equation}\label{pp:eq:curve}
 F_{{\rm std},i}(z)
 =\frac{\sigma_i(a^2)F_{t,i}(z)}{3F_{t,i}(z)+1}.
\end{equation}
Thus its image lies in the specified principal RM component
$\HH^6/\Gamma_{\rm std}$, with a compatible neat refinement when a fine
moduli interpretation is required.
\end{proposition}
\begin{proof}
The ideal identities $(u)=\mathfrak p_7$, $(\beta)=\mathfrak p_3$ and
$(d)=\mathfrak p_3\mathfrak p_7^5$ give the unit
\[
 \varepsilon=\frac{d}{\beta u^5}
 =28+215t+335t^2+82t^3-87t^4-35t^5,
 \qquad N_{L/\Ql}(\varepsilon)=1.
\]
They may also be checked directly from $f$: $N(u)=7$,
$N(\beta)=-27$ and $\beta$ vanishes modulo the degree-three factor at $3$.
Under multiplication by $d$, the local conditions defining the common
subgroup in Proposition~\ref{prop:commonprincipal} say
\[
 y\in\beta O,\qquad y\in u^2O,\qquad x-3y\in u^3O.
\]
Since $(\beta)$ and $(u)$ are coprime, their simultaneous solution is
\[
 dP=Ou^3(1,0)+O\beta u^2(3,1).
\]
Now
\[
 dT=\varepsilon
 \begin{pmatrix}u^3&3\beta u^2\\0&\beta u^2\end{pmatrix},
 \qquad
 T^{-1}=\begin{pmatrix}\beta u^2&-3\beta u^2\\0&u^3\end{pmatrix}.
\]
The unit factor leaves the module unchanged, proving
\eqref{pp:eq:basis} and $O^2\subset P$. Also $\delta d=\varepsilon$,
so $\delta O=\mathfrak D_L^{-1}$ and
$[P:O^2]=|N(\delta)|^{-1}=453789$.

In the basis $T$, the pairing is
$E_\delta((r,s),(r',s'))=\operatorname{Tr}_{L/\Ql}
(\delta(r's-rs'))$. Exact arithmetic gives
\[
 \delta=-\frac{11}{21}-\frac{209}{21}t-\frac{55}{3}t^2
       -\frac{110}{21}t^3+\frac{33}{7}t^4+2t^5.
\]
The integral matrix $(\operatorname{Tr}(\delta t^{i+j}))_{0\leq i,j\leq5}$
has determinant $-1$, so its alternating trace form is unimodular.
Disjoint rational isolating intervals for all six roots of $f$, followed
by rational interval evaluation of this polynomial, certify the stated
signs. The exact identities and signs are checked in
\path{experiments/principal_normalization_v111/explicit_basis.py}
and separately, using rational polynomial arithmetic, in
\path{experiments/principal_normalization_v111/review/independent_basis_check.py}.

At a period $z\in\HH^6$, interpreted componentwise, the basis columns
evaluate to $az$ and $b(3z+1)$. Division by $b(3z+1)$ gives the period
\[
 \zeta=\frac{a}{b}\frac{z}{3z+1}
\]
on the standard $O^2$ module with form $E_\delta$. Its imaginary-part
signs are those of $\delta$; $E_\delta$ is positive on this mixed
component. Multiplication of the complex coordinate by $\delta$ gives
$w=\delta\zeta=a^2z/(3z+1)$ and lattice $Ow+\delta O$. The corresponding
homology identification is
\[
 C=\operatorname{diag}(1,\delta)T^{-1},\qquad \det C=1.
\]
It maps $P$ onto $P_{\rm std}$ and preserves $E_1$ exactly. Thus this
is a polarized isomorphism, not just an isogeny comparison.
The principal module $P_{\rm std}$ is self-dual because $\delta O$ is
the codifferent. Its transition stabilizer for $g\mapsto g^{-\mathsf T}$
is precisely \eqref{pp:eq:group}. The form $E_1$ is positive on its
all-upper-half-plane domain.

The composite period matrix is
\[
 W=\begin{pmatrix}a^2&0\\3&1\end{pmatrix}\in\GL_2^+(L),
 \qquad \det W=a^2.
\]
It gives \eqref{pp:eq:curve}; its denominators do not vanish on $\HH$.
The stabilizer calculation below verifies equivariance on the same
connected cover as Theorem~\ref{thm:principalfamily}. That theorem's
algebraization identifies the resulting moduli map with the original
principal family, preserving its degree-$453789$ isogeny and its three
specialized $K$-actions.
\end{proof}

\begin{proposition}[Explicit transition stabilizer]\label{pp:prop:stabilizer}
For $g=\left(\begin{smallmatrix}r&s\\v&q\end{smallmatrix}\right)
\in\SL_2(O)$ put $Q(g)=v+3r-3q-9s$. Then
\[
 \Gamma_H=\{g\in\SL_2(O):s\in\beta O,\ Q(g)\in uO\},
 \qquad
 W\Gamma_HW^{-1}=\Gamma_{\rm std}\cap W\SL_2(O)W^{-1}.
\]
\end{proposition}
\begin{proof}
The homology basis matrix is $T$, the period matrix is $S=T^{\mathsf T}$,
and the transition action is $g^{-\mathsf T}$. Direct multiplication gives
\[
 T^{-1}g^{-\mathsf T}T
 =\begin{pmatrix}
 q+3s&-(\beta/u)Q(g)\\-(u/\beta)s&r-3s
 \end{pmatrix}.
\]
Its integrality is equivalent to the two displayed congruences, because
$u$ and $\beta$ are coprime. Determinant one then gives a lattice
automorphism. On the final periods,
\[
 WgW^{-1}
 =\begin{pmatrix}r-3s&a^2s\\a^{-2}Q(g)&q+3s\end{pmatrix}.
\]
The identities $a^2/\delta=u/\beta$ and $\delta a^{-2}=\beta/u$ prove
the intersection assertion. Consequently \eqref{pp:eq:curve} is
equivariant for $Wj(\Delta')W^{-1}\subset\Gamma_{\rm std}$.
The indices $224$ and $28$ are those of
Proposition~\ref{prop:principalstabilizer}; no new degree after neat
refinement is inferred.
\end{proof}

\subsection{Marked and unmarked distinctions}
\begin{proposition}[Distinct principal points with fixed real multiplication]
\label{pp:prop:markeddistinct}
Let $B_X$, for $X=A,B,C$, be the three specified principal fibers of
Theorem~\ref{thm:principalfamily}, with their fixed $\OL$ actions.
For $X\ne Y$ one has
\[
 \operatorname{Hom}_{\OL}(B_X,B_Y)\otimes_{\mathbf Z}\Ql=0.
\]
Consequently their images in the principal RM component
$\HH^6/\Gamma_{\rm std}$ are pairwise distinct, also after forgetting
auxiliary level while retaining the $\OL$ marking.
\end{proposition}
\begin{proof}
The actual geometric CM types, expressed using $z=\zeta_{42}$, are
\[
 \Phi_A=\{1,5,17,19,29,31\},\quad
 \Phi_B=\{1,5,23,25,29,31\},\quad
 \Phi_C=\{1,5,11,13,17,19\}.
\]
Their $M$ actions restrict to the same prescribed $L$ action:
the generators for $B,C$ satisfy $Z^2-tZ+I=0$, and the certified
generator $Z_A=(A_t+tI)/(t^2-1)$ of
Proposition~\externalref{supplement}{prop:Amaximalpair} satisfies the same identity and
$Z_A^3=A_t$. Passing to transposes gives the homology actions without
changing the scalar $t$. Thus $z+z^{-1}$ acts as the fixed scalar $t$
in every case. These are the normalized geometric types, not companion types
or representatives re-marked by arbitrary Galois automorphisms.
The lattice inclusions of Theorem~\ref{thm:principalfamily} are the
identity on rational $L^2$, and the coordinate normalization in
Proposition~\ref{pp:prop:explicit} commutes with $L$ at each real
place. Hence the same marked rational Hodge structures, with their
rational $M$ actions, describe the principal fibers.

For completeness, consider two rank-one $M$ rational Hodge structures
with this fixed $L$ restriction. After identifying each rational
vector space with $M$, every rational linear map is uniquely
\[
 F=\sum_{g\in G}m_g g,\qquad m_g\in M,\qquad
 G=\operatorname{Gal}(M/\Ql).
\]
Indeed, after splitting $M$ over $\mathbf C$, each summand $Mg$
occupies the matrix positions of its embedding permutation. These
supports are disjoint and their dimensions sum to $[M:\Ql]^2$.
Commutation with every $a\in L$ gives
$m_g(g(a)-a)=0$. Since $M$ is a field, only
$g\in\operatorname{Gal}(M/L)=\{1,c\}$ can remain, where $c$ is
complex conjugation. If $m_g\ne0$, every conjugate of $m_g$ is
nonzero. Its entries cannot cancel those of another summand.
Preserving the Hodge subspaces therefore requires the two types to
be equal for $g=1$, or complementary for $g=c$.

The three displayed types are pairwise unequal, and all contain $1$
and omit $41$, so no two are complementary. Thus every $L$-linear
Hodge map between distinct fibers is zero. An $\OL$-linear
homomorphism of abelian varieties induces such a map on rational
homology, and a homomorphism with zero map on homology is zero.
This proves the assertion without a primitivity assumption.
An identification in the specified principal RM moduli component
would be an $\OL$-linear polarized isomorphism, which is excluded.
The unmarked isogeny relation $19\Phi_B=\Phi_C$ from
Proposition~\ref{prop:knownCM} changes the $L$ marking and gives no
contradiction. The next proposition treats the unmarked principal
images; neither result identifies Mostaed's named image.
\end{proof}


\begin{proposition}[Three distinct unmarked principal fibers]
\label{pp:prop:unmarkeddistinct}
The three specified principally polarized fibers $B_A,B_B,B_C$ of
Theorem~\ref{thm:principalfamily} are pairwise nonisomorphic as principally
polarized complex abelian varieties, after forgetting both the fixed
$\OL$ action and auxiliary level. Thus they give three distinct points
of the coarse moduli space of principally polarized abelian sixfolds.
\end{proposition}
\begin{proof}
By Proposition~\ref{prop:knownCM}, the underlying $A$ fiber is isogenous
to $J(y^2=x^7-1)^2$, whereas the $B,C$ fibers are isogenous to the simple
sixfold $X_{21}$. The principal-model isogenies preserve this distinction,
so only $B_B$ versus $B_C$ remains.

Write $z=\zeta_{42}$, $t=z+z^{-1}$, $u=t-1$ and
\[
 \beta=-t^5+t^4+6t^3-2t^2-8t-2,\qquad
 a=(\beta u^2)^{-1},\quad b=u^{-3}.
\]
The actual common lattice of Proposition~\ref{pp:prop:explicit} is
$P=\OL(a,0)+\OL(3b,b)$, with the restriction of $E_1$.
For $X=B,C$, let $(c_X,d_X)$ be the bottom row of the normalized
vertex matrix $X_t$. Explicitly,
\[
 c_B=t^{-1},\quad d_B=t,\qquad
 c_C=t^{-1}(-1+5t+4t^2-t^3-t^4),\quad
 d_C=-t+4t^2+4t^3-t^4-t^5.
\]
On homology, $z$ acts by $X_t^T$. The identification
$J_X:M\longrightarrow L^2$ given by
$J_X(v_0+v_1z)=(c_Xv_1,v_0+d_Xv_1)$ therefore has
\[
 I_X=J_X^{-1}(P)=\OL a w_X+\OL b(1+3w_X),\qquad
 w_X=(z-d_X)/c_X.
\]
In these coordinates its actual polarization is
\[
 E_X(v,w)=\Tr_{M/\Ql}(\xi_Xv\bar w),\qquad
 \xi_X=-\frac{c_X}{z-z^{-1}}.
\]
Indeed the relative trace equals $c_X(v_0w_1-v_1w_0)$, which
recovers $E_1(J_Xv,J_Xw)$ after applying $\Tr_{L/\Ql}$.

The disjoint-support argument in
Proposition~\ref{pp:prop:markeddistinct} applies without imposing
$L$-linearity: a term $m_g\sigma_g$ can preserve the Hodge subspaces
only if $\Phi_B=g\Phi_C$. For the actual types displayed there, the
unique such exponent is $g=31$ modulo $42$. Hence every possible
isomorphism has the form $F(v)=m\sigma_{31}(v)$. Its integrality and
polarization conditions imply, respectively,
\[
 m\in J:=\frac{I_C}{\sigma_{31}(aw_B)},\qquad
 m\bar m=r:=\frac{\sigma_{31}(\xi_B)}{\xi_C}.
\]
The first condition uses only the image of the first actual lattice
generator; it requires no invertibility of a nonmaximal-order ideal.
The second follows by moving $\sigma_{31}$ through the absolute trace.
In particular $q(m):=\Tr_{M/\Ql}(m\bar m/r)=12$ is necessary.

We record the finite obstruction explicitly. Let $Q$ be the rational
Gram matrix of $q$ on the basis of $J$ induced by
$t^ja w_C,t^jb(1+3w_C)$, $0\leq j<6$, divided by
$\sigma_{31}(aw_B)$. The certificate
\path{certificates/principal_unmarked_bc_v122.py} reconstructs this
matrix from the displayed field and polarization formulas. Its data
file \path{certificates/principal_unmarked_bc_data.json} supplies an
integral matrix $U$ with determinant $-1$ and the exact identity
$U^TQU=H$, where
\begingroup
\setlength{\arraycolsep}{2.3pt}
\scriptsize
\[
H=\left(\begin{array}{*{12}{r}}
14&-7&0&0&-6&-2&3&2&0&0&-3&6\\
-7&14&0&0&3&4&-6&-1&0&0&-3&-3\\
0&0&14&0&4&-1&-2&-6&0&-7&2&-4\\
0&0&0&14&-1&-6&2&-2&-7&0&1&1\\
-6&3&4&-1&16&4&-8&-8&2&-2&-1&2\\
-2&4&-1&-6&4&16&-8&4&3&2&2&2\\
3&-6&-2&2&-8&-8&16&4&-1&4&-1&-1\\
2&-1&-6&-2&-8&4&4&16&4&3&-1&2\\
0&0&0&-7&2&3&-1&4&14&0&1&-2\\
0&0&-7&0&-2&2&4&3&0&14&2&2\\
-3&-3&2&1&-1&2&-1&-1&1&2&22&-11\\
6&-3&-4&1&2&2&-1&2&-2&2&-11&22
\end{array}\right).
\]
\endgroup
Exact rational factorization proves $H$ positive definite, and direct
inversion gives $12(H^{-1})_{ii}<4$ for every $i$. Cauchy--Schwarz
therefore places every integer vector of norm at most $12$ in
$\{-1,0,1\}^{12}$. Exact enumeration of this box gives only the
origin in that closed ball. The primary certificate instead uses
rational completion of squares, visiting $59$ nodes with the same
result. The separate implementation
\path{independent/principal_unmarked_bc_v122.py} reconstructs $Q$
by cyclotomic companion matrices and checks all $3^{12}=531441$ box
vectors by an integer split quadratic form. Both paths verify the
unimodular change of basis and the complete bounds.
In fact the lattice minimum is $14$: $H$ is even integral, the
closed ball of norm $12$ contains only zero, and its first basis vector
has norm $14$.

Thus the necessary trace norm $12$ is impossible, excluding every
polarized isomorphism $B_B\simeq B_C$. This leaves their established
unmarked isogeny relation unchanged. The stronger unpolarized distinction
is proved in Proposition~\ref{pp:prop:unpolarizeddistinct} below; neither
argument asserts injectivity of the whole principal curve or identifies
it with Mostaed's named image.
\end{proof}


\subsection{Scalar-polarized maps}
\begin{proposition}[The smallest scalar-polarized $B$--$C$ isogeny]
\label{pp:prop:bcisogeny}
For the specified principal fibers $(B_B,\lambda_B)$ and
$(B_C,\lambda_C)$, the smallest positive integer $k$ for which an
isogeny $f:B_B\to B_C$ satisfies $f^*\lambda_C=k\lambda_B$ is $2$.
There are exactly six such isogenies at $k=2$, each of degree $64$.
They have the same kernel, a maximal isotropic subgroup of $B_B[2]$
isomorphic to $(\Z/2\Z)^6$, and form one orbit under the polarized
automorphisms of $B_C$. The origin-preserving polarized automorphism
groups of both fibers are the explicit cyclic groups of order six
generated, in their respective CM coordinates, by multiplication by $z^7$.
\end{proposition}
\begin{proof}
Use $I_B,I_C,\xi_B,\xi_C,r,J,Q,H$ and the unique Hodge twist
$\sigma_{31}$ from Proposition~\ref{pp:prop:unmarkeddistinct}.
A rational Hodge map $m\sigma_{31}$ is an isogeny with multiplier $k$
if and only if
\[
 m\sigma_{31}(I_B)\subseteq I_C,\qquad m\bar m=kr.
\]
The inclusion must be tested on the full integral lattices; membership
in $J$ alone is insufficient. The norm equation is nonzero, so it makes
the rational map invertible. An integral Hodge map between these
abelian varieties is algebraic, as in the homology comparison above.
On the common integral symplectic form $E$, its matrix satisfies
$T^tET=kE$, so Pfaffians give its degree $\det T=k^6$.
Moreover, any rational scalar multiplier is integral: unimodularity
makes $E$ primitive, so integer linear combinations of its entries
equal $1$, and the corresponding combinations of the entries of
$T^tET$ equal $k$. Positivity then gives $k>0$.

For $k=2$ the necessary trace norm is $24$. The exact inverse-diagonal
bounds for $H$ give
\[
 (|n_1|,\ldots,|n_{12}|)\leq(1,1,1,1,2,1,1,1,1,1,1,1).
\]
The certificate \path{certificates/principal_bc_isogeny_v123.py}
and the separate companion-matrix implementation
\path{independent/principal_bc_isogeny_v123.py} prove completeness
by rational suffix recursion and Schur-complement prefix minima,
respectively. The independent path records $175$ vectors in the closed
trace-norm ball of value $24$, and both find $84$ on its boundary.
Of the latter, $72$ fail the
full relative norm equation, and six more fail lattice integrality.
The remaining six satisfy the full norm equation, all $144$ matrix
integrality conditions, $T^tET=2E$ and $\det T=64$.

Explicitly, one map is $F_0(v)=m_0\sigma_{31}(v)$, where
\[
\begin{split}
127m_0={}&19-98z-85z^2-37z^3+15z^4+48z^5-56z^6\\
        &{}+24z^7+7z^8+105z^9+82z^{10}+2z^{11}.
\end{split}
\]
Its integral matrix is retained in
\path{certificates/principal_bc_isogeny_data.json}; the implementations
reconstruct it directly from the actual bases of $I_B,I_C$ and verify
these identities. The complete six-map list is
$F_j=z^{7j}F_0$, $0\leq j<6$. Multiplication by $z^7$ preserves each
actual integral lattice and its polarization and has exact order six;
both implementations check these statements directly.
Thus all six maps have the same kernel.

The identity $T^tET=2E$ gives the integral matrix
$2T^{-1}=E^{-1}T^tE$. Hence the cokernel of $T$ is killed by $2$;
its order $64$ makes it $(\Z/2\Z)^6$. It identifies with the kernel
of the isogeny in $B_B[2]$. For representatives $v,w$ of kernel points,
$2E_B(v,w)=E_C(F_0v,F_0w)$ is integral, so their $2$-Weil pairing is
trivial. Its order $2^6$ makes this isotropic subgroup maximal.
The independent certificate also computes the Smith form
$\operatorname{diag}(1,1,1,1,1,1,2,2,2,2,2,2)$ for every map.

Postcomposition with any polarized automorphism of $B_C$ permutes the
six maps faithfully, since $F_0$ is surjective. Its orbit of $F_0$
therefore has the size of that automorphism group and at most six
elements. The six displayed units already occur, proving that they
are the full group. The same argument by precomposition applies to
$B_B$: if $F_0a=F_0b$, the connected image of the homomorphism $a-b$
lies in the finite kernel of $F_0$, so $a=b$.

Finally multiplier one would have degree one and give a polarized
isomorphism, excluded by Proposition~\ref{pp:prop:unmarkeddistinct}.
Thus multiplier two is minimal. This does not assert a minimum degree
among arbitrary unpolarized isogenies, nor a fixed-$\OL$-linear map;
the twist $\sigma_{31}$ changes the prescribed real multiplication.
\end{proof}

\subsection{Unpolarized isomorphisms and minimum degree}
\begin{proposition}[Three distinct underlying abelian varieties]
\label{pp:prop:unpolarizeddistinct}
The three specified principal models $B_A,B_B,B_C$ of
Theorem~\ref{thm:principalfamily} are pairwise nonisomorphic as complex
abelian varieties, even after forgetting their polarizations and real
multiplication markings.
\end{proposition}
\begin{proof}
The accepted simple-factor decomposition distinguishes $B_A$ from
$B_B,B_C$. For the latter pair, use the actual CM-coordinate lattices
$I_B,I_C$ and the rational Hodge twist $\sigma_{31}$ from
Proposition~\ref{pp:prop:unmarkeddistinct}. Write
$S_X=\OM I_X$ for their maximal-order saturations, and let
$F_0=m_0\sigma_{31}$ be the explicit degree-$64$ isogeny of
Proposition~\ref{pp:prop:bcisogeny}. Put
\[
 s=z^6+z^{12}+z^{24}-z^{18}-z^{30}-z^{36},\qquad
 \alpha=\frac{1+s}{2},\qquad \varpi=1-z^{14}.
\]
Then $s^2=-7$, $\alpha\bar\alpha=2$, and exact lattice calculations give
\begin{equation}\label{pp:eq:unpolarizedlocaldata}
 F_0(S_B)=\alpha S_C,\qquad [S_X:I_X]=27,\qquad
 [\OM:R]=27,\quad R=(I_C:I_C),\quad
 v-\bar v\in\varpi\OM\quad(v\in R).
\end{equation}
The first equality includes an integral determinant-one basis change;
an equality of indices alone would not suffice. The last condition is
verified on a basis of the \emph{full} multiplier order, not merely a
selected suborder.

Here is a finite exact description of the computation. In the basis
$1,z,\ldots,z^{11}$, the saturation is the lattice generated by
$z^i I_X$, $0\leq i<12$. For each basis element of $M$, write its
multiplication matrix on $I_C$ and collect the $144$ entries as a column
of a rational matrix $A$. The full multiplier order consists exactly
of the coefficient vectors $x$ for which $Ax\in\Z^{144}$; it is the
dual of the lattice generated by the rows of $A$.
The certificate \path{certificates/principal_bc_unpolarized_v125.py}
reconstructs these lattices by integer Hermite forms after clearing
denominators. The separate certificate
\path{independent/principal_bc_unpolarized_v125.py} checks the
saturations by containment, stability under $z$ and generation modulo
$3$, and the full multiplier order by linear integrality constraints
modulo $3$. Both verify \eqref{pp:eq:unpolarizedlocaldata} on the actual
C56 lattices and the displayed $F_0$.

The ideal $\mathfrak P=\varpi\OM$ is the prime above $3$ with residue
field $\mathbb F_{729}$. Indeed $\zeta_3=z^{14}$ reduces to $1$, and
$3$ has order $6$ modulo $7$, so $\Phi_7$ is irreducible over
$\mathbb F_3$. Complex conjugation acts on this field by $x\mapsto
x^{27}$, since $27\equiv-1\pmod7$.

Suppose now that $F=m\sigma_{31}$ were an unpolarized isomorphism
$B_B\to B_C$. It would send $I_B$ onto $I_C$ and hence $S_B$ onto
$S_C$. Comparing this with \eqref{pp:eq:unpolarizedlocaldata} shows that
\[
 \eta S_C=S_C,\qquad \eta=\alpha m/m_0.
\]
Multiplication by $\eta$ and by $\eta^{-1}$ therefore has an integral
matrix on the full lattice $S_C$. Their characteristic polynomials
show that both are algebraic integers: $\eta\in\OM^\times$.
This uses no class-number-one or principal-ideal choice.

Let $\Z_{(3)}$ denote localization at the rational prime $3$.
Since $\det F_0=64$, the maps $F$ and $F_0$ are both isomorphisms
from $I_B\otimes\Z_{(3)}$ to $I_C\otimes\Z_{(3)}$.
Thus multiplication by $u=m/m_0$ belongs to
$(R\otimes\Z_{(3)})^\times$. By the last condition in
\eqref{pp:eq:unpolarizedlocaldata}, its nonzero residue is fixed by
complex conjugation. Consequently
\[
 \frac{\eta}{\bar\eta}\equiv
 \frac{\alpha}{\bar\alpha}\equiv s\pmod{\mathfrak P}.
\]
The second congruence uses $s^2=-1$ in characteristic $3$:
$(1+s)/(1-s)=s$. The residue of $s$ has exact order $4$.
On the other hand, the CM unit $\eta/\bar\eta$ has absolute value $1$
at every complex embedding, so Kronecker's theorem makes it a root
of unity in $M$, hence an element of $\mu_{42}$. In characteristic
$3$ its reduction has order dividing $14$. The two orders are
incompatible, proving the contradiction.
The degree-$64$ isogenies remain valid. The next proposition determines
the least degree among all unpolarized isogenies.
\end{proof}

\begin{proposition}[The least unpolarized isogeny degree]
\label{pp:prop:mindegree}
Among all isogenies $B_B\to B_C$ of the specified principal models,
with no polarization compatibility imposed, the least degree is $7$.
There is an explicit such isogeny with cyclic kernel of order $7$.
\end{proposition}
\begin{proof}
Retain $I_X,S_X,m_0,\alpha$ from the preceding propositions and put
$\beta=1+z-z^3$. Its absolute norm is $7$. Define
\[
 m_7=z^5\beta\alpha^{-1}m_0,\qquad F_7=m_7\sigma_{31}.
\]
Explicitly,
\[
\begin{split}
127m_7={}&-40-41z-55z^2-9z^3+62z^4+46z^5+31z^6\\
       &{}+23z^7+12z^8+53z^9+68z^{10}-51z^{11}.
\end{split}
\]
The matrix $I_C^{-1}M(m_7)\sigma_{31}(I_B)$ is integral with
determinant $7$. Thus the rational Hodge map $F_7$ induces an algebraic
isogeny of degree $7$, with cyclic kernel. Its exact integral matrix is
retained in \path{certificates/principal_bc_min_degree_data.json}.
The certificate \path{certificates/principal_bc_min_degree_v126.py}
checks every entry and the determinant; the separate companion
implementation \path{independent/principal_bc_min_degree_v126.py}
reconstructs the map independently and checks the Smith form.

For the lower bound, any isogeny $F=m\sigma_{31}$ sends $I_B$ into
$I_C$, and hence sends $S_B$ into $S_C$. The exact identity
$F_0(S_B)=\alpha S_C$ gives
\[
 F(S_B)=\eta S_C\subseteq S_C,\qquad \eta=\alpha m/m_0.
\]
Its integral multiplication matrix makes $\eta$ an algebraic integer.
Since $[S_B:I_B]=[S_C:I_C]=27$, lattice indices yield
\[
 \deg F=[I_C:F(I_B)]=[S_C:F(S_B)]
       =[S_C:\eta S_C]=N_{M/\mathbb Q}(\eta).
\]
The last norm is positive because $M$ is CM. If it were between $2$
and $6$, the ideal $(\eta)$ would have a prime ideal factor of norm
$2$, $3$, $4$ or $5$. None exists. Indeed $\OM=\Z[z]$, and exact
polynomial arithmetic gives
\[
 \gcd(\Phi_{42},X^4-X)=1\quad\text{over }\mathbb F_2,
 \qquad
 \gcd(\Phi_{42},X^p-X)=1\quad\text{over }\mathbb F_p,
 \quad p=3,5.
\]
These exclude residue degrees one or two above $2$ and degree one
above $3$ or $5$, exhausting the possible prime-power norms at most
$6$. The primary certificate checks the three gcds; the independent
certificate evaluates $\Phi_{42}$ on $\mathbb F_4$, $\mathbb F_3$
and $\mathbb F_5$ and verifies that there are no roots.
Degree one is excluded by Proposition~\ref{pp:prop:unpolarizeddistinct}.
Together with $F_7$, this proves the asserted minimum.

The map $F_7$ is not scalar-polarized for the specified principal
polarizations: a positive rational scalar multiplier would be an
integer $k$ and give degree $k^6$, which cannot equal $7$.
This is compatible with the minimum scalar multiplier $2$ and degree
$64$ in Proposition~\ref{pp:prop:bcisogeny}. We assert no finite census
of all unpolarized isogenies.
\end{proof}

\subsection{The integral Hom ideal}
\begin{proposition}[Integral maps and an ideal class of order two]
\label{pp:prop:homideal}
Set $R=(I_C:I_C)=\OL+\mathfrak P$, where
$\mathfrak P=(1-z^{14})\OM$, and define
\[
 \mathfrak a=\alpha\OL+\mathfrak P,
 \qquad \gamma=z^6(1-z^2).
\]
All integral homomorphisms $B_B\to B_C$ are, uniquely,
\[
 F_\eta=\frac{m_0}{\alpha}\eta\sigma_{31},\qquad
 \eta\in\mathfrak a.
\]
The normalized module $\mathfrak a$ is an invertible fractional
$R$-ideal under target postcomposition, with
\[
 [\OM:\mathfrak a]=27,
 \qquad \mathfrak a\bar{\mathfrak a}=R,
 \qquad \mathfrak a^2=\gamma R,
 \qquad [\mathfrak a]\in\operatorname{Pic}(R)
       \text{ of exact order }2.
\]
For every nonzero $\eta$, the map $F_\eta$ is an isogeny of degree
$N_{M/\mathbb Q}(\eta)$. If
$E_{B,q}(v,w)=\operatorname{Tr}_{M/\mathbb Q}(\xi_B qv\bar w)$,
then its actual pullback on the source lattice is
\[
 F_\eta^*E_C=E_{B,\,\sigma_{19}(\eta\bar\eta)}.
\]
In particular the degree-$7$ map of Proposition~\ref{pp:prop:mindegree}
has pullback polarization type $(1,1,1,1,1,7)$.
\end{proposition}
\begin{proof}
Every rational Hodge map has the unique form $m\sigma_{31}$.
The map $(m_0/\alpha)\sigma_{31}$ identifies $S_B$ with $S_C$
by Proposition~\ref{pp:prop:unpolarizeddistinct}.
Thus an integral map must have $\eta S_C\subseteq S_C$, so
$\eta\in\OM$. Away from $3$, both original lattices equal their
saturations locally, because their saturation indices are $27$;
there is no further local restriction there.
At $3$, the map $F_0$ is a lattice isomorphism and $\alpha$ is a unit.
The remaining condition is exactly
\[
 \eta/\alpha\in R\otimes\Z_{(3)}.
\]
Since $R=\OL+\mathfrak P$ has residue field $\mathbb F_{27}$
inside $\OM/\mathfrak P=\mathbb F_{729}$, this says that the
residue of $\eta$ lies in the line $\alpha\mathbb F_{27}$.
Its full inverse image in $\OM$ is precisely $\mathfrak a$.
This proves the description of all maps, including sufficiency.

These module equalities also have direct finite certificates. Collecting
the $144$ entries of the map for each coefficient of $\eta$ gives
rational linear integrality constraints. The primary certificate
\path{certificates/principal_bc_hom_v127.py} takes the exact dual
of their row lattice; the result is $\mathfrak a$, not just a
sublattice of candidate maps. The independent implementation
\path{independent/principal_bc_hom_v127.py} derives the same module
from linear conditions modulo $3$, with containment and index checks.
The primary additionally checks $R=\OL+\mathfrak P$; both verify
the full products
$\mathfrak a\bar{\mathfrak a}=R$ and
$\mathfrak a^2=\gamma R$ on integral bases.
They check that $\gamma\in\OM^\times$ has absolute norm $1$.
The first product gives an explicit inverse fractional ideal.
If $\mathfrak a=\delta R$ were principal, the equality
$[\OM:\mathfrak a]=[\OM:R]=27$ would give
$N_{M/\mathbb Q}(\delta)=1$. Since $\delta\in\mathfrak a$,
the corresponding map would be an isomorphism, contradicting
Proposition~\ref{pp:prop:unpolarizeddistinct}. Hence the class is
nontrivial, and the second product proves that its order is exactly two.
The full Picard group is determined in Proposition~\ref{pp:prop:picardorder} below.

The degree formula is the saturation-index argument of
Proposition~\ref{pp:prop:mindegree}. For the polarization, the identity
$m_0\bar m_0=2r$ and $\alpha\bar\alpha=2$ gives
$m\bar m=r\eta\bar\eta$. Moving $\sigma_{31}$ through the trace
therefore gives
\[
 E_C(F_\eta v,F_\eta w)
 =\operatorname{Tr}_{M/\mathbb Q}
   \bigl(\xi_B\sigma_{19}(\eta\bar\eta)v\bar w\bigr),
\]
since $\sigma_{19}=\sigma_{31}^{-1}$. The inverse twist is essential
on source coordinates. For nonzero $\eta$ its multiplier is totally
positive; integrality follows from the integral map.
The independent checker verifies all $78$ quadratic coefficients
of this identity on a basis of $\mathfrak a$, including the mixed terms.
For the degree-$7$ witness, both paths check the alternating Smith
invariants $(1^{10},7^2)$, giving the displayed polarization type.
The formula describes exactly the pullbacks arising from this Hom
module; it does not assert that every totally positive element is realized.
\end{proof}

\subsection{Ideal classes and principal polarizations}
\begin{proposition}[The full Picard group of the principal-model order]
\label{pp:prop:picardorder}
For the order $R=\OL+\mathfrak P$ in
Proposition~\ref{pp:prop:homideal},
\[
 \operatorname{Pic}(R)\cong\mathbb Z/2\mathbb Z,
 \qquad
 \operatorname{Pic}(R)=\{[R],[\mathfrak a]\},
 \qquad \mathfrak a=\alpha\OL+\mathfrak P.
\]
Here $\operatorname{Pic}(R)$ denotes invertible fractional $R$-ideals
modulo principal fractional ideals, without an orientation or positivity
condition.
\end{proposition}
\begin{proof}
Write $O=\OM$. The conductor and residue computations above give
\[
 O/\mathfrak P=\mathbb F_{729},\qquad
 R/\mathfrak P=\mathbb F_{27}.
\]
Masley's Main Theorem gives class number one for
$M=\mathbb Q(\zeta_{42})=\mathbb Q(\zeta_{21})$:
$\varphi(21)=12\le24$, and $21$ is absent from its exceptional table
\cite[p.~179]{Masley1976}. This concerns the maximal order $O$;
the following argument determines
the class group of the proper order $R$.

Let $I$ be an invertible fractional $R$-ideal. Since $IO$ is principal,
rescale $I$ so that $IO=O$. Then
\[
 \mathfrak P\subseteq I\subseteq O:
 \qquad \mathfrak P=\mathfrak P(IO)\subseteq I
\]
because $\mathfrak P O\subseteq R$.
Away from the conductor maximal ideal $\mathfrak P$ of $R$,
the order is maximal and the normalized ideal equals $O$ locally.
At $\mathfrak P$, invertibility gives
$I_{\mathfrak P}=uR_{\mathfrak P}$, where normalization makes
$u$ a unit of $O_{\mathfrak P}$. Therefore $I/\mathfrak P$ is a
nonzero $\mathbb F_{27}$-line in $\mathbb F_{729}$.
Conversely, the inverse image $J_a$ of any such line
$a\mathbb F_{27}$ is locally principal: at $\mathfrak P$ it is
$uR_{\mathfrak P}$ for a lift $u$ of $a$, and elsewhere it is $O$.
Thus every such $J_a$ is invertible, and
$J_aJ_b=J_{ab}$ by the same local calculation.

Any isomorphism between fractional ideals is multiplication by an
element of $M^\times$. Between normalized ideals it must preserve
their $O$-saturations, so it is multiplication by a global $O$-unit.
Consequently this classifies every invertible ideal class and gives
\[
 \operatorname{Pic}(R)\cong
 \frac{\mathbb F_{729}^{\times}}
 {\mathbb F_{27}^{\times}\,\operatorname{im}(O^{\times})}.
\]
The intermediate quotient
$Q=\mathbb F_{729}^{\times}/\mathbb F_{27}^{\times}$ is cyclic
of order $28$. Conjugation induces $x\mapsto x^{27}$ on the residue
field, and the ratio map
\[
 Q\longrightarrow\mathbb F_{729}^{\times},
 \qquad [x]\longmapsto x/x^{27}
\]
is injective, since its kernel consists of the fixed-field units.

For every $\varepsilon\in O^{\times}$, Kronecker's theorem gives
$\varepsilon/\bar\varepsilon\in\mu_{42}$. The reduction of
$\mu_{42}$ has order $14$, so the image of the entire global unit
group in $Q$ is contained in the subgroup of order $14$.
It equals that subgroup. Indeed $z$ has residue order $14$,
and its ratio $z^2$ has order $7$. The global unit
$\gamma=z^6(1-z^2)$ has norm $1$ and
\[
 \gamma/\bar\gamma=-z^{14}\equiv-1\pmod{\mathfrak P},
\]
so its class has order $2$. These two classes generate a subgroup
of order $14$. The remaining quotient therefore has order $2$.

Finally, $\alpha\bar\alpha=2$ makes the residue of $\alpha$
nonzero. With $s=\sqrt{-7}$ chosen as above,
\[
 \alpha/\bar\alpha\equiv s\pmod{\mathfrak P},
 \qquad s^2\equiv-1\pmod{\mathfrak P}.
\]
Its ratio has order $4$, so its class lies outside the global-unit
image of order $14$. Its residue line has inverse image
$\mathfrak a=\alpha\OL+\mathfrak P$. This identifies the
nonidentity class and agrees with the ideal product identities of
Proposition~\ref{pp:prop:homideal}.

The finite inputs are checked by
\path{certificates/principal_order_class_v129.py}, using Frobenius
and polynomial gcds over $\mathbb F_3$ and exact rational resultants.
The separate implementation
\path{independent/principal_order_class_v129.py} checks all $39$
possible monic divisors of degrees at most three, enumerates the
$729$ residue elements, and computes the global norms by rational
multiplication matrices. Both identify the residue field using
\[
 \Phi_{42}(X)\bmod3
 =(X^6-X^5+X^4-X^3+X^2-X+1)^2.
\]
These finite checks supplement the class-number theorem and the
ideal and global-unit arguments above; they do not enumerate
the infinite group of global units.
\end{proof}

\begin{proposition}[The ordinary ideal-class monoid of the selected order]
\label{pp:prop:idealclassmonoid}
Let $O=\OM$ and retain $R$ and $\mathfrak a$ from
Proposition~\ref{pp:prop:picardorder}. Let $\operatorname{ICM}(R)$
be the monoid of nonzero fractional $R$-ideals that are full
finitely generated $\mathbb Z$-lattices in $M$, modulo $R$-linear
isomorphism, with ordinary ideal multiplication. Then
\[
 \operatorname{ICM}(R)=\{[R],[\mathfrak a],[O]\},\qquad
 [\mathfrak a]^2=[R],\quad [\mathfrak a][O]=[O]^2=[O].
\]
Its unit group is $\operatorname{Pic}(R)\cong\mathbb Z/2\mathbb Z$;
$[O]$ is the remaining, absorbing idempotent. The multiplier rings
$(I:I)=\{x\in M:xI\subseteq I\}$ are $R$ for the first two
classes and $O$ for the third. These are ordinary rank-one classes:
the zero ideal is excluded, and equivalence is $R$-linear, without
orientation or semilinear identifications. This statement for the
selected order makes no geometric realization or polarization claim.
\end{proposition}
\begin{proof}
Use the conductor $\mathfrak P$, residue fields
$k=R/\mathfrak P=\mathbb F_{27}$ and
$k_2=O/\mathfrak P=\mathbb F_{729}$, and
$\operatorname{Pic}(O)=1$ established in the preceding proposition.
For any ideal $I$ in the stated class, $IO$ is a nonzero fractional
ideal of the Dedekind domain $O$. It is principal, even when $I$
is not invertible over $R$. A scalar rescaling therefore gives
$IO=O$ and $I\subseteq O$. Write $1=\sum_j i_j o_j$ with
$i_j\in I$ and $o_j\in O$. For $p\in\mathfrak P$,
\[
 p=\sum_j i_j(po_j)\in I,
\]
since $po_j\in\mathfrak P\subseteq R$. Thus
$\mathfrak P\subseteq I\subseteq O$, and $I/\mathfrak P$
is a nonzero $k$-subspace of $k_2$; otherwise $I=\mathfrak P$
would contradict $IO=O$. Since $[k_2:k]=2$, this subspace is
either a line or all of $k_2$. The latter case gives $I=O$.

For $t\in k_2^\times$, let $J_t$ be the inverse image of $tk$.
A lift of $t$, together with $\mathfrak P$, generates $O$ as an
$O$-ideal, so $J_tO=O$. The product $J_tJ_u$ also has saturation
$O$ and hence contains $\mathfrak P$ by the same argument.
Its residue is exactly $tuk$, giving
\[
 J_tJ_u=J_{tu},\qquad J_1=R,\qquad J_tJ_{t^{-1}}=R.
\]
Every line lift is therefore invertible over $R$.

An $R$-linear isomorphism between these full rank-one ideals
extends over $\operatorname{Frac}(R)=M$ to multiplication by
$c\in M^\times$. Between normalized ideals it satisfies $cO=O$,
so $c\in O^\times$. The preceding proposition gives exactly two
line classes, represented by $R$ and $\mathfrak a$. The class of
$O$ is distinct, since an $O$-unit cannot carry $O$ to a proper
line lift. Normalization gives $[I][O]=[O]$ for every class,
and the remaining products follow from $\operatorname{Pic}(R)$.

Finally, $xI\subseteq I$ for a normalized ideal implies
$xO\subseteq O$, hence $x\in O$. For $I=J_t$, preservation of
the residue line $tk$ is equivalent to $x\bmod\mathfrak P\in k$,
so $(J_t:J_t)=R$; also $(O:O)=O$. These rings are unchanged by
scalar rescaling. As $R\ne O$, the third class is not invertible:
an inverse ideal $H$ would give $xR=xIH\subseteq IH=R$ for
every $x\in(I:I)$, forcing $(I:I)=R$.
\end{proof}

\begin{proposition}[Marked principal polarizations for the selected order]
\label{pp:prop:markedpolarizations}
Put $O=\OM$, and retain $R$ and $\mathfrak a$ from
Proposition~\ref{pp:prop:homideal}. Fix a marked CM type $\Phi$ of $M$,
obtained from $\Phi_0=\{1,5,11,13,17,19\}$ by conjugating a subset
$T$ of its six embeddings. Among full fractional $R$-ideals, up to
$M$-linear polarized isomorphism, the principal polarizations whose
Rosati involution on $M$ is complex conjugation are as follows:
\[
\begin{array}{c|ccc|c}
 & [R]&[\mathfrak a]&[O]&\text{total at }\Phi\\ \hline
|T|\text{ even}&2&2&0&4\\
|T|\text{ odd}&0&0&1&1
\end{array}
\]
The origin-preserving marked polarized automorphism groups are
$\mu_6$ in the first row and $\mu_{42}$ in the second. Thus there
are $160$ marked polarized classes over the $64$ marked CM types.
This count makes no identification between different types and no
assertion that these objects lie on a specified curve.
\end{proposition}
\begin{proof}
Write $I^\vee=\{x\in M:\Tr_{M/\Ql}(xI)\subseteq\Z\}$.
An $M$-compatible alternating form is
$E_\xi(x,y)=\Tr_{M/\Ql}(\xi x\bar y)$, with $\bar\xi=-\xi$.
Its dual lattice is
\[
 I^{\#_\xi}=\xi^{-1}(\bar I)^\vee
             =\xi^{-1}(R^\vee:\bar I).
\]
Indeed $x\bar I\subseteq R^\vee$ is equivalent to
$\Tr(x\bar I)\subseteq\Z$, since $\bar I R=\bar I$.
For invertible $I$, this colon equals $R^\vee\bar I^{-1}$.
Replacing $R^\vee$ by the maximal-order codifferent here would be
incorrect for a nonmaximal order.

Set $d=f'(t)$, retain $\beta$ from Proposition~\ref{pp:prop:explicit},
and put
\[
 \kappa_O=\frac1{d(z-\bar z)},\qquad
 \kappa_R=\frac1{d\beta(z-\bar z)}.
\]
The identity $\beta O=(1-z^{14})O$ gives
$O=\OL\oplus\OL z$ and $R=\OL\oplus\OL\beta z$.
Monogenic relative trace duality, followed by trace transitivity,
therefore gives $O^\vee=\kappa_O O$ and $R^\vee=\kappa_R R$.
Since $\mathfrak a\bar{\mathfrak a}=R$, one also has
$(\bar{\mathfrak a})^\vee=\kappa_R\mathfrak a$.
Consequently principal compatible parameters on $R,\mathfrak a$
are $\kappa_R u$, and those on $O$ are $\kappa_O u$, with
$u\in E_L=\OL^\times$. For the nonmaximal models this follows
because the real units of their multiplier ring $R$ are exactly $E_L$.
The conductor identity is certified by
\path{certificates/cm_polarization_conductor_v135.py}: it verifies
$\beta=(1-z^{14})a_0$ with
$a_0=1+z+z^5+z^6-z^8+z^{11}$ an integral unit of norm one.

With the convention $E(x,Jx)>0$, positivity means
$\operatorname{Im}\phi(\xi)>0$ for every $\phi\in\Phi$.
At $\Phi_0$, the signs of $d$ are $(+,-,+,-,+,-)$ and
$\operatorname{Im}\phi(z)>0$. Thus $\kappa_O$ has odd negative
sign parity. Since $N_{L/\Ql}(\beta)=-27$, $\kappa_R$ has even
parity. Conjugating an embedding reverses its imaginary sign.
Lemma~\externalref{supplement}{lem:narrowclass} gives exactly the even-parity subgroup
as the real-unit signature image, and
$[E_L^+:E_L^2]=2$. This proves the asserted existence and
nonexistence of positive parameters.

We next compute the relevant unit norms. For $e\in R^\times$,
Kronecker's theorem gives $e/\bar e\in\mu_{42}$.
Its residue modulo $\mathfrak P$ is one because $R/\mathfrak P$
is the conjugation-fixed field $\mathbb F_{27}$.
The reduction of $\mu_{42}$ has order $14$, so its kernel is
$\mu_3$. Choose $\zeta\in\mu_6$ with $\zeta^2=e/\bar e$.
Then $e/\zeta\in E_L$. Conversely $\mu_6E_L\subset R^\times$,
giving
\[
 R^\times=\mu_6E_L,\qquad N_{M/L}(R^\times)=E_L^2.
\]
By Theorem~\externalref{supplement}{thm:unitnorm}, $N_{M/L}(O^\times)=E_L^+$.
For a fixed positive parameter the polarization classes form a torsor
under $E_L^+/N_{M/L}(S^\times)$, where $S=(I:I)$.
Proposition~\ref{pp:prop:idealclassmonoid} exhausts the three scalar
classes and their multiplier rings, so these quotients prove the table.
Polarized automorphisms have norm one and are roots of unity.
The roots in $R$ are exactly $\mu_6$, while those in $O$ are
$\mu_{42}$. Finally $32\cdot4+32\cdot1=160$.
\end{proof}

\begin{proposition}[All principal polarizations on the specified B/C models]
\label{pp:prop:actualpolarizations}
Let $B_B,B_C$ be the two specified principal models, with their
original polarizations $\lambda_B,\lambda_C$. Put
$\varepsilon_+=4-t^2$. Their scalar ideal classes are respectively
$[\mathfrak a]$ and $[R]$. On each underlying variety $B_X$,
exactly two isomorphism classes of principal polarizations occur,
represented by $\lambda_X$ and $\lambda_X\circ\iota(\varepsilon_+)$.
After forgetting all CM and RM markings these still give exactly
four polarized isomorphism classes over the two underlying varieties.
Each has origin-preserving polarized automorphism group $\mu_6$.
\end{proposition}
\begin{proof}
Use the exact $I_B,I_C,\xi_B,\xi_C$ coordinates of
Proposition~\ref{pp:prop:unmarkeddistinct}, and $m_0,\alpha$ from
Proposition~\ref{pp:prop:homideal}. The scalars
\[
\begin{split}
 c_B={}&8+11z+4z^2+4z^3+8z^4-5z^5-22z^6-24z^7\\
       &{}-2z^8+20z^9+18z^{10}+3z^{11},\qquad
 c_C=\alpha\sigma_{31}(c_B)/m_0
\end{split}
\]
satisfy $c_BI_B=\mathfrak a$ and $c_CI_C=R$.
The certificate \path{certificates/actual_bc_polarizations_v135.py}
verifies integral unimodular changes of basis for both equalities,
and transports each form to
$\xi'_X=\xi_X/(c_X\bar c_X)=\kappa_Ru_X$ with $u_X\in E_L$.
It checks the exact alternating matrices and their pullbacks against
the original principal matrices. The complete coefficients and basis
changes are retained in
\path{certificates/actual_bc_polarizations_data.json}.
These are finite identities using one arithmetic implementation;
the following classification is a human argument.

There is also a structural reason for $c_C$.
For $D=(m_0/\alpha)\sigma_{31}(I_B)$, the complete Hom calculation
gives $(I_C:D)=\mathfrak a$. The proved B identification makes
$D=(m_0/\alpha)\sigma_{31}(c_B)^{-1}\bar{\mathfrak a}$ invertible;
here $\sigma_{31}(\mathfrak a)=\bar{\mathfrak a}$.
Multiplying by $D$ and using
$\mathfrak a\bar{\mathfrak a}=R$ gives the C identification.

For $\gamma=z^6(1-z^2)$ one has
\[
 \varepsilon_+=\gamma\bar\gamma=4-t^2,\qquad
 \gamma/\bar\gamma=-z^{14}\notin\mu_{21}.
\]
Thus $\varepsilon_+$ is a totally positive real unit.
If it were $v^2$ with $v\in E_L$, Kronecker's theorem would give
$\gamma/v\in\mu_{42}$, forcing
$\gamma/\bar\gamma\in\mu_{42}^2=\mu_{21}$, a contradiction.
It represents the nontrivial class of $E_L^+/E_L^2$.
Proposition~\ref{pp:prop:markedpolarizations} now supplies the two
compatible principal polarization classes on each model.

These are all principal polarizations, even without a marking.
The actual B and C types in Proposition~\ref{pp:prop:markeddistinct}
each have trivial stabilizer in $\operatorname{Gal}(M/\Ql)$.
The disjoint permutation supports of
$\operatorname{End}_{\Ql}(M)=\bigoplus_\sigma M\sigma$
then show $\operatorname{End}^0(B_X)=M$.
For any polarization $\lambda=\lambda_Xa$, its Rosati involution is
$x\mapsto a^{-1}\bar x a=\bar x$, since $M$ is commutative.
Thus every polarization is compatible, and every automorphism lies
in the multiplier group $R^\times$. Its norm subgroup is $E_L^2$,
so the two classes remain distinct without markings.
Proposition~\ref{pp:prop:unpolarizeddistinct} excludes all cross
identifications between B and C. The automorphism groups follow
from Proposition~\ref{pp:prop:markedpolarizations}.
No corresponding claim is made for the induced A type.
\end{proof}


% !TeX root = ../../packet-and-principal-models.tex
\section{The unit twist and the principal curve image}
\begin{proposition}[A global unit twist and its coarse involution]
\label{pp:prop:unitfamily}
On the actual family $(\mathcal B,\iota,\lambda_0)$ of
Theorem~\ref{thm:principalfamily}, the homomorphism
$\lambda_1=\lambda_0\circ\iota(\varepsilon_+)$ is another
principal polarization over the same base $C'$.
After forgetting auxiliary level, the two polarized families are
nonisomorphic over the identity of $C'$, and twisting twice returns
their polarized-family isomorphism class. On the full coarse
component of Proposition~\ref{pp:prop:explicit}, this operation is
the nontrivial involution
\[
 \tau_+(w)_i=\frac{w_i}{\sigma_i(\varepsilon_+)},\qquad
 F_{+,i}(z)=\frac{\sigma_i(a^2/\varepsilon_+)F_{t,i}(z)}
                         {3F_{t,i}(z)+1}.
\]
The specified B/C fibers receive their second polarization classes
from Proposition~\ref{pp:prop:actualpolarizations}.
\end{proposition}
\begin{proof}
Global RM self-adjointness makes $\lambda_1$ symmetric.
Its Riemann form multiplies each real eigenspace block of the
original positive form by $\sigma_i(\varepsilon_+)>0$.
Since $\varepsilon_+^{-1}\in\OL$, it is a principal polarization.
This is a global algebraic polarization homomorphism; no choice
of a global ample line bundle is required. The exact marking
comparison in Proposition~\ref{pp:prop:markeddistinct} identifies
$z+z^{-1}$ with the fixed global $t$, so the B/C twists are literal
instances of the preceding proposition.

In the standard module $P_{\rm std}$ put
$D=\operatorname{diag}(\varepsilon_+^{-1},1)$.
It preserves the lattice and satisfies
$E_{\varepsilon_+}(Dx,Dy)=E_1(x,y)$.
Evaluation at $w$ after $D$ is evaluation at $w/\varepsilon_+$.
This gives the displayed period formula with its stated direction.
Conjugating $\left(\begin{smallmatrix}r&\delta s\\v/\delta&h\end{smallmatrix}\right)$
by $D$ gives
$\left(\begin{smallmatrix}r&\delta s/\varepsilon_+\\
\varepsilon_+v/\delta&h\end{smallmatrix}\right)$,
so $D$ normalizes $\Gamma_{\rm std}$. Moreover
$D^2=\varepsilon_+^{-1}\operatorname{diag}(\varepsilon_+^{-1},\varepsilon_+)$;
the second factor belongs to $\Gamma_{\rm std}$ and the scalar
acts trivially on the domain. This proves the coarse involution.
Equivalently, twisting twice is pullback by $\iota(\varepsilon_+)$.
The distinct B/C polarization classes witness nontriviality of
both the coarse action and the two families over the identity base.
\end{proof}

\begin{proposition}[The unit involution preserves the principal curve]
\label{pp:prop:unitcurve}
The involution $\tau_+$ preserves the entire image $V_{\rm std}$ of
\eqref{pp:eq:curve}, and its restriction to that image is nontrivial.
More precisely, let $a_\triangle,b_\triangle,c_\triangle$ be the
original triangle generators, so that
$A=a_\triangle^2$, $B=b_\triangle$, $C=c_\triangle^2$, and put
$g=BACc_\triangle$. For the actual coarse moduli map
$f_0(z)=[WF_t(z)]$ one has
\[
 \tau_+f_0(z)=f_0(g^{-1}z).
\]
\end{proposition}
\begin{proof}
The full-triangle equivariance of the actual constructed map is essential.
McMullen's Section~4 identifies the quadratic trace-field extension and
its nontrivial automorphism, which sends the triangle generators to
$(-a_\triangle,b_\triangle,-c_\triangle)$.
Thus the two extensions of each real embedding of $L$ give identical
projective matrices and ordered target vertices. The vertex-respecting
triangle maps in \cite[Thm.~6.3, Eq.~(6.1)]{McMullen2026} therefore agree
for the two lifts. Keeping one from each pair yields the six components
of the actual $F$, with equivariance for the full triangle group.
Conjugating target coordinates by $h_t$ preserves this equivariance.
This argument would not follow from arbitrary $\Delta_0$-equivariance.

Set
\[
 s_0=-2-3t+4t^2+t^3-t^4,\quad
 h=s_0(I+C_t),\quad k=B_tA_tC_th.
\]
Exact field arithmetic gives $s_0^2=2-t$ and integral inverses for
$s_0$ and $\varepsilon_+=4-t^2$. Consequently
$\det h=\det k=\varepsilon_+$ and $h,k\in\GL_2^+(\OL)$.
If $\theta=2\cos(\pi/42)$, Cayley--Hamilton gives
$I+c_\triangle^2=\theta c_\triangle$.
Therefore $h$ is projectively the normalized $c_\triangle$, and
$F_t(gz)=kF_t(z)$, componentwise. The first-component base matrix
is $h_{t,1}^{-1}k_1h_{t,1}$ projectively; it is not $k_1$ in the
original base marking.

For $D=\operatorname{diag}(\varepsilon_+^{-1},1)$ the explicit witness
satisfies
\[
 T^{-1}k^{-\mathsf T}T\in\GL_2(\OL),\qquad
 U=\operatorname{diag}(\delta,1)^{-1}
       DWkW^{-1}\operatorname{diag}(\delta,1)\in\SL_2(\OL).
\]
The exact coefficient arrays, integral inverses and determinant identities
are checked by \path{certificates/curve_twist_image_v138.py} with
\path{certificates/curve_twist_image_data.json}.
Only this fixed word is needed; no search-completeness assertion enters.
In particular $\gamma=DWkW^{-1}$ belongs to the specified
$\Gamma_{\rm std}$. Hence
\[
 [DWF_t(z)]=[\gamma Wk^{-1}F_t(z)]
          =[WF_t(g^{-1}z)].
\]
Since $g^{-1}$ maps $\HH$ bijectively to itself, this proves equality
of the entire images. They are the actual algebraic images by
Theorem~\ref{thm:principalfamily} and
Proposition~\ref{pp:prop:unitfamily}.
The moved B/C points in Proposition~\ref{pp:prop:actualpolarizations}
show that the restricted involution is nontrivial.
\end{proof}

\begin{proposition}[Two principal classes over each specified RM variety]
\label{pp:prop:sixprincipal}
For each $X=A,B,C$, the underlying marked variety $(B_X,\iota_L)$ has
exactly two RM-compatible principal-polarization isomorphism classes,
represented by $\lambda_X$ and
$\lambda_X\circ\iota_L(\varepsilon_+)$.
Their origin-preserving RM-polarized automorphism groups are $\mu_6$.
The resulting six marked objects are pairwise distinct and all lie
on $V_{\rm std}$.
\end{proposition}
\begin{proof}
Identify rational homology with the rank-one module over the chosen
$M$. This remains valid for the induced A type. Since $M/L$ is
quadratic, every $L$-linear rational endomorphism has the unique form
$m_a+m_b\,\overline{(\cdot)}$, for $a,b\in M$.
After complexification the first summand preserves each embedding line,
while the second exchanges every line with its conjugate.
Each CM type contains exactly one line from every conjugate pair;
if $b\ne0$, all its embedding values are nonzero. Thus the second
summand cannot preserve the Hodge decomposition or cancel with the
first. It follows that
\[
 \operatorname{End}^0_L(B_X)=M,\qquad
 \operatorname{End}_L(B_X)=R,
\]
where the integral equality uses Proposition~\ref{prop:survivingorder}
for all three selected field actions. The first equality is a
centralizer statement, not an assertion that the full endomorphism
algebra at A is $M$.

If $\lambda$ is another RM-compatible principal polarization, then
$a=\lambda_X^{-1}\lambda$ is an integral automorphism commuting
with $L$. Symmetry makes $a$ fixed by the original Rosati, hence
$a\in E_L$, and positivity gives $a\in E_L^+$.
Conversely each such unit supplies a compatible principal polarization.
The $L$-linear automorphisms are precisely $R^\times$, whose norms
are $E_L^2$. Therefore the classes are
$E_L^+/E_L^2=\{1,\varepsilon_+\}$ by
Propositions~\ref{pp:prop:markedpolarizations}
and~\ref{pp:prop:actualpolarizations}.
The same norm computation gives the stated automorphism groups.
Proposition~\ref{pp:prop:unpolarizeddistinct} excludes identifications
between different $X$, while Proposition~\ref{pp:prop:unitcurve}
realizes each second class on the same image at $g^{-1}z_X$.
\end{proof}

\begin{proposition}[Unmarked separation at the induced A fiber]
\label{pp:prop:unmarkeda}
The principally polarized varieties
$(B_A,\lambda_A)$ and $(B_A,\lambda_A\circ\iota_L(\varepsilon_+))$
are nonisomorphic after forgetting all CM and RM markings.
\end{proposition}
\begin{proof}
Use the actual selected generator $Z_A=(A_t+tI)/(t^2-1)$.
Its lower-left and lower-right entries are
\[
 c=\frac{t^3-1}{t(t^2-1)},\qquad
 d_0=\frac{-t^5+5t^3+t^2-5t}{t^2-1}.
\]
In $M=\Ql(z)$, set $w=(z-d_0)/c$.
The row $(w,1)$ is an eigenvector of $Z_A^{\mathsf T}$, so
$(x,y)\mapsto xw+y$ identifies the actual rational homology with $M$.
On the common principal basis of Proposition~\ref{pp:prop:explicit},
its lattice and polarization parameter are
\[
 I_A=\sum_{j=0}^5\mathbf Z\,awt^j
       +\sum_{j=0}^5\mathbf Z\,b(1+3w)t^j,\qquad
 \xi_A=-\frac{c}{z-\bar z}.
\]
The trace form $\operatorname{Tr}_{M/\Ql}(\xi_Ax\bar y)$
reproduces the actual principal alternating matrix $E$ exactly.

In $\operatorname{End}_{\Ql}(M)=\bigoplus_\sigma M\sigma$,
the summands have disjoint permutation supports on the complex
embedding lines. A nonzero summand is a Hodge endomorphism precisely
when its automorphism stabilizes the actual type
$\{1,5,17,19,29,31\}$. Its stabilizer is $\{1,29\}$, and therefore
\[
 \operatorname{End}^0(B_A)=M+M\sigma_{29}.
\]
In particular the full algebra, of rational dimension $24$, is used;
restricting to multiplication or to one semilinear summand would
not prove the assertion.

Let $U_1,\ldots,U_{24}$ be the rational homology matrices for
$z^j$ and $z^j\sigma_{29}$, $0\le j<12$, in the displayed basis
of $I_A$. Form the $144\times24$ matrix $C$ by flattening them.
The full integral endomorphism coefficient lattice is
\[
 \mathcal L=\{v\in\Ql^{24}:Cv\in\mathbf Z^{144}\}.
\]
Equivalently, it is the rational dual of the integer lattice
generated by the constraint rows. This gives a complete lattice
calculation, including cancellations between the two summands.

Let $Q_+$ be the homology action of $\varepsilon_+$.
An unmarked polarized isomorphism $U$ would satisfy
$U^{\mathsf T}EU=EQ_+$, hence
\[
 \operatorname{Tr}_{\Ql}(Q_+^{-1}U^\dagger U)=12,\qquad
 U^\dagger=E^{-1}U^{\mathsf T}E.
\]
This is a positive definite quadratic form on the real Hodge
endomorphism algebra: the positive Hodge metric makes $\dagger$
the adjoint, and $Q_+$ is positive self-adjoint.
Its exact Gram matrix on $\mathcal L$ has an integral unimodular
reduction $H$ with positive rational LDL data $(\mu,d)$.
Every vector in its closed norm-$12$ ball satisfies
$|x_i|^2\le12(H^{-1})_{ii}$. The exact decomposition
\[
 x^{\mathsf T}Hx=\sum_i d_i
       \left(x_i+\sum_{j>i}\mu_{ji}x_j\right)^2
\]
therefore supplies finite coordinate bounds and exhaustive
suffix pruning using only rational comparisons.

The certificate \path{certificates/a_unmarked_polarization_v140.py}
reconstructs the actual model, full dual lattice and Gram matrix,
checks the integral unimodular reduction, positive LDL data and
bounds in \path{certificates/a_unmarked_polarization_data.json},
and exhausts this ball in $347$ nodes: only zero occurs.
The separate implementation
\path{independent/a_unmarked_polarization_v140.py} works directly
on $L^2$ homology, constructs $\sigma_{29}$ by a cyclic vector,
proves fullness by integer row saturation, and recomputes the
Gram and LDL in reverse coordinate order. Its $925$ nodes again
leave only zero. Thus the required norm-$12$ vector does not exist.
\end{proof}

\begin{proposition}[The four B/C polarization pairs]
\label{pp:prop:polarizationsquare}
Put $B_e=(B_B,\lambda_B\varepsilon_+^e)$ and
$C_f=(B_C,\lambda_C\varepsilon_+^f)$, for $e,f\in\{0,1\}$.
For every pair $B_e\to C_f$ the least positive integer scalar
polarization multiplier is $2$, attained by exactly six isogenies,
all of degree $64$.
With $\gamma_0=z^6(1-z^2)$ and
$v_{ef}=\sigma_{31}(\gamma_0)^e\gamma_0^{-f}$,
their normalized Hom parameters are
\[
 \eta=\begin{cases}
 \alpha v_{ef}z^{7j},&e=f,\\
 \bar\alpha v_{ef}z^{7j},&e\ne f,
 \end{cases}\qquad 0\le j<6.
\]
The six maps share one kernel for each pair. The pairs $(0,0),(1,1)$
have the same kernel $K_+$, while $(0,1),(1,0)$ have a second kernel
$K_-$. These are complementary order-$64$ subgroups of $B_B[2]$.
\end{proposition}
\begin{proof}
Proposition~\ref{pp:prop:homideal} gives every rational homomorphism
as $F_\eta=(m_0/\alpha)\eta\sigma_{31}$, integral exactly when
$\eta\in\mathfrak a$. Its source pullback multiplier is
$\sigma_{19}(\eta\bar\eta)$. For the two twisted forms the required
multiplier-$n$ equation is consequently
\[
 \eta\bar\eta=n\frac{\sigma_{31}(\varepsilon_+)^e}
                         {\varepsilon_+^f}.
\]
Since $\gamma_0\bar\gamma_0=\varepsilon_+$,
$N_{M/L}(v_{ef})$ is the unit factor on the right.

The polynomial $\Phi_{42}$ modulo $2$ is the product of the two
irreducible sextics
\[
 (X^6+X^4+X^2+X+1)(X^6+X^5+X^4+X^2+1).
\]
Together with $\alpha\bar\alpha=2$ and
$\alpha+\bar\alpha=1$, this identifies the two conjugate primes
above $2$ as $(\alpha),(\bar\alpha)$.
Every integral solution for $n=2$ therefore has one of these two
ideals. Division by $\alpha v_{ef}$ or $\bar\alpha v_{ef}$ gives
an integral unit of relative norm one, hence a root of unity.
Thus the complete candidates are the two families
$\alpha v_{ef}z^j$, $\bar\alpha v_{ef}z^j$, $0\le j<42$.

For membership in $\mathfrak a=\alpha\OL+\mathfrak P$,
consider the quotient $\mathbf F_{729}^\times/\mathbf F_{27}^\times$
of order $28$. The class of $z$ has order $7$.
The class of $\gamma_0$ is the unique element of order two:
$\gamma_0/\bar\gamma_0=-z^{14}$ reduces to $-1$.
The class of $\bar\alpha/\alpha$ is that same element, since
$\alpha^2-\alpha+2=0$ modulo $3$ gives
$\alpha^4=-1$ and $(\alpha/\bar\alpha)^2=-1$ there.
The automorphism $\sigma_{31}$ fixes this unique order-two class.
Consequently the displayed branch cancels the order-two class
exactly as stated, and the remaining exponent must be divisible by $7$.
There are exactly six maps. Their full norm, ideal membership,
integral homology matrices and polarization identities are verified in
\path{certificates/bc_polarization_square_v141.py} with its retained data.

Multiplier one would imply an isomorphism of the underlying B/C
varieties, excluded by Proposition~\ref{pp:prop:unpolarizeddistinct}.
Thus the displayed multiplier two is minimal. Multiplication by
$z^{7j}\in R^\times$ preserves each target form, so the six maps
for a fixed pair have the same kernel.

Let $T_{ef}$ denote the integral homology matrix for $j=0$.
The certificate verifies $\det T_{ef}=64$ and the four actual
polarization equations. The polarized adjoint relation implies
that every kernel lies in $B_B[2]$. It is therefore identified with
$\ker(T_{ef}\bmod2)$ on the common source basis.
All four ranks are $6$; the stacked ranks for $(00,11)$ and
$(01,10)$ are $6$, and the stacked rank for $(00,01)$ is $12$.
These exact finite calculations prove the asserted equality and
complementarity of kernels.
\end{proof}

\begin{remark}[Image and level boundaries]\label{pp:rem:unitimage}
The two parametrized moduli maps differ at the specified B/C parameters,
but their images agree by Proposition~\ref{pp:prop:unitcurve}.
The diagonal matrix in the square calculation need not lie in a chosen
neat subgroup, so no involution preserving an arbitrary fine-level
marking is asserted. These constructions neither add a point to the
original trace-polarized packet nor identify a principal image with
Mostaed's literal normalization.
\end{remark}

\begin{corollary}[Six Weil points on the specified principal image]
\label{pp:cor:weilimage}
Let $\mathcal W_{K,\rm std}$ denote the intrinsic $K$-Weil locus in the
specified principal RM component. Then
$V_{\rm std}\cap\mathcal W_{K,\rm std}$ contains the six distinct
marked points of Proposition~\ref{pp:prop:sixprincipal}.
All have $K$-signature $(3,3)$ and Hermitian discriminant $-1$.
The actual CM types and Hodge algebras are those of their respective
A/B/C underlying varieties. All six remain distinct as principally polarized varieties after
forgetting the RM and CM markings.
\end{corollary}
\begin{proof}
Keep the original lattice $P$, chosen field actions and integral
$\mathcal O_K$ generators of Proposition~\ref{prop:commonprincipal}.
Twisting replaces $E_1$ by
$E_{\varepsilon_+}(v,w)=E_1(\varepsilon_+v,w)$, where
$\varepsilon_+I_2$ is an integral automorphism of $P$.
The chosen $M$ action and therefore the actual CM type are unchanged.
When the new point is represented at $g^{-1}z_X$, these are the
transported actions through the explicit orbit identification above;
they are not inferred from a companion CM type.

The integral $K$ action commutes with $\varepsilon_+$ and its Rosati
remains conjugation. The complex $K$ action is unchanged, so the
signature is still $(3,3)$. On the six-dimensional $K$ homology space,
$M=LK$ and $L\cap K=\Ql$ give
\[
 \det_K(m_{\varepsilon_+})
 =N_{M/K}(\varepsilon_+)=N_{L/\Ql}(\varepsilon_+)=1.
\]
The last norm is also checked in the exact transport certificate.
Thus twisting the Hermitian form in one variable leaves its determinant
and discriminant unchanged. Positivity and principality were established
in Proposition~\ref{pp:prop:unitfamily}.
Changing polarization does not change the Hodge structure, so the
previous divisor/Weil generator and algebraicity conclusions apply.
Finally, the four B/C polarized objects remain distinct without markings
by Proposition~\ref{pp:prop:actualpolarizations}, and
Proposition~\ref{pp:prop:unmarkeda} separates the two A objects.
The underlying A variety is different from both B/C varieties by
Proposition~\ref{pp:prop:unpolarizeddistinct}. Hence all six remain
distinct without markings; no classification of all A polarizations
is needed.
\end{proof}

\begin{remark}[The integral comparison still required]
\label{pp:rem:obstruction}
For any full $O$ lattice $P'\subset L^2$ self-dual for $E_1$, its
determinant fractional ideal is $\bigwedge_O^2P'=\mathfrak D_L^{-1}$.
Indeed at every finite prime a local basis matrix $T_{\mathfrak p}$
turns the pairing into the standard alternating form multiplied by
$\det T_{\mathfrak p}$; perfection of the trace pairing is equivalent
to $(\det T_{\mathfrak p})O_{\mathfrak p}
=\mathfrak D_{L,\mathfrak p}^{-1}$. These equalities determine the
fractional ideal globally. Proposition~\ref{prop:noprincipal} shows
that the codifferent has no totally positive generator. Thus no
$T'\in\GL_2^+(L)$ satisfies $T'O^2=P'$. The positive matrix $W$ above
compares $P$ to the different module $P_{\rm std}$ and does not contradict
this obstruction.

The original trace-polarized fibers do not become principal by passing
to a cover alone. The chosen quotient kernel has orbit size $28$, so its
marked descent requires its stabilizer; this does not rule out descent
of an unmarked target or a suitable unspecified finite-index subgroup.
In particular Corollary~\ref{pp:cor:weilimage} concerns the explicit
principal image just constructed. Identifying it with Mostaed's named
$V$ requires an integral polarization/level comparison and equality of
curve images or direct point membership. Neither that identification
nor disjointness of the two images follows here. The original
standard lattice, a principal isogenous model, and literal membership
in a prescribed principally polarized locus remain distinct assertions.
\end{remark}


% !TeX root = ../../packet-and-principal-models.tex
\section{Scope and questions for further research}\label{sec:conclusion}
Under the execution premise, the fixed marked packet meets the normalized
curve at precisely three quotient points. The common lattice construction
and the unit twist give the stated principal models and six Weil points.
These conclusions keep the integral lattice, polarization and field marking
visible; they do not identify Mostaed's literal principal locus.

There are two natural continuations. First, determine whether the same
finite-domain method handles other norm generators and levels. Such an
extension needs a new arithmetic domain and new geometric certificates.
Second, identify the constructed principal family with a specified moduli
normalization, or prove that the integral data obstruct that identification.
The latter question must fix the lattice, level and polarization module.

The finite quotient family of Paper II starts from a primitive Fermat
factor, not from further points of this packet. Its Jacobian and descent
questions therefore do not constitute unfinished cases of the three-point
intersection theorem. The supporting Hodge note likewise concerns abstract
fixed-field types and does not add points to this curve.

\begingroup\small
% !TeX root = ../packet-and-principal-models.tex
\begin{thebibliography}{99}
\bibitem{DeligneHodge2006}
P. Deligne,
\emph{The Hodge Conjecture}, in J. Carlson, A. Jaffe and A. Wiles (eds.),
\emph{The Millennium Prize Problems}, Clay Mathematics Institute/AMS,
2006, 45--53.

\bibitem{Goodson2024}
H. Goodson,
\emph{An exploration of degeneracy in abelian varieties of Fermat type},
arXiv:2211.03909v2, 1 August 2024.
The cited numbering refers to this version.
\url{https://arxiv.org/abs/2211.03909v2}

\bibitem{LiMantovanPries2021}
W.~Li, E.~Mantovan and R.~Pries,
\emph{Data for Shimura varieties intersecting the Torelli locus},
arXiv:2105.02286v1 (2021); Pacific J. Math. \textbf{343} (2026), 179--210.
Proposition~4.5, Proposition~6.11 and Lemma~6.12 refer to the
inspected arXiv version.
\url{https://arxiv.org/abs/2105.02286v1}

\bibitem{Markman2025}
E. Markman,
\emph{Cycles on abelian $2n$-folds of Weil type from secant sheaves on abelian $n$-folds},
arXiv:2502.03415v2 (2025), Section~1.1, p.~3, and Theorem~1.5.1, p.~9.
The theorem proves algebraicity of Hodge--Weil classes for polarized
abelian sixfolds of Weil type with discriminant $-1$, without a
principal-polarization hypothesis.
\url{https://arxiv.org/pdf/2502.03415v2}

\bibitem{Masley1976}
J. M. Masley,
Solution of small class number problems for cyclotomic fields,
\emph{Compositio Mathematica} \textbf{33} (1976), 179--186,
\url{https://www.numdam.org/item/CM_1976__33_2_179_0.pdf}.

\bibitem{McMullen2026}
C. T. McMullen,
\emph{Triangle groups and Hilbert modular varieties},
J. Math. Soc. Japan 78 (2026), no.~3, 729--749; DOI 10.2969/jmsj/94179417.
The theorem and section locators used here refer to the author's version
dated 27 June 2023, typeset 27 March 2026, Sections~4 and~6
(printed pp.~14--15 and 19--21).
\url{https://people.math.harvard.edu/~ctm/papers/home/text/papers/quat/quat.pdf}

\bibitem{MilneModuli2011}
J. S. Milne,
\emph{Shimura Varieties and Moduli}, manuscript dated 30 April 2011,
Theorems 3.6 and 7.13, pp.~15 and 38--39. The numbering here refers to
this author version.
\url{https://www.jmilne.org/math/xnotes/svh.pdf}

\bibitem{MilneCM2020}
J. S. Milne,
\emph{Complex Multiplication}, version 0.10, 14 July 2020,
Section 3.11 and Propositions 3.12--3.13 and 3.17, pp.~29--31.
\url{https://www.jmilne.org/math/CourseNotes/CM.pdf}

\bibitem{MilneHodge2020}
J. S. Milne,
\emph{Hodge classes on abelian varieties},
arXiv:2010.08857, 2020.
\url{https://arxiv.org/abs/2010.08857}

\bibitem{Mostaed2026}
A. Mostaed,
\emph{McMullen's Curve, the Weil Locus, and the Hodge Conjecture for Abelian Sixfolds},
arXiv:2603.20268v1, 14 March 2026.
\url{https://arxiv.org/abs/2603.20268}

\bibitem{OpenAICM2026}
OpenAI math repository,
\emph{The rational Hodge conjecture for CM abelian varieties},
public manuscript dated 6 October 2026, introduction (main CM theorem)
and tensor-reduction section (auxiliary-factor contractions).
The comparison uses repository commit
\texttt{4eaee87548250b2682affda7d05d947c05554bae};
the general theorem and its formalization are not proof inputs here.
\url{https://github.com/openai/math/tree/4eaee87548250b2682affda7d05d947c05554bae/preprints/The-rational-Hodge-conjecture-for-CM-abelian-varieties-October-6-2026}

\bibitem{Saller2026}
K. Saller,
\emph{Split Weil structures on abelian sixfolds and the Hodge conjecture},
version 2, Zenodo record 23290280 (2026), Remarks~4.7--4.8.
\url{https://doi.org/10.5281/zenodo.23290280}

\bibitem{Schoen1998}
C. Schoen,
\emph{Addendum to ``Hodge classes on self-products of a variety with an automorphism''},
Compositio Math. 114 (1998), no.~3, 329--336.

\end{thebibliography}

\endgroup
\end{document}
